% ---------------------------------------------------------------------------
% deficiency.tex --- assembling a CRNT verdict on the Edelstein motif.
%
% Source fragment: loaded by SI/si.tex, exported standalone by
% Figure/export_tikz.py.
%
% The point of the figure is panel (c): both classical theorems come back
% inconclusive, and the package names the hypothesis that failed.  So (c)
% carries the vermillion accent and (a),(b) stay neutral/green -- the reader's
% eye should land on the verdict, not on the arithmetic that produced it.
%
% Every fact is shown once.  What a picture can carry -- the linkage classes,
% the signs of S, the conservation vector and the row of zeros it produces --
% is drawn, not restated in prose; the panels keep one short line each.
%
% Width: laid out on the 16.40 cm text block of this manuscript (see
% Figure/arch.tex) and included unscaled.
% ---------------------------------------------------------------------------
\IfFileExists{synstyle.tex}{% ---------------------------------------------------------------------------
% synstyle.tex --- shared look for every TikZ figure of the SynKit CRN paper.
%
% Input it from a figure file:
%     \IfFileExists{synstyle.tex}{% ---------------------------------------------------------------------------
% synstyle.tex --- shared look for every TikZ figure of the SynKit CRN paper.
%
% Input it from a figure file:
%     \IfFileExists{synstyle.tex}{% ---------------------------------------------------------------------------
% synstyle.tex --- shared look for every TikZ figure of the SynKit CRN paper.
%
% Input it from a figure file:
%     \IfFileExists{synstyle.tex}{\input{synstyle}}{\input{Figure/synstyle}}
% so that the figure compiles both from Figure/ and from the paper root.
%
% The visual vocabulary follows ../../../../Style/style.tex: the
% colour-blind-safe Okabe--Ito accents, quiet neutral rules and wash, circular
% panel badges, and a compact technical type scale are shared by the companion
% papers.  Semantic aliases below keep the CRN meaning explicit.
%
% DESIGN CONTRACT -----------------------------------------------------------
% 1. Geometry.  Every figure is laid out on a canvas exactly \scW wide (the
%    sn-jnl text width, 372 pt) and is included unscaled, so one point of type
%    is the same physical size in all five figures.  Open each picture with
%    \scCanvas{<height in cm>} to pin the bounding box.
% 2. Colour is semantic, never decorative.  Four roles only:
%       scAccent (green)       structure, the normal/expected path
%       scMark   (orange)      catalyst, certificate, or highlight
%       scAlert  (vermillion)  failure, violation, or uncovered set
%       scRef    (blue)        secondary or reference series
%    A panel gets at most one accent.  Titles are always ink; panel identity
%    is also encoded by a letter inside a circular badge, not by colour alone.
% 3. Type scale.  Three sizes, no others, and nothing below 6 pt:
%       \scTitle (7 pt bold)  panel titles
%       \scBody  (6.5 pt)     labels, values, in-figure prose
%       \scSmall (6 pt)       units, footnotes, dense graph labels
% ---------------------------------------------------------------------------
% Load guard.  The paper inputs every figure source inline, so this file is
% reached once per figure; without the guard the second \newlength below
% would abort the run.
\ifdefined\SynKitStyleLoaded\endinput\fi
\newcommand{\SynKitStyleLoaded}{}

% Load in the document preamble. XeLaTeX/LuaLaTeX and an installed Arial
% family are required; fontspec reports an error instead of substituting.
\RequirePackage[no-math]{fontspec}
\newfontfamily\scArial{Arial}[NFSSFamily=scArial,Ligatures=TeX]
% Keep manuscript mathematics unchanged; select Arial letters and numerals
% only inside figures. TeX supplies specialist mathematical symbols.
\RequirePackage[subdued,italic]{mathastext}
\MTDeclareVersion{synarial}{TU}{scArial}{m}{n}
\SetMathAlphabet{\Mathsf}{synarial}{TU}{scArial}{m}{n}
\SetMathAlphabet{\Mathtt}{synarial}{TU}{scArial}{m}{n}

\usetikzlibrary{positioning,calc,fit,backgrounds,arrows.meta,shapes.geometric,
                shapes.misc,decorations.pathreplacing,patterns,math}

% --- canvas ---------------------------------------------------------------
% 372 pt text width minus the 1 pt standalone border on each side.
\newlength{\scW}\setlength{\scW}{13.0326cm}
\providecommand{\scWnum}{13.0326}
% \scCanvas{<height>} pins the picture to the full text width, y running down
% from 0, so every figure shares one coordinate convention.  It *sets* the
% bounding box rather than contributing to it, so the emitted page is exactly
% the text width and every figure is included at 1:1 -- otherwise a stray
% invisible extent (pgfplots reserves a little around an axis) rescales one
% figure relative to the others and the same \scSmall prints at two sizes.
% Anything drawn outside the canvas is therefore cropped: keep content inside,
% and check with
%     \providecommand{\scCanvas}[1]{\path[draw=red] (0,0) rectangle (\scWnum,-#1);}
\providecommand{\scCanvas}[1]{%
  \path[use as bounding box] (0,0) rectangle (\scWnum,-#1);}

% --- palette --------------------------------------------------------------
% Canonical companion-paper values (../../../../Style/style.tex).
\definecolor{npgOrange}{HTML}{E69F00}
\definecolor{npgSky}{HTML}{56B4E9}
\definecolor{npgGreen}{HTML}{009E73}
\definecolor{npgYellow}{HTML}{F0E442}
\definecolor{npgBlue}{HTML}{0072B2}
\definecolor{npgVermillion}{HTML}{E64B35}
\definecolor{npgPurple}{HTML}{CC79A7}
\definecolor{figInk}{HTML}{1A1A1A}
\definecolor{figMuted}{HTML}{666666}
\definecolor{figHair}{HTML}{D9D9D9}
\definecolor{figWash}{HTML}{F7F7F7}

\colorlet{scInk}{figInk}              % primary text
\colorlet{scInkM}{figMuted}           % secondary text
\colorlet{scPaper}{figWash}           % quiet cards and note strips
\colorlet{scSpF}{npgSky!12}           % species fill
\colorlet{scSpE}{npgBlue}             % species edge
\colorlet{scRxF}{npgOrange!12}        % reaction fill
\colorlet{scRxE}{npgOrange!88!black}  % reaction edge
\colorlet{scReac}{npgBlue!78!black}   % reactant edge
\colorlet{scProd}{npgGreen!82!black}  % product edge
\colorlet{scHiF}{npgOrange!13}        % highlight fill
\colorlet{scHiE}{npgOrange!82!black}  % highlight edge
% Stage / accent tints.  stDE remains available for the fifth plot series.
\colorlet{stAF}{npgGreen!9}\colorlet{stAE}{npgGreen}
\colorlet{stBF}{npgOrange!11}\colorlet{stBE}{npgOrange}
\colorlet{stCF}{npgBlue!9}\colorlet{stCE}{npgBlue}
\colorlet{stDF}{npgPurple!10}\colorlet{stDE}{npgPurple!85!black}
\colorlet{stEF}{npgVermillion!9}\colorlet{stEE}{npgVermillion}

% Semantic aliases.  Figures should speak in these, not in stAE/stBE/...
\colorlet{scAccent}{stAE}\colorlet{scAccentF}{stAF}   % structure / normal
\colorlet{scMark}{stBE}\colorlet{scMarkF}{stBF}       % highlight / certificate
\colorlet{scRef}{stCE}\colorlet{scRefF}{stCF}         % secondary series
\colorlet{scAlert}{stEE}\colorlet{scAlertF}{stEF}     % failure / violation
\colorlet{scRule}{figHair}                            % card and axis hairlines
\colorlet{scRuleL}{figHair!55}                        % grid, dividers

% --- text styles ----------------------------------------------------------
% Three sizes, fixed in points so they do not drift with the document class.
% Arial is shared by labels, headings, panel badges, and chemistry labels.
\providecommand{\scTitle}{\scArial\fontsize{7}{8.4}\selectfont\bfseries\color{scInk}}
\providecommand{\scBody}{\scArial\fontsize{6.5}{7.8}\selectfont\color{scInkM}}
\providecommand{\scSmall}{\scArial\fontsize{6}{7.2}\selectfont\color{scInkM}}
\providecommand{\scSub}[1]{{\scSmall\itshape #1}}
\providecommand{\scPanel}[1]{{\scTitle(#1)}}
% Numerals that carry a result: ink, not colour, so the eye reads the value.
\providecommand{\scNum}{\scArial\fontsize{10}{11}\selectfont\bfseries\color{scInk}}

% --- shared TikZ styles ---------------------------------------------------
\tikzset{
  % text never hyphenates inside a figure
  every node/.append style={execute at begin node={%
      \hyphenpenalty=10000 \exhyphenpenalty=10000 \hbadness=10000 }},
  syn/.style={font=\scArial,color=scInk,
              execute at begin picture={\MTversion*{synarial}},
              >={Stealth[round,length=2.4mm,width=1.9mm]}},
  % containers ------------------------------------------------------------
  % One card chrome for the whole paper: white, 3 pt radius, one hairline.
  panel/.style={rounded corners=2.5pt,fill=none,draw=scRule,line width=0.45pt},
  card/.style={rounded corners=2.2pt,fill=scPaper,draw=scRule,line width=0.45pt,
               inner sep=5pt,align=left,anchor=north west},
  % a quiet strip for a panel's closing fact
  note/.style={rounded corners=2.5pt,fill=scPaper,draw=scRuleL,line width=0.4pt,
               inner xsep=6pt,inner ysep=3.5pt,align=left,font=\scSmall},
  % small labelled value pill
  chip/.style={rounded corners=2pt,fill=scPaper,draw=scRule,line width=0.4pt,
               inner xsep=4pt,inner ysep=2pt,font=\scSmall},
  % Circular letter badge: the same accessible panel marker used in Style/*.
  hbullet/.style={circle,fill=scInk,text=white,draw=none,
                  font=\scArial\fontsize{6}{6}\selectfont\bfseries,
                  minimum size=4.2mm,inner sep=0pt},
  ctag/.style={font=\scTitle,anchor=north west,text=scInk},
  % CRN primitives --------------------------------------------------------
  species/.style={rounded rectangle,fill=scSpF,draw=scSpE,line width=0.45pt,
                  inner xsep=2.4pt,inner ysep=1.6pt,font=\scSmall,
                  text=scInk,minimum height=3.6mm},
  cofactor/.style={species,fill=scRefF,draw=scRef},
  catalyst/.style={species,fill=scHiF,draw=scHiE},
  rxn/.style={rectangle,fill=scRxF,draw=scRxE,line width=0.45pt,
              inner sep=1.2pt,font=\scSmall,text=scInk,
              minimum width=3.0mm,minimum height=3.0mm},
  place/.style={circle,fill=scSpF,draw=scSpE,line width=0.45pt,
                minimum size=4.0mm,inner sep=0pt},
  trans/.style={rectangle,fill=scRxE,draw=scRxE,minimum width=1.1mm,
                minimum height=3.6mm,inner sep=0pt},
  token/.style={circle,fill=scInk,minimum size=1.1mm,inner sep=0pt},
  % edges -----------------------------------------------------------------
  reac/.style={->,line width=0.5pt,draw=scReac,shorten >=0.8pt,shorten <=0.8pt},
  prod/.style={->,line width=0.5pt,draw=scProd,shorten >=0.8pt,shorten <=0.8pt},
  cat/.style={->,line width=0.5pt,draw=scHiE,shorten >=0.8pt,shorten <=0.8pt},
  flow/.style={-{Stealth[round,length=2.2mm,width=1.8mm]},line width=1.1pt,
               draw=scInkM!45},
  flab/.style={font=\scSmall,text=scInkM,inner sep=2pt},
  stoi/.style={font=\scSmall,text=scInkM,inner sep=0.8pt},
  % verdict marks: shape carries the meaning, colour only reinforces it
  vpass/.style={circle,fill=scAccentF,draw=scAccent,text=scAccent,
                font=\fontsize{6}{6}\selectfont\bfseries,
                minimum size=3.2mm,inner sep=0pt},
  vfail/.style={circle,fill=scAlertF,draw=scAlert,text=scAlert,
                font=\fontsize{6}{6}\selectfont\bfseries,
                minimum size=3.2mm,inner sep=0pt},
  callout/.style={font=\scSmall,text=scInkM,align=left},
}

% \scHead[<badge colour>]{<anchor coord>}{<letter>}{<title>}{<subtitle or empty>}
% One header treatment for every panel in the paper.
\providecommand{\scHead}[5][scInk]{%
  \node[hbullet,fill=#1,anchor=north west] (scbadge#3) at ($#2+(0,-0.015)$)
    {\MakeUppercase{#3}};
  \node[ctag,anchor=west] at ($(scbadge#3.east)+(0.14,0)$) {#4};
  \ifx\relax#5\relax\else
    \node[anchor=north west,font=\scSmall,text=scInkM]
      at ($#2+(0.56,-0.245)$) {#5};
  \fi}

% --- shared pgfplots look (only if the figure loaded pgfplots) ------------
\ifdefined\pgfplotsset
  \pgfplotsset{compat=1.18}
  \pgfplotsset{
    synaxis/.style={
      scale only axis,
      tick align=outside, tick pos=left,
      axis line style={draw=scInkM!55,line width=0.4pt},
      % One grid, on the value axis only.  No minor grid: at these sizes it
      % turns the plot frame into a grey wash.
      ymajorgrids=true, xmajorgrids=false,
      major grid style={draw=scRuleL,line width=0.3pt},
      major tick length=1.6pt,
      label style={font=\scBody,color=scInk},
      tick label style={font=\scSmall,color=scInkM},
      legend style={font=\scSmall,draw=none,fill=none,inner sep=1pt,
                    row sep=0.4pt,legend columns=1},
      legend cell align=left,
      every axis plot/.append style={line width=0.7pt,mark size=1.2pt},
    },
  }
\fi
}{% ---------------------------------------------------------------------------
% synstyle.tex --- shared look for every TikZ figure of the SynKit CRN paper.
%
% Input it from a figure file:
%     \IfFileExists{synstyle.tex}{\input{synstyle}}{\input{Figure/synstyle}}
% so that the figure compiles both from Figure/ and from the paper root.
%
% The visual vocabulary follows ../../../../Style/style.tex: the
% colour-blind-safe Okabe--Ito accents, quiet neutral rules and wash, circular
% panel badges, and a compact technical type scale are shared by the companion
% papers.  Semantic aliases below keep the CRN meaning explicit.
%
% DESIGN CONTRACT -----------------------------------------------------------
% 1. Geometry.  Every figure is laid out on a canvas exactly \scW wide (the
%    sn-jnl text width, 372 pt) and is included unscaled, so one point of type
%    is the same physical size in all five figures.  Open each picture with
%    \scCanvas{<height in cm>} to pin the bounding box.
% 2. Colour is semantic, never decorative.  Four roles only:
%       scAccent (green)       structure, the normal/expected path
%       scMark   (orange)      catalyst, certificate, or highlight
%       scAlert  (vermillion)  failure, violation, or uncovered set
%       scRef    (blue)        secondary or reference series
%    A panel gets at most one accent.  Titles are always ink; panel identity
%    is also encoded by a letter inside a circular badge, not by colour alone.
% 3. Type scale.  Three sizes, no others, and nothing below 6 pt:
%       \scTitle (7 pt bold)  panel titles
%       \scBody  (6.5 pt)     labels, values, in-figure prose
%       \scSmall (6 pt)       units, footnotes, dense graph labels
% ---------------------------------------------------------------------------
% Load guard.  The paper inputs every figure source inline, so this file is
% reached once per figure; without the guard the second \newlength below
% would abort the run.
\ifdefined\SynKitStyleLoaded\endinput\fi
\newcommand{\SynKitStyleLoaded}{}

% Load in the document preamble. XeLaTeX/LuaLaTeX and an installed Arial
% family are required; fontspec reports an error instead of substituting.
\RequirePackage[no-math]{fontspec}
\newfontfamily\scArial{Arial}[NFSSFamily=scArial,Ligatures=TeX]
% Keep manuscript mathematics unchanged; select Arial letters and numerals
% only inside figures. TeX supplies specialist mathematical symbols.
\RequirePackage[subdued,italic]{mathastext}
\MTDeclareVersion{synarial}{TU}{scArial}{m}{n}
\SetMathAlphabet{\Mathsf}{synarial}{TU}{scArial}{m}{n}
\SetMathAlphabet{\Mathtt}{synarial}{TU}{scArial}{m}{n}

\usetikzlibrary{positioning,calc,fit,backgrounds,arrows.meta,shapes.geometric,
                shapes.misc,decorations.pathreplacing,patterns,math}

% --- canvas ---------------------------------------------------------------
% 372 pt text width minus the 1 pt standalone border on each side.
\newlength{\scW}\setlength{\scW}{13.0326cm}
\providecommand{\scWnum}{13.0326}
% \scCanvas{<height>} pins the picture to the full text width, y running down
% from 0, so every figure shares one coordinate convention.  It *sets* the
% bounding box rather than contributing to it, so the emitted page is exactly
% the text width and every figure is included at 1:1 -- otherwise a stray
% invisible extent (pgfplots reserves a little around an axis) rescales one
% figure relative to the others and the same \scSmall prints at two sizes.
% Anything drawn outside the canvas is therefore cropped: keep content inside,
% and check with
%     \providecommand{\scCanvas}[1]{\path[draw=red] (0,0) rectangle (\scWnum,-#1);}
\providecommand{\scCanvas}[1]{%
  \path[use as bounding box] (0,0) rectangle (\scWnum,-#1);}

% --- palette --------------------------------------------------------------
% Canonical companion-paper values (../../../../Style/style.tex).
\definecolor{npgOrange}{HTML}{E69F00}
\definecolor{npgSky}{HTML}{56B4E9}
\definecolor{npgGreen}{HTML}{009E73}
\definecolor{npgYellow}{HTML}{F0E442}
\definecolor{npgBlue}{HTML}{0072B2}
\definecolor{npgVermillion}{HTML}{E64B35}
\definecolor{npgPurple}{HTML}{CC79A7}
\definecolor{figInk}{HTML}{1A1A1A}
\definecolor{figMuted}{HTML}{666666}
\definecolor{figHair}{HTML}{D9D9D9}
\definecolor{figWash}{HTML}{F7F7F7}

\colorlet{scInk}{figInk}              % primary text
\colorlet{scInkM}{figMuted}           % secondary text
\colorlet{scPaper}{figWash}           % quiet cards and note strips
\colorlet{scSpF}{npgSky!12}           % species fill
\colorlet{scSpE}{npgBlue}             % species edge
\colorlet{scRxF}{npgOrange!12}        % reaction fill
\colorlet{scRxE}{npgOrange!88!black}  % reaction edge
\colorlet{scReac}{npgBlue!78!black}   % reactant edge
\colorlet{scProd}{npgGreen!82!black}  % product edge
\colorlet{scHiF}{npgOrange!13}        % highlight fill
\colorlet{scHiE}{npgOrange!82!black}  % highlight edge
% Stage / accent tints.  stDE remains available for the fifth plot series.
\colorlet{stAF}{npgGreen!9}\colorlet{stAE}{npgGreen}
\colorlet{stBF}{npgOrange!11}\colorlet{stBE}{npgOrange}
\colorlet{stCF}{npgBlue!9}\colorlet{stCE}{npgBlue}
\colorlet{stDF}{npgPurple!10}\colorlet{stDE}{npgPurple!85!black}
\colorlet{stEF}{npgVermillion!9}\colorlet{stEE}{npgVermillion}

% Semantic aliases.  Figures should speak in these, not in stAE/stBE/...
\colorlet{scAccent}{stAE}\colorlet{scAccentF}{stAF}   % structure / normal
\colorlet{scMark}{stBE}\colorlet{scMarkF}{stBF}       % highlight / certificate
\colorlet{scRef}{stCE}\colorlet{scRefF}{stCF}         % secondary series
\colorlet{scAlert}{stEE}\colorlet{scAlertF}{stEF}     % failure / violation
\colorlet{scRule}{figHair}                            % card and axis hairlines
\colorlet{scRuleL}{figHair!55}                        % grid, dividers

% --- text styles ----------------------------------------------------------
% Three sizes, fixed in points so they do not drift with the document class.
% Arial is shared by labels, headings, panel badges, and chemistry labels.
\providecommand{\scTitle}{\scArial\fontsize{7}{8.4}\selectfont\bfseries\color{scInk}}
\providecommand{\scBody}{\scArial\fontsize{6.5}{7.8}\selectfont\color{scInkM}}
\providecommand{\scSmall}{\scArial\fontsize{6}{7.2}\selectfont\color{scInkM}}
\providecommand{\scSub}[1]{{\scSmall\itshape #1}}
\providecommand{\scPanel}[1]{{\scTitle(#1)}}
% Numerals that carry a result: ink, not colour, so the eye reads the value.
\providecommand{\scNum}{\scArial\fontsize{10}{11}\selectfont\bfseries\color{scInk}}

% --- shared TikZ styles ---------------------------------------------------
\tikzset{
  % text never hyphenates inside a figure
  every node/.append style={execute at begin node={%
      \hyphenpenalty=10000 \exhyphenpenalty=10000 \hbadness=10000 }},
  syn/.style={font=\scArial,color=scInk,
              execute at begin picture={\MTversion*{synarial}},
              >={Stealth[round,length=2.4mm,width=1.9mm]}},
  % containers ------------------------------------------------------------
  % One card chrome for the whole paper: white, 3 pt radius, one hairline.
  panel/.style={rounded corners=2.5pt,fill=none,draw=scRule,line width=0.45pt},
  card/.style={rounded corners=2.2pt,fill=scPaper,draw=scRule,line width=0.45pt,
               inner sep=5pt,align=left,anchor=north west},
  % a quiet strip for a panel's closing fact
  note/.style={rounded corners=2.5pt,fill=scPaper,draw=scRuleL,line width=0.4pt,
               inner xsep=6pt,inner ysep=3.5pt,align=left,font=\scSmall},
  % small labelled value pill
  chip/.style={rounded corners=2pt,fill=scPaper,draw=scRule,line width=0.4pt,
               inner xsep=4pt,inner ysep=2pt,font=\scSmall},
  % Circular letter badge: the same accessible panel marker used in Style/*.
  hbullet/.style={circle,fill=scInk,text=white,draw=none,
                  font=\scArial\fontsize{6}{6}\selectfont\bfseries,
                  minimum size=4.2mm,inner sep=0pt},
  ctag/.style={font=\scTitle,anchor=north west,text=scInk},
  % CRN primitives --------------------------------------------------------
  species/.style={rounded rectangle,fill=scSpF,draw=scSpE,line width=0.45pt,
                  inner xsep=2.4pt,inner ysep=1.6pt,font=\scSmall,
                  text=scInk,minimum height=3.6mm},
  cofactor/.style={species,fill=scRefF,draw=scRef},
  catalyst/.style={species,fill=scHiF,draw=scHiE},
  rxn/.style={rectangle,fill=scRxF,draw=scRxE,line width=0.45pt,
              inner sep=1.2pt,font=\scSmall,text=scInk,
              minimum width=3.0mm,minimum height=3.0mm},
  place/.style={circle,fill=scSpF,draw=scSpE,line width=0.45pt,
                minimum size=4.0mm,inner sep=0pt},
  trans/.style={rectangle,fill=scRxE,draw=scRxE,minimum width=1.1mm,
                minimum height=3.6mm,inner sep=0pt},
  token/.style={circle,fill=scInk,minimum size=1.1mm,inner sep=0pt},
  % edges -----------------------------------------------------------------
  reac/.style={->,line width=0.5pt,draw=scReac,shorten >=0.8pt,shorten <=0.8pt},
  prod/.style={->,line width=0.5pt,draw=scProd,shorten >=0.8pt,shorten <=0.8pt},
  cat/.style={->,line width=0.5pt,draw=scHiE,shorten >=0.8pt,shorten <=0.8pt},
  flow/.style={-{Stealth[round,length=2.2mm,width=1.8mm]},line width=1.1pt,
               draw=scInkM!45},
  flab/.style={font=\scSmall,text=scInkM,inner sep=2pt},
  stoi/.style={font=\scSmall,text=scInkM,inner sep=0.8pt},
  % verdict marks: shape carries the meaning, colour only reinforces it
  vpass/.style={circle,fill=scAccentF,draw=scAccent,text=scAccent,
                font=\fontsize{6}{6}\selectfont\bfseries,
                minimum size=3.2mm,inner sep=0pt},
  vfail/.style={circle,fill=scAlertF,draw=scAlert,text=scAlert,
                font=\fontsize{6}{6}\selectfont\bfseries,
                minimum size=3.2mm,inner sep=0pt},
  callout/.style={font=\scSmall,text=scInkM,align=left},
}

% \scHead[<badge colour>]{<anchor coord>}{<letter>}{<title>}{<subtitle or empty>}
% One header treatment for every panel in the paper.
\providecommand{\scHead}[5][scInk]{%
  \node[hbullet,fill=#1,anchor=north west] (scbadge#3) at ($#2+(0,-0.015)$)
    {\MakeUppercase{#3}};
  \node[ctag,anchor=west] at ($(scbadge#3.east)+(0.14,0)$) {#4};
  \ifx\relax#5\relax\else
    \node[anchor=north west,font=\scSmall,text=scInkM]
      at ($#2+(0.56,-0.245)$) {#5};
  \fi}

% --- shared pgfplots look (only if the figure loaded pgfplots) ------------
\ifdefined\pgfplotsset
  \pgfplotsset{compat=1.18}
  \pgfplotsset{
    synaxis/.style={
      scale only axis,
      tick align=outside, tick pos=left,
      axis line style={draw=scInkM!55,line width=0.4pt},
      % One grid, on the value axis only.  No minor grid: at these sizes it
      % turns the plot frame into a grey wash.
      ymajorgrids=true, xmajorgrids=false,
      major grid style={draw=scRuleL,line width=0.3pt},
      major tick length=1.6pt,
      label style={font=\scBody,color=scInk},
      tick label style={font=\scSmall,color=scInkM},
      legend style={font=\scSmall,draw=none,fill=none,inner sep=1pt,
                    row sep=0.4pt,legend columns=1},
      legend cell align=left,
      every axis plot/.append style={line width=0.7pt,mark size=1.2pt},
    },
  }
\fi
}
% so that the figure compiles both from Figure/ and from the paper root.
%
% The visual vocabulary follows ../../../../Style/style.tex: the
% colour-blind-safe Okabe--Ito accents, quiet neutral rules and wash, circular
% panel badges, and a compact technical type scale are shared by the companion
% papers.  Semantic aliases below keep the CRN meaning explicit.
%
% DESIGN CONTRACT -----------------------------------------------------------
% 1. Geometry.  Every figure is laid out on a canvas exactly \scW wide (the
%    sn-jnl text width, 372 pt) and is included unscaled, so one point of type
%    is the same physical size in all five figures.  Open each picture with
%    \scCanvas{<height in cm>} to pin the bounding box.
% 2. Colour is semantic, never decorative.  Four roles only:
%       scAccent (green)       structure, the normal/expected path
%       scMark   (orange)      catalyst, certificate, or highlight
%       scAlert  (vermillion)  failure, violation, or uncovered set
%       scRef    (blue)        secondary or reference series
%    A panel gets at most one accent.  Titles are always ink; panel identity
%    is also encoded by a letter inside a circular badge, not by colour alone.
% 3. Type scale.  Three sizes, no others, and nothing below 6 pt:
%       \scTitle (7 pt bold)  panel titles
%       \scBody  (6.5 pt)     labels, values, in-figure prose
%       \scSmall (6 pt)       units, footnotes, dense graph labels
% ---------------------------------------------------------------------------
% Load guard.  The paper inputs every figure source inline, so this file is
% reached once per figure; without the guard the second \newlength below
% would abort the run.
\ifdefined\SynKitStyleLoaded\endinput\fi
\newcommand{\SynKitStyleLoaded}{}

% Load in the document preamble. XeLaTeX/LuaLaTeX and an installed Arial
% family are required; fontspec reports an error instead of substituting.
\RequirePackage[no-math]{fontspec}
\newfontfamily\scArial{Arial}[NFSSFamily=scArial,Ligatures=TeX]
% Keep manuscript mathematics unchanged; select Arial letters and numerals
% only inside figures. TeX supplies specialist mathematical symbols.
\RequirePackage[subdued,italic]{mathastext}
\MTDeclareVersion{synarial}{TU}{scArial}{m}{n}
\SetMathAlphabet{\Mathsf}{synarial}{TU}{scArial}{m}{n}
\SetMathAlphabet{\Mathtt}{synarial}{TU}{scArial}{m}{n}

\usetikzlibrary{positioning,calc,fit,backgrounds,arrows.meta,shapes.geometric,
                shapes.misc,decorations.pathreplacing,patterns,math}

% --- canvas ---------------------------------------------------------------
% 372 pt text width minus the 1 pt standalone border on each side.
\newlength{\scW}\setlength{\scW}{13.0326cm}
\providecommand{\scWnum}{13.0326}
% \scCanvas{<height>} pins the picture to the full text width, y running down
% from 0, so every figure shares one coordinate convention.  It *sets* the
% bounding box rather than contributing to it, so the emitted page is exactly
% the text width and every figure is included at 1:1 -- otherwise a stray
% invisible extent (pgfplots reserves a little around an axis) rescales one
% figure relative to the others and the same \scSmall prints at two sizes.
% Anything drawn outside the canvas is therefore cropped: keep content inside,
% and check with
%     \providecommand{\scCanvas}[1]{\path[draw=red] (0,0) rectangle (\scWnum,-#1);}
\providecommand{\scCanvas}[1]{%
  \path[use as bounding box] (0,0) rectangle (\scWnum,-#1);}

% --- palette --------------------------------------------------------------
% Canonical companion-paper values (../../../../Style/style.tex).
\definecolor{npgOrange}{HTML}{E69F00}
\definecolor{npgSky}{HTML}{56B4E9}
\definecolor{npgGreen}{HTML}{009E73}
\definecolor{npgYellow}{HTML}{F0E442}
\definecolor{npgBlue}{HTML}{0072B2}
\definecolor{npgVermillion}{HTML}{E64B35}
\definecolor{npgPurple}{HTML}{CC79A7}
\definecolor{figInk}{HTML}{1A1A1A}
\definecolor{figMuted}{HTML}{666666}
\definecolor{figHair}{HTML}{D9D9D9}
\definecolor{figWash}{HTML}{F7F7F7}

\colorlet{scInk}{figInk}              % primary text
\colorlet{scInkM}{figMuted}           % secondary text
\colorlet{scPaper}{figWash}           % quiet cards and note strips
\colorlet{scSpF}{npgSky!12}           % species fill
\colorlet{scSpE}{npgBlue}             % species edge
\colorlet{scRxF}{npgOrange!12}        % reaction fill
\colorlet{scRxE}{npgOrange!88!black}  % reaction edge
\colorlet{scReac}{npgBlue!78!black}   % reactant edge
\colorlet{scProd}{npgGreen!82!black}  % product edge
\colorlet{scHiF}{npgOrange!13}        % highlight fill
\colorlet{scHiE}{npgOrange!82!black}  % highlight edge
% Stage / accent tints.  stDE remains available for the fifth plot series.
\colorlet{stAF}{npgGreen!9}\colorlet{stAE}{npgGreen}
\colorlet{stBF}{npgOrange!11}\colorlet{stBE}{npgOrange}
\colorlet{stCF}{npgBlue!9}\colorlet{stCE}{npgBlue}
\colorlet{stDF}{npgPurple!10}\colorlet{stDE}{npgPurple!85!black}
\colorlet{stEF}{npgVermillion!9}\colorlet{stEE}{npgVermillion}

% Semantic aliases.  Figures should speak in these, not in stAE/stBE/...
\colorlet{scAccent}{stAE}\colorlet{scAccentF}{stAF}   % structure / normal
\colorlet{scMark}{stBE}\colorlet{scMarkF}{stBF}       % highlight / certificate
\colorlet{scRef}{stCE}\colorlet{scRefF}{stCF}         % secondary series
\colorlet{scAlert}{stEE}\colorlet{scAlertF}{stEF}     % failure / violation
\colorlet{scRule}{figHair}                            % card and axis hairlines
\colorlet{scRuleL}{figHair!55}                        % grid, dividers

% --- text styles ----------------------------------------------------------
% Three sizes, fixed in points so they do not drift with the document class.
% Arial is shared by labels, headings, panel badges, and chemistry labels.
\providecommand{\scTitle}{\scArial\fontsize{7}{8.4}\selectfont\bfseries\color{scInk}}
\providecommand{\scBody}{\scArial\fontsize{6.5}{7.8}\selectfont\color{scInkM}}
\providecommand{\scSmall}{\scArial\fontsize{6}{7.2}\selectfont\color{scInkM}}
\providecommand{\scSub}[1]{{\scSmall\itshape #1}}
\providecommand{\scPanel}[1]{{\scTitle(#1)}}
% Numerals that carry a result: ink, not colour, so the eye reads the value.
\providecommand{\scNum}{\scArial\fontsize{10}{11}\selectfont\bfseries\color{scInk}}

% --- shared TikZ styles ---------------------------------------------------
\tikzset{
  % text never hyphenates inside a figure
  every node/.append style={execute at begin node={%
      \hyphenpenalty=10000 \exhyphenpenalty=10000 \hbadness=10000 }},
  syn/.style={font=\scArial,color=scInk,
              execute at begin picture={\MTversion*{synarial}},
              >={Stealth[round,length=2.4mm,width=1.9mm]}},
  % containers ------------------------------------------------------------
  % One card chrome for the whole paper: white, 3 pt radius, one hairline.
  panel/.style={rounded corners=2.5pt,fill=none,draw=scRule,line width=0.45pt},
  card/.style={rounded corners=2.2pt,fill=scPaper,draw=scRule,line width=0.45pt,
               inner sep=5pt,align=left,anchor=north west},
  % a quiet strip for a panel's closing fact
  note/.style={rounded corners=2.5pt,fill=scPaper,draw=scRuleL,line width=0.4pt,
               inner xsep=6pt,inner ysep=3.5pt,align=left,font=\scSmall},
  % small labelled value pill
  chip/.style={rounded corners=2pt,fill=scPaper,draw=scRule,line width=0.4pt,
               inner xsep=4pt,inner ysep=2pt,font=\scSmall},
  % Circular letter badge: the same accessible panel marker used in Style/*.
  hbullet/.style={circle,fill=scInk,text=white,draw=none,
                  font=\scArial\fontsize{6}{6}\selectfont\bfseries,
                  minimum size=4.2mm,inner sep=0pt},
  ctag/.style={font=\scTitle,anchor=north west,text=scInk},
  % CRN primitives --------------------------------------------------------
  species/.style={rounded rectangle,fill=scSpF,draw=scSpE,line width=0.45pt,
                  inner xsep=2.4pt,inner ysep=1.6pt,font=\scSmall,
                  text=scInk,minimum height=3.6mm},
  cofactor/.style={species,fill=scRefF,draw=scRef},
  catalyst/.style={species,fill=scHiF,draw=scHiE},
  rxn/.style={rectangle,fill=scRxF,draw=scRxE,line width=0.45pt,
              inner sep=1.2pt,font=\scSmall,text=scInk,
              minimum width=3.0mm,minimum height=3.0mm},
  place/.style={circle,fill=scSpF,draw=scSpE,line width=0.45pt,
                minimum size=4.0mm,inner sep=0pt},
  trans/.style={rectangle,fill=scRxE,draw=scRxE,minimum width=1.1mm,
                minimum height=3.6mm,inner sep=0pt},
  token/.style={circle,fill=scInk,minimum size=1.1mm,inner sep=0pt},
  % edges -----------------------------------------------------------------
  reac/.style={->,line width=0.5pt,draw=scReac,shorten >=0.8pt,shorten <=0.8pt},
  prod/.style={->,line width=0.5pt,draw=scProd,shorten >=0.8pt,shorten <=0.8pt},
  cat/.style={->,line width=0.5pt,draw=scHiE,shorten >=0.8pt,shorten <=0.8pt},
  flow/.style={-{Stealth[round,length=2.2mm,width=1.8mm]},line width=1.1pt,
               draw=scInkM!45},
  flab/.style={font=\scSmall,text=scInkM,inner sep=2pt},
  stoi/.style={font=\scSmall,text=scInkM,inner sep=0.8pt},
  % verdict marks: shape carries the meaning, colour only reinforces it
  vpass/.style={circle,fill=scAccentF,draw=scAccent,text=scAccent,
                font=\fontsize{6}{6}\selectfont\bfseries,
                minimum size=3.2mm,inner sep=0pt},
  vfail/.style={circle,fill=scAlertF,draw=scAlert,text=scAlert,
                font=\fontsize{6}{6}\selectfont\bfseries,
                minimum size=3.2mm,inner sep=0pt},
  callout/.style={font=\scSmall,text=scInkM,align=left},
}

% \scHead[<badge colour>]{<anchor coord>}{<letter>}{<title>}{<subtitle or empty>}
% One header treatment for every panel in the paper.
\providecommand{\scHead}[5][scInk]{%
  \node[hbullet,fill=#1,anchor=north west] (scbadge#3) at ($#2+(0,-0.015)$)
    {\MakeUppercase{#3}};
  \node[ctag,anchor=west] at ($(scbadge#3.east)+(0.14,0)$) {#4};
  \ifx\relax#5\relax\else
    \node[anchor=north west,font=\scSmall,text=scInkM]
      at ($#2+(0.56,-0.245)$) {#5};
  \fi}

% --- shared pgfplots look (only if the figure loaded pgfplots) ------------
\ifdefined\pgfplotsset
  \pgfplotsset{compat=1.18}
  \pgfplotsset{
    synaxis/.style={
      scale only axis,
      tick align=outside, tick pos=left,
      axis line style={draw=scInkM!55,line width=0.4pt},
      % One grid, on the value axis only.  No minor grid: at these sizes it
      % turns the plot frame into a grey wash.
      ymajorgrids=true, xmajorgrids=false,
      major grid style={draw=scRuleL,line width=0.3pt},
      major tick length=1.6pt,
      label style={font=\scBody,color=scInk},
      tick label style={font=\scSmall,color=scInkM},
      legend style={font=\scSmall,draw=none,fill=none,inner sep=1pt,
                    row sep=0.4pt,legend columns=1},
      legend cell align=left,
      every axis plot/.append style={line width=0.7pt,mark size=1.2pt},
    },
  }
\fi
}{% ---------------------------------------------------------------------------
% synstyle.tex --- shared look for every TikZ figure of the SynKit CRN paper.
%
% Input it from a figure file:
%     \IfFileExists{synstyle.tex}{% ---------------------------------------------------------------------------
% synstyle.tex --- shared look for every TikZ figure of the SynKit CRN paper.
%
% Input it from a figure file:
%     \IfFileExists{synstyle.tex}{\input{synstyle}}{\input{Figure/synstyle}}
% so that the figure compiles both from Figure/ and from the paper root.
%
% The visual vocabulary follows ../../../../Style/style.tex: the
% colour-blind-safe Okabe--Ito accents, quiet neutral rules and wash, circular
% panel badges, and a compact technical type scale are shared by the companion
% papers.  Semantic aliases below keep the CRN meaning explicit.
%
% DESIGN CONTRACT -----------------------------------------------------------
% 1. Geometry.  Every figure is laid out on a canvas exactly \scW wide (the
%    sn-jnl text width, 372 pt) and is included unscaled, so one point of type
%    is the same physical size in all five figures.  Open each picture with
%    \scCanvas{<height in cm>} to pin the bounding box.
% 2. Colour is semantic, never decorative.  Four roles only:
%       scAccent (green)       structure, the normal/expected path
%       scMark   (orange)      catalyst, certificate, or highlight
%       scAlert  (vermillion)  failure, violation, or uncovered set
%       scRef    (blue)        secondary or reference series
%    A panel gets at most one accent.  Titles are always ink; panel identity
%    is also encoded by a letter inside a circular badge, not by colour alone.
% 3. Type scale.  Three sizes, no others, and nothing below 6 pt:
%       \scTitle (7 pt bold)  panel titles
%       \scBody  (6.5 pt)     labels, values, in-figure prose
%       \scSmall (6 pt)       units, footnotes, dense graph labels
% ---------------------------------------------------------------------------
% Load guard.  The paper inputs every figure source inline, so this file is
% reached once per figure; without the guard the second \newlength below
% would abort the run.
\ifdefined\SynKitStyleLoaded\endinput\fi
\newcommand{\SynKitStyleLoaded}{}

% Load in the document preamble. XeLaTeX/LuaLaTeX and an installed Arial
% family are required; fontspec reports an error instead of substituting.
\RequirePackage[no-math]{fontspec}
\newfontfamily\scArial{Arial}[NFSSFamily=scArial,Ligatures=TeX]
% Keep manuscript mathematics unchanged; select Arial letters and numerals
% only inside figures. TeX supplies specialist mathematical symbols.
\RequirePackage[subdued,italic]{mathastext}
\MTDeclareVersion{synarial}{TU}{scArial}{m}{n}
\SetMathAlphabet{\Mathsf}{synarial}{TU}{scArial}{m}{n}
\SetMathAlphabet{\Mathtt}{synarial}{TU}{scArial}{m}{n}

\usetikzlibrary{positioning,calc,fit,backgrounds,arrows.meta,shapes.geometric,
                shapes.misc,decorations.pathreplacing,patterns,math}

% --- canvas ---------------------------------------------------------------
% 372 pt text width minus the 1 pt standalone border on each side.
\newlength{\scW}\setlength{\scW}{13.0326cm}
\providecommand{\scWnum}{13.0326}
% \scCanvas{<height>} pins the picture to the full text width, y running down
% from 0, so every figure shares one coordinate convention.  It *sets* the
% bounding box rather than contributing to it, so the emitted page is exactly
% the text width and every figure is included at 1:1 -- otherwise a stray
% invisible extent (pgfplots reserves a little around an axis) rescales one
% figure relative to the others and the same \scSmall prints at two sizes.
% Anything drawn outside the canvas is therefore cropped: keep content inside,
% and check with
%     \providecommand{\scCanvas}[1]{\path[draw=red] (0,0) rectangle (\scWnum,-#1);}
\providecommand{\scCanvas}[1]{%
  \path[use as bounding box] (0,0) rectangle (\scWnum,-#1);}

% --- palette --------------------------------------------------------------
% Canonical companion-paper values (../../../../Style/style.tex).
\definecolor{npgOrange}{HTML}{E69F00}
\definecolor{npgSky}{HTML}{56B4E9}
\definecolor{npgGreen}{HTML}{009E73}
\definecolor{npgYellow}{HTML}{F0E442}
\definecolor{npgBlue}{HTML}{0072B2}
\definecolor{npgVermillion}{HTML}{E64B35}
\definecolor{npgPurple}{HTML}{CC79A7}
\definecolor{figInk}{HTML}{1A1A1A}
\definecolor{figMuted}{HTML}{666666}
\definecolor{figHair}{HTML}{D9D9D9}
\definecolor{figWash}{HTML}{F7F7F7}

\colorlet{scInk}{figInk}              % primary text
\colorlet{scInkM}{figMuted}           % secondary text
\colorlet{scPaper}{figWash}           % quiet cards and note strips
\colorlet{scSpF}{npgSky!12}           % species fill
\colorlet{scSpE}{npgBlue}             % species edge
\colorlet{scRxF}{npgOrange!12}        % reaction fill
\colorlet{scRxE}{npgOrange!88!black}  % reaction edge
\colorlet{scReac}{npgBlue!78!black}   % reactant edge
\colorlet{scProd}{npgGreen!82!black}  % product edge
\colorlet{scHiF}{npgOrange!13}        % highlight fill
\colorlet{scHiE}{npgOrange!82!black}  % highlight edge
% Stage / accent tints.  stDE remains available for the fifth plot series.
\colorlet{stAF}{npgGreen!9}\colorlet{stAE}{npgGreen}
\colorlet{stBF}{npgOrange!11}\colorlet{stBE}{npgOrange}
\colorlet{stCF}{npgBlue!9}\colorlet{stCE}{npgBlue}
\colorlet{stDF}{npgPurple!10}\colorlet{stDE}{npgPurple!85!black}
\colorlet{stEF}{npgVermillion!9}\colorlet{stEE}{npgVermillion}

% Semantic aliases.  Figures should speak in these, not in stAE/stBE/...
\colorlet{scAccent}{stAE}\colorlet{scAccentF}{stAF}   % structure / normal
\colorlet{scMark}{stBE}\colorlet{scMarkF}{stBF}       % highlight / certificate
\colorlet{scRef}{stCE}\colorlet{scRefF}{stCF}         % secondary series
\colorlet{scAlert}{stEE}\colorlet{scAlertF}{stEF}     % failure / violation
\colorlet{scRule}{figHair}                            % card and axis hairlines
\colorlet{scRuleL}{figHair!55}                        % grid, dividers

% --- text styles ----------------------------------------------------------
% Three sizes, fixed in points so they do not drift with the document class.
% Arial is shared by labels, headings, panel badges, and chemistry labels.
\providecommand{\scTitle}{\scArial\fontsize{7}{8.4}\selectfont\bfseries\color{scInk}}
\providecommand{\scBody}{\scArial\fontsize{6.5}{7.8}\selectfont\color{scInkM}}
\providecommand{\scSmall}{\scArial\fontsize{6}{7.2}\selectfont\color{scInkM}}
\providecommand{\scSub}[1]{{\scSmall\itshape #1}}
\providecommand{\scPanel}[1]{{\scTitle(#1)}}
% Numerals that carry a result: ink, not colour, so the eye reads the value.
\providecommand{\scNum}{\scArial\fontsize{10}{11}\selectfont\bfseries\color{scInk}}

% --- shared TikZ styles ---------------------------------------------------
\tikzset{
  % text never hyphenates inside a figure
  every node/.append style={execute at begin node={%
      \hyphenpenalty=10000 \exhyphenpenalty=10000 \hbadness=10000 }},
  syn/.style={font=\scArial,color=scInk,
              execute at begin picture={\MTversion*{synarial}},
              >={Stealth[round,length=2.4mm,width=1.9mm]}},
  % containers ------------------------------------------------------------
  % One card chrome for the whole paper: white, 3 pt radius, one hairline.
  panel/.style={rounded corners=2.5pt,fill=none,draw=scRule,line width=0.45pt},
  card/.style={rounded corners=2.2pt,fill=scPaper,draw=scRule,line width=0.45pt,
               inner sep=5pt,align=left,anchor=north west},
  % a quiet strip for a panel's closing fact
  note/.style={rounded corners=2.5pt,fill=scPaper,draw=scRuleL,line width=0.4pt,
               inner xsep=6pt,inner ysep=3.5pt,align=left,font=\scSmall},
  % small labelled value pill
  chip/.style={rounded corners=2pt,fill=scPaper,draw=scRule,line width=0.4pt,
               inner xsep=4pt,inner ysep=2pt,font=\scSmall},
  % Circular letter badge: the same accessible panel marker used in Style/*.
  hbullet/.style={circle,fill=scInk,text=white,draw=none,
                  font=\scArial\fontsize{6}{6}\selectfont\bfseries,
                  minimum size=4.2mm,inner sep=0pt},
  ctag/.style={font=\scTitle,anchor=north west,text=scInk},
  % CRN primitives --------------------------------------------------------
  species/.style={rounded rectangle,fill=scSpF,draw=scSpE,line width=0.45pt,
                  inner xsep=2.4pt,inner ysep=1.6pt,font=\scSmall,
                  text=scInk,minimum height=3.6mm},
  cofactor/.style={species,fill=scRefF,draw=scRef},
  catalyst/.style={species,fill=scHiF,draw=scHiE},
  rxn/.style={rectangle,fill=scRxF,draw=scRxE,line width=0.45pt,
              inner sep=1.2pt,font=\scSmall,text=scInk,
              minimum width=3.0mm,minimum height=3.0mm},
  place/.style={circle,fill=scSpF,draw=scSpE,line width=0.45pt,
                minimum size=4.0mm,inner sep=0pt},
  trans/.style={rectangle,fill=scRxE,draw=scRxE,minimum width=1.1mm,
                minimum height=3.6mm,inner sep=0pt},
  token/.style={circle,fill=scInk,minimum size=1.1mm,inner sep=0pt},
  % edges -----------------------------------------------------------------
  reac/.style={->,line width=0.5pt,draw=scReac,shorten >=0.8pt,shorten <=0.8pt},
  prod/.style={->,line width=0.5pt,draw=scProd,shorten >=0.8pt,shorten <=0.8pt},
  cat/.style={->,line width=0.5pt,draw=scHiE,shorten >=0.8pt,shorten <=0.8pt},
  flow/.style={-{Stealth[round,length=2.2mm,width=1.8mm]},line width=1.1pt,
               draw=scInkM!45},
  flab/.style={font=\scSmall,text=scInkM,inner sep=2pt},
  stoi/.style={font=\scSmall,text=scInkM,inner sep=0.8pt},
  % verdict marks: shape carries the meaning, colour only reinforces it
  vpass/.style={circle,fill=scAccentF,draw=scAccent,text=scAccent,
                font=\fontsize{6}{6}\selectfont\bfseries,
                minimum size=3.2mm,inner sep=0pt},
  vfail/.style={circle,fill=scAlertF,draw=scAlert,text=scAlert,
                font=\fontsize{6}{6}\selectfont\bfseries,
                minimum size=3.2mm,inner sep=0pt},
  callout/.style={font=\scSmall,text=scInkM,align=left},
}

% \scHead[<badge colour>]{<anchor coord>}{<letter>}{<title>}{<subtitle or empty>}
% One header treatment for every panel in the paper.
\providecommand{\scHead}[5][scInk]{%
  \node[hbullet,fill=#1,anchor=north west] (scbadge#3) at ($#2+(0,-0.015)$)
    {\MakeUppercase{#3}};
  \node[ctag,anchor=west] at ($(scbadge#3.east)+(0.14,0)$) {#4};
  \ifx\relax#5\relax\else
    \node[anchor=north west,font=\scSmall,text=scInkM]
      at ($#2+(0.56,-0.245)$) {#5};
  \fi}

% --- shared pgfplots look (only if the figure loaded pgfplots) ------------
\ifdefined\pgfplotsset
  \pgfplotsset{compat=1.18}
  \pgfplotsset{
    synaxis/.style={
      scale only axis,
      tick align=outside, tick pos=left,
      axis line style={draw=scInkM!55,line width=0.4pt},
      % One grid, on the value axis only.  No minor grid: at these sizes it
      % turns the plot frame into a grey wash.
      ymajorgrids=true, xmajorgrids=false,
      major grid style={draw=scRuleL,line width=0.3pt},
      major tick length=1.6pt,
      label style={font=\scBody,color=scInk},
      tick label style={font=\scSmall,color=scInkM},
      legend style={font=\scSmall,draw=none,fill=none,inner sep=1pt,
                    row sep=0.4pt,legend columns=1},
      legend cell align=left,
      every axis plot/.append style={line width=0.7pt,mark size=1.2pt},
    },
  }
\fi
}{% ---------------------------------------------------------------------------
% synstyle.tex --- shared look for every TikZ figure of the SynKit CRN paper.
%
% Input it from a figure file:
%     \IfFileExists{synstyle.tex}{\input{synstyle}}{\input{Figure/synstyle}}
% so that the figure compiles both from Figure/ and from the paper root.
%
% The visual vocabulary follows ../../../../Style/style.tex: the
% colour-blind-safe Okabe--Ito accents, quiet neutral rules and wash, circular
% panel badges, and a compact technical type scale are shared by the companion
% papers.  Semantic aliases below keep the CRN meaning explicit.
%
% DESIGN CONTRACT -----------------------------------------------------------
% 1. Geometry.  Every figure is laid out on a canvas exactly \scW wide (the
%    sn-jnl text width, 372 pt) and is included unscaled, so one point of type
%    is the same physical size in all five figures.  Open each picture with
%    \scCanvas{<height in cm>} to pin the bounding box.
% 2. Colour is semantic, never decorative.  Four roles only:
%       scAccent (green)       structure, the normal/expected path
%       scMark   (orange)      catalyst, certificate, or highlight
%       scAlert  (vermillion)  failure, violation, or uncovered set
%       scRef    (blue)        secondary or reference series
%    A panel gets at most one accent.  Titles are always ink; panel identity
%    is also encoded by a letter inside a circular badge, not by colour alone.
% 3. Type scale.  Three sizes, no others, and nothing below 6 pt:
%       \scTitle (7 pt bold)  panel titles
%       \scBody  (6.5 pt)     labels, values, in-figure prose
%       \scSmall (6 pt)       units, footnotes, dense graph labels
% ---------------------------------------------------------------------------
% Load guard.  The paper inputs every figure source inline, so this file is
% reached once per figure; without the guard the second \newlength below
% would abort the run.
\ifdefined\SynKitStyleLoaded\endinput\fi
\newcommand{\SynKitStyleLoaded}{}

% Load in the document preamble. XeLaTeX/LuaLaTeX and an installed Arial
% family are required; fontspec reports an error instead of substituting.
\RequirePackage[no-math]{fontspec}
\newfontfamily\scArial{Arial}[NFSSFamily=scArial,Ligatures=TeX]
% Keep manuscript mathematics unchanged; select Arial letters and numerals
% only inside figures. TeX supplies specialist mathematical symbols.
\RequirePackage[subdued,italic]{mathastext}
\MTDeclareVersion{synarial}{TU}{scArial}{m}{n}
\SetMathAlphabet{\Mathsf}{synarial}{TU}{scArial}{m}{n}
\SetMathAlphabet{\Mathtt}{synarial}{TU}{scArial}{m}{n}

\usetikzlibrary{positioning,calc,fit,backgrounds,arrows.meta,shapes.geometric,
                shapes.misc,decorations.pathreplacing,patterns,math}

% --- canvas ---------------------------------------------------------------
% 372 pt text width minus the 1 pt standalone border on each side.
\newlength{\scW}\setlength{\scW}{13.0326cm}
\providecommand{\scWnum}{13.0326}
% \scCanvas{<height>} pins the picture to the full text width, y running down
% from 0, so every figure shares one coordinate convention.  It *sets* the
% bounding box rather than contributing to it, so the emitted page is exactly
% the text width and every figure is included at 1:1 -- otherwise a stray
% invisible extent (pgfplots reserves a little around an axis) rescales one
% figure relative to the others and the same \scSmall prints at two sizes.
% Anything drawn outside the canvas is therefore cropped: keep content inside,
% and check with
%     \providecommand{\scCanvas}[1]{\path[draw=red] (0,0) rectangle (\scWnum,-#1);}
\providecommand{\scCanvas}[1]{%
  \path[use as bounding box] (0,0) rectangle (\scWnum,-#1);}

% --- palette --------------------------------------------------------------
% Canonical companion-paper values (../../../../Style/style.tex).
\definecolor{npgOrange}{HTML}{E69F00}
\definecolor{npgSky}{HTML}{56B4E9}
\definecolor{npgGreen}{HTML}{009E73}
\definecolor{npgYellow}{HTML}{F0E442}
\definecolor{npgBlue}{HTML}{0072B2}
\definecolor{npgVermillion}{HTML}{E64B35}
\definecolor{npgPurple}{HTML}{CC79A7}
\definecolor{figInk}{HTML}{1A1A1A}
\definecolor{figMuted}{HTML}{666666}
\definecolor{figHair}{HTML}{D9D9D9}
\definecolor{figWash}{HTML}{F7F7F7}

\colorlet{scInk}{figInk}              % primary text
\colorlet{scInkM}{figMuted}           % secondary text
\colorlet{scPaper}{figWash}           % quiet cards and note strips
\colorlet{scSpF}{npgSky!12}           % species fill
\colorlet{scSpE}{npgBlue}             % species edge
\colorlet{scRxF}{npgOrange!12}        % reaction fill
\colorlet{scRxE}{npgOrange!88!black}  % reaction edge
\colorlet{scReac}{npgBlue!78!black}   % reactant edge
\colorlet{scProd}{npgGreen!82!black}  % product edge
\colorlet{scHiF}{npgOrange!13}        % highlight fill
\colorlet{scHiE}{npgOrange!82!black}  % highlight edge
% Stage / accent tints.  stDE remains available for the fifth plot series.
\colorlet{stAF}{npgGreen!9}\colorlet{stAE}{npgGreen}
\colorlet{stBF}{npgOrange!11}\colorlet{stBE}{npgOrange}
\colorlet{stCF}{npgBlue!9}\colorlet{stCE}{npgBlue}
\colorlet{stDF}{npgPurple!10}\colorlet{stDE}{npgPurple!85!black}
\colorlet{stEF}{npgVermillion!9}\colorlet{stEE}{npgVermillion}

% Semantic aliases.  Figures should speak in these, not in stAE/stBE/...
\colorlet{scAccent}{stAE}\colorlet{scAccentF}{stAF}   % structure / normal
\colorlet{scMark}{stBE}\colorlet{scMarkF}{stBF}       % highlight / certificate
\colorlet{scRef}{stCE}\colorlet{scRefF}{stCF}         % secondary series
\colorlet{scAlert}{stEE}\colorlet{scAlertF}{stEF}     % failure / violation
\colorlet{scRule}{figHair}                            % card and axis hairlines
\colorlet{scRuleL}{figHair!55}                        % grid, dividers

% --- text styles ----------------------------------------------------------
% Three sizes, fixed in points so they do not drift with the document class.
% Arial is shared by labels, headings, panel badges, and chemistry labels.
\providecommand{\scTitle}{\scArial\fontsize{7}{8.4}\selectfont\bfseries\color{scInk}}
\providecommand{\scBody}{\scArial\fontsize{6.5}{7.8}\selectfont\color{scInkM}}
\providecommand{\scSmall}{\scArial\fontsize{6}{7.2}\selectfont\color{scInkM}}
\providecommand{\scSub}[1]{{\scSmall\itshape #1}}
\providecommand{\scPanel}[1]{{\scTitle(#1)}}
% Numerals that carry a result: ink, not colour, so the eye reads the value.
\providecommand{\scNum}{\scArial\fontsize{10}{11}\selectfont\bfseries\color{scInk}}

% --- shared TikZ styles ---------------------------------------------------
\tikzset{
  % text never hyphenates inside a figure
  every node/.append style={execute at begin node={%
      \hyphenpenalty=10000 \exhyphenpenalty=10000 \hbadness=10000 }},
  syn/.style={font=\scArial,color=scInk,
              execute at begin picture={\MTversion*{synarial}},
              >={Stealth[round,length=2.4mm,width=1.9mm]}},
  % containers ------------------------------------------------------------
  % One card chrome for the whole paper: white, 3 pt radius, one hairline.
  panel/.style={rounded corners=2.5pt,fill=none,draw=scRule,line width=0.45pt},
  card/.style={rounded corners=2.2pt,fill=scPaper,draw=scRule,line width=0.45pt,
               inner sep=5pt,align=left,anchor=north west},
  % a quiet strip for a panel's closing fact
  note/.style={rounded corners=2.5pt,fill=scPaper,draw=scRuleL,line width=0.4pt,
               inner xsep=6pt,inner ysep=3.5pt,align=left,font=\scSmall},
  % small labelled value pill
  chip/.style={rounded corners=2pt,fill=scPaper,draw=scRule,line width=0.4pt,
               inner xsep=4pt,inner ysep=2pt,font=\scSmall},
  % Circular letter badge: the same accessible panel marker used in Style/*.
  hbullet/.style={circle,fill=scInk,text=white,draw=none,
                  font=\scArial\fontsize{6}{6}\selectfont\bfseries,
                  minimum size=4.2mm,inner sep=0pt},
  ctag/.style={font=\scTitle,anchor=north west,text=scInk},
  % CRN primitives --------------------------------------------------------
  species/.style={rounded rectangle,fill=scSpF,draw=scSpE,line width=0.45pt,
                  inner xsep=2.4pt,inner ysep=1.6pt,font=\scSmall,
                  text=scInk,minimum height=3.6mm},
  cofactor/.style={species,fill=scRefF,draw=scRef},
  catalyst/.style={species,fill=scHiF,draw=scHiE},
  rxn/.style={rectangle,fill=scRxF,draw=scRxE,line width=0.45pt,
              inner sep=1.2pt,font=\scSmall,text=scInk,
              minimum width=3.0mm,minimum height=3.0mm},
  place/.style={circle,fill=scSpF,draw=scSpE,line width=0.45pt,
                minimum size=4.0mm,inner sep=0pt},
  trans/.style={rectangle,fill=scRxE,draw=scRxE,minimum width=1.1mm,
                minimum height=3.6mm,inner sep=0pt},
  token/.style={circle,fill=scInk,minimum size=1.1mm,inner sep=0pt},
  % edges -----------------------------------------------------------------
  reac/.style={->,line width=0.5pt,draw=scReac,shorten >=0.8pt,shorten <=0.8pt},
  prod/.style={->,line width=0.5pt,draw=scProd,shorten >=0.8pt,shorten <=0.8pt},
  cat/.style={->,line width=0.5pt,draw=scHiE,shorten >=0.8pt,shorten <=0.8pt},
  flow/.style={-{Stealth[round,length=2.2mm,width=1.8mm]},line width=1.1pt,
               draw=scInkM!45},
  flab/.style={font=\scSmall,text=scInkM,inner sep=2pt},
  stoi/.style={font=\scSmall,text=scInkM,inner sep=0.8pt},
  % verdict marks: shape carries the meaning, colour only reinforces it
  vpass/.style={circle,fill=scAccentF,draw=scAccent,text=scAccent,
                font=\fontsize{6}{6}\selectfont\bfseries,
                minimum size=3.2mm,inner sep=0pt},
  vfail/.style={circle,fill=scAlertF,draw=scAlert,text=scAlert,
                font=\fontsize{6}{6}\selectfont\bfseries,
                minimum size=3.2mm,inner sep=0pt},
  callout/.style={font=\scSmall,text=scInkM,align=left},
}

% \scHead[<badge colour>]{<anchor coord>}{<letter>}{<title>}{<subtitle or empty>}
% One header treatment for every panel in the paper.
\providecommand{\scHead}[5][scInk]{%
  \node[hbullet,fill=#1,anchor=north west] (scbadge#3) at ($#2+(0,-0.015)$)
    {\MakeUppercase{#3}};
  \node[ctag,anchor=west] at ($(scbadge#3.east)+(0.14,0)$) {#4};
  \ifx\relax#5\relax\else
    \node[anchor=north west,font=\scSmall,text=scInkM]
      at ($#2+(0.56,-0.245)$) {#5};
  \fi}

% --- shared pgfplots look (only if the figure loaded pgfplots) ------------
\ifdefined\pgfplotsset
  \pgfplotsset{compat=1.18}
  \pgfplotsset{
    synaxis/.style={
      scale only axis,
      tick align=outside, tick pos=left,
      axis line style={draw=scInkM!55,line width=0.4pt},
      % One grid, on the value axis only.  No minor grid: at these sizes it
      % turns the plot frame into a grey wash.
      ymajorgrids=true, xmajorgrids=false,
      major grid style={draw=scRuleL,line width=0.3pt},
      major tick length=1.6pt,
      label style={font=\scBody,color=scInk},
      tick label style={font=\scSmall,color=scInkM},
      legend style={font=\scSmall,draw=none,fill=none,inner sep=1pt,
                    row sep=0.4pt,legend columns=1},
      legend cell align=left,
      every axis plot/.append style={line width=0.7pt,mark size=1.2pt},
    },
  }
\fi
}
% so that the figure compiles both from Figure/ and from the paper root.
%
% The visual vocabulary follows ../../../../Style/style.tex: the
% colour-blind-safe Okabe--Ito accents, quiet neutral rules and wash, circular
% panel badges, and a compact technical type scale are shared by the companion
% papers.  Semantic aliases below keep the CRN meaning explicit.
%
% DESIGN CONTRACT -----------------------------------------------------------
% 1. Geometry.  Every figure is laid out on a canvas exactly \scW wide (the
%    sn-jnl text width, 372 pt) and is included unscaled, so one point of type
%    is the same physical size in all five figures.  Open each picture with
%    \scCanvas{<height in cm>} to pin the bounding box.
% 2. Colour is semantic, never decorative.  Four roles only:
%       scAccent (green)       structure, the normal/expected path
%       scMark   (orange)      catalyst, certificate, or highlight
%       scAlert  (vermillion)  failure, violation, or uncovered set
%       scRef    (blue)        secondary or reference series
%    A panel gets at most one accent.  Titles are always ink; panel identity
%    is also encoded by a letter inside a circular badge, not by colour alone.
% 3. Type scale.  Three sizes, no others, and nothing below 6 pt:
%       \scTitle (7 pt bold)  panel titles
%       \scBody  (6.5 pt)     labels, values, in-figure prose
%       \scSmall (6 pt)       units, footnotes, dense graph labels
% ---------------------------------------------------------------------------
% Load guard.  The paper inputs every figure source inline, so this file is
% reached once per figure; without the guard the second \newlength below
% would abort the run.
\ifdefined\SynKitStyleLoaded\endinput\fi
\newcommand{\SynKitStyleLoaded}{}

% Load in the document preamble. XeLaTeX/LuaLaTeX and an installed Arial
% family are required; fontspec reports an error instead of substituting.
\RequirePackage[no-math]{fontspec}
\newfontfamily\scArial{Arial}[NFSSFamily=scArial,Ligatures=TeX]
% Keep manuscript mathematics unchanged; select Arial letters and numerals
% only inside figures. TeX supplies specialist mathematical symbols.
\RequirePackage[subdued,italic]{mathastext}
\MTDeclareVersion{synarial}{TU}{scArial}{m}{n}
\SetMathAlphabet{\Mathsf}{synarial}{TU}{scArial}{m}{n}
\SetMathAlphabet{\Mathtt}{synarial}{TU}{scArial}{m}{n}

\usetikzlibrary{positioning,calc,fit,backgrounds,arrows.meta,shapes.geometric,
                shapes.misc,decorations.pathreplacing,patterns,math}

% --- canvas ---------------------------------------------------------------
% 372 pt text width minus the 1 pt standalone border on each side.
\newlength{\scW}\setlength{\scW}{13.0326cm}
\providecommand{\scWnum}{13.0326}
% \scCanvas{<height>} pins the picture to the full text width, y running down
% from 0, so every figure shares one coordinate convention.  It *sets* the
% bounding box rather than contributing to it, so the emitted page is exactly
% the text width and every figure is included at 1:1 -- otherwise a stray
% invisible extent (pgfplots reserves a little around an axis) rescales one
% figure relative to the others and the same \scSmall prints at two sizes.
% Anything drawn outside the canvas is therefore cropped: keep content inside,
% and check with
%     \providecommand{\scCanvas}[1]{\path[draw=red] (0,0) rectangle (\scWnum,-#1);}
\providecommand{\scCanvas}[1]{%
  \path[use as bounding box] (0,0) rectangle (\scWnum,-#1);}

% --- palette --------------------------------------------------------------
% Canonical companion-paper values (../../../../Style/style.tex).
\definecolor{npgOrange}{HTML}{E69F00}
\definecolor{npgSky}{HTML}{56B4E9}
\definecolor{npgGreen}{HTML}{009E73}
\definecolor{npgYellow}{HTML}{F0E442}
\definecolor{npgBlue}{HTML}{0072B2}
\definecolor{npgVermillion}{HTML}{E64B35}
\definecolor{npgPurple}{HTML}{CC79A7}
\definecolor{figInk}{HTML}{1A1A1A}
\definecolor{figMuted}{HTML}{666666}
\definecolor{figHair}{HTML}{D9D9D9}
\definecolor{figWash}{HTML}{F7F7F7}

\colorlet{scInk}{figInk}              % primary text
\colorlet{scInkM}{figMuted}           % secondary text
\colorlet{scPaper}{figWash}           % quiet cards and note strips
\colorlet{scSpF}{npgSky!12}           % species fill
\colorlet{scSpE}{npgBlue}             % species edge
\colorlet{scRxF}{npgOrange!12}        % reaction fill
\colorlet{scRxE}{npgOrange!88!black}  % reaction edge
\colorlet{scReac}{npgBlue!78!black}   % reactant edge
\colorlet{scProd}{npgGreen!82!black}  % product edge
\colorlet{scHiF}{npgOrange!13}        % highlight fill
\colorlet{scHiE}{npgOrange!82!black}  % highlight edge
% Stage / accent tints.  stDE remains available for the fifth plot series.
\colorlet{stAF}{npgGreen!9}\colorlet{stAE}{npgGreen}
\colorlet{stBF}{npgOrange!11}\colorlet{stBE}{npgOrange}
\colorlet{stCF}{npgBlue!9}\colorlet{stCE}{npgBlue}
\colorlet{stDF}{npgPurple!10}\colorlet{stDE}{npgPurple!85!black}
\colorlet{stEF}{npgVermillion!9}\colorlet{stEE}{npgVermillion}

% Semantic aliases.  Figures should speak in these, not in stAE/stBE/...
\colorlet{scAccent}{stAE}\colorlet{scAccentF}{stAF}   % structure / normal
\colorlet{scMark}{stBE}\colorlet{scMarkF}{stBF}       % highlight / certificate
\colorlet{scRef}{stCE}\colorlet{scRefF}{stCF}         % secondary series
\colorlet{scAlert}{stEE}\colorlet{scAlertF}{stEF}     % failure / violation
\colorlet{scRule}{figHair}                            % card and axis hairlines
\colorlet{scRuleL}{figHair!55}                        % grid, dividers

% --- text styles ----------------------------------------------------------
% Three sizes, fixed in points so they do not drift with the document class.
% Arial is shared by labels, headings, panel badges, and chemistry labels.
\providecommand{\scTitle}{\scArial\fontsize{7}{8.4}\selectfont\bfseries\color{scInk}}
\providecommand{\scBody}{\scArial\fontsize{6.5}{7.8}\selectfont\color{scInkM}}
\providecommand{\scSmall}{\scArial\fontsize{6}{7.2}\selectfont\color{scInkM}}
\providecommand{\scSub}[1]{{\scSmall\itshape #1}}
\providecommand{\scPanel}[1]{{\scTitle(#1)}}
% Numerals that carry a result: ink, not colour, so the eye reads the value.
\providecommand{\scNum}{\scArial\fontsize{10}{11}\selectfont\bfseries\color{scInk}}

% --- shared TikZ styles ---------------------------------------------------
\tikzset{
  % text never hyphenates inside a figure
  every node/.append style={execute at begin node={%
      \hyphenpenalty=10000 \exhyphenpenalty=10000 \hbadness=10000 }},
  syn/.style={font=\scArial,color=scInk,
              execute at begin picture={\MTversion*{synarial}},
              >={Stealth[round,length=2.4mm,width=1.9mm]}},
  % containers ------------------------------------------------------------
  % One card chrome for the whole paper: white, 3 pt radius, one hairline.
  panel/.style={rounded corners=2.5pt,fill=none,draw=scRule,line width=0.45pt},
  card/.style={rounded corners=2.2pt,fill=scPaper,draw=scRule,line width=0.45pt,
               inner sep=5pt,align=left,anchor=north west},
  % a quiet strip for a panel's closing fact
  note/.style={rounded corners=2.5pt,fill=scPaper,draw=scRuleL,line width=0.4pt,
               inner xsep=6pt,inner ysep=3.5pt,align=left,font=\scSmall},
  % small labelled value pill
  chip/.style={rounded corners=2pt,fill=scPaper,draw=scRule,line width=0.4pt,
               inner xsep=4pt,inner ysep=2pt,font=\scSmall},
  % Circular letter badge: the same accessible panel marker used in Style/*.
  hbullet/.style={circle,fill=scInk,text=white,draw=none,
                  font=\scArial\fontsize{6}{6}\selectfont\bfseries,
                  minimum size=4.2mm,inner sep=0pt},
  ctag/.style={font=\scTitle,anchor=north west,text=scInk},
  % CRN primitives --------------------------------------------------------
  species/.style={rounded rectangle,fill=scSpF,draw=scSpE,line width=0.45pt,
                  inner xsep=2.4pt,inner ysep=1.6pt,font=\scSmall,
                  text=scInk,minimum height=3.6mm},
  cofactor/.style={species,fill=scRefF,draw=scRef},
  catalyst/.style={species,fill=scHiF,draw=scHiE},
  rxn/.style={rectangle,fill=scRxF,draw=scRxE,line width=0.45pt,
              inner sep=1.2pt,font=\scSmall,text=scInk,
              minimum width=3.0mm,minimum height=3.0mm},
  place/.style={circle,fill=scSpF,draw=scSpE,line width=0.45pt,
                minimum size=4.0mm,inner sep=0pt},
  trans/.style={rectangle,fill=scRxE,draw=scRxE,minimum width=1.1mm,
                minimum height=3.6mm,inner sep=0pt},
  token/.style={circle,fill=scInk,minimum size=1.1mm,inner sep=0pt},
  % edges -----------------------------------------------------------------
  reac/.style={->,line width=0.5pt,draw=scReac,shorten >=0.8pt,shorten <=0.8pt},
  prod/.style={->,line width=0.5pt,draw=scProd,shorten >=0.8pt,shorten <=0.8pt},
  cat/.style={->,line width=0.5pt,draw=scHiE,shorten >=0.8pt,shorten <=0.8pt},
  flow/.style={-{Stealth[round,length=2.2mm,width=1.8mm]},line width=1.1pt,
               draw=scInkM!45},
  flab/.style={font=\scSmall,text=scInkM,inner sep=2pt},
  stoi/.style={font=\scSmall,text=scInkM,inner sep=0.8pt},
  % verdict marks: shape carries the meaning, colour only reinforces it
  vpass/.style={circle,fill=scAccentF,draw=scAccent,text=scAccent,
                font=\fontsize{6}{6}\selectfont\bfseries,
                minimum size=3.2mm,inner sep=0pt},
  vfail/.style={circle,fill=scAlertF,draw=scAlert,text=scAlert,
                font=\fontsize{6}{6}\selectfont\bfseries,
                minimum size=3.2mm,inner sep=0pt},
  callout/.style={font=\scSmall,text=scInkM,align=left},
}

% \scHead[<badge colour>]{<anchor coord>}{<letter>}{<title>}{<subtitle or empty>}
% One header treatment for every panel in the paper.
\providecommand{\scHead}[5][scInk]{%
  \node[hbullet,fill=#1,anchor=north west] (scbadge#3) at ($#2+(0,-0.015)$)
    {\MakeUppercase{#3}};
  \node[ctag,anchor=west] at ($(scbadge#3.east)+(0.14,0)$) {#4};
  \ifx\relax#5\relax\else
    \node[anchor=north west,font=\scSmall,text=scInkM]
      at ($#2+(0.56,-0.245)$) {#5};
  \fi}

% --- shared pgfplots look (only if the figure loaded pgfplots) ------------
\ifdefined\pgfplotsset
  \pgfplotsset{compat=1.18}
  \pgfplotsset{
    synaxis/.style={
      scale only axis,
      tick align=outside, tick pos=left,
      axis line style={draw=scInkM!55,line width=0.4pt},
      % One grid, on the value axis only.  No minor grid: at these sizes it
      % turns the plot frame into a grey wash.
      ymajorgrids=true, xmajorgrids=false,
      major grid style={draw=scRuleL,line width=0.3pt},
      major tick length=1.6pt,
      label style={font=\scBody,color=scInk},
      tick label style={font=\scSmall,color=scInkM},
      legend style={font=\scSmall,draw=none,fill=none,inner sep=1pt,
                    row sep=0.4pt,legend columns=1},
      legend cell align=left,
      every axis plot/.append style={line width=0.7pt,mark size=1.2pt},
    },
  }
\fi
}
% so that the figure compiles both from Figure/ and from the paper root.
%
% The visual vocabulary follows ../../../../Style/style.tex: the
% colour-blind-safe Okabe--Ito accents, quiet neutral rules and wash, circular
% panel badges, and a compact technical type scale are shared by the companion
% papers.  Semantic aliases below keep the CRN meaning explicit.
%
% DESIGN CONTRACT -----------------------------------------------------------
% 1. Geometry.  Every figure is laid out on a canvas exactly \scW wide (the
%    sn-jnl text width, 372 pt) and is included unscaled, so one point of type
%    is the same physical size in all five figures.  Open each picture with
%    \scCanvas{<height in cm>} to pin the bounding box.
% 2. Colour is semantic, never decorative.  Four roles only:
%       scAccent (green)       structure, the normal/expected path
%       scMark   (orange)      catalyst, certificate, or highlight
%       scAlert  (vermillion)  failure, violation, or uncovered set
%       scRef    (blue)        secondary or reference series
%    A panel gets at most one accent.  Titles are always ink; panel identity
%    is also encoded by a letter inside a circular badge, not by colour alone.
% 3. Type scale.  Three sizes, no others, and nothing below 6 pt:
%       \scTitle (7 pt bold)  panel titles
%       \scBody  (6.5 pt)     labels, values, in-figure prose
%       \scSmall (6 pt)       units, footnotes, dense graph labels
% ---------------------------------------------------------------------------
% Load guard.  The paper inputs every figure source inline, so this file is
% reached once per figure; without the guard the second \newlength below
% would abort the run.
\ifdefined\SynKitStyleLoaded\endinput\fi
\newcommand{\SynKitStyleLoaded}{}

% Load in the document preamble. XeLaTeX/LuaLaTeX and an installed Arial
% family are required; fontspec reports an error instead of substituting.
\RequirePackage[no-math]{fontspec}
\newfontfamily\scArial{Arial}[NFSSFamily=scArial,Ligatures=TeX]
% Keep manuscript mathematics unchanged; select Arial letters and numerals
% only inside figures. TeX supplies specialist mathematical symbols.
\RequirePackage[subdued,italic]{mathastext}
\MTDeclareVersion{synarial}{TU}{scArial}{m}{n}
\SetMathAlphabet{\Mathsf}{synarial}{TU}{scArial}{m}{n}
\SetMathAlphabet{\Mathtt}{synarial}{TU}{scArial}{m}{n}

\usetikzlibrary{positioning,calc,fit,backgrounds,arrows.meta,shapes.geometric,
                shapes.misc,decorations.pathreplacing,patterns,math}

% --- canvas ---------------------------------------------------------------
% 372 pt text width minus the 1 pt standalone border on each side.
\newlength{\scW}\setlength{\scW}{13.0326cm}
\providecommand{\scWnum}{13.0326}
% \scCanvas{<height>} pins the picture to the full text width, y running down
% from 0, so every figure shares one coordinate convention.  It *sets* the
% bounding box rather than contributing to it, so the emitted page is exactly
% the text width and every figure is included at 1:1 -- otherwise a stray
% invisible extent (pgfplots reserves a little around an axis) rescales one
% figure relative to the others and the same \scSmall prints at two sizes.
% Anything drawn outside the canvas is therefore cropped: keep content inside,
% and check with
%     \providecommand{\scCanvas}[1]{\path[draw=red] (0,0) rectangle (\scWnum,-#1);}
\providecommand{\scCanvas}[1]{%
  \path[use as bounding box] (0,0) rectangle (\scWnum,-#1);}

% --- palette --------------------------------------------------------------
% Canonical companion-paper values (../../../../Style/style.tex).
\definecolor{npgOrange}{HTML}{E69F00}
\definecolor{npgSky}{HTML}{56B4E9}
\definecolor{npgGreen}{HTML}{009E73}
\definecolor{npgYellow}{HTML}{F0E442}
\definecolor{npgBlue}{HTML}{0072B2}
\definecolor{npgVermillion}{HTML}{E64B35}
\definecolor{npgPurple}{HTML}{CC79A7}
\definecolor{figInk}{HTML}{1A1A1A}
\definecolor{figMuted}{HTML}{666666}
\definecolor{figHair}{HTML}{D9D9D9}
\definecolor{figWash}{HTML}{F7F7F7}

\colorlet{scInk}{figInk}              % primary text
\colorlet{scInkM}{figMuted}           % secondary text
\colorlet{scPaper}{figWash}           % quiet cards and note strips
\colorlet{scSpF}{npgSky!12}           % species fill
\colorlet{scSpE}{npgBlue}             % species edge
\colorlet{scRxF}{npgOrange!12}        % reaction fill
\colorlet{scRxE}{npgOrange!88!black}  % reaction edge
\colorlet{scReac}{npgBlue!78!black}   % reactant edge
\colorlet{scProd}{npgGreen!82!black}  % product edge
\colorlet{scHiF}{npgOrange!13}        % highlight fill
\colorlet{scHiE}{npgOrange!82!black}  % highlight edge
% Stage / accent tints.  stDE remains available for the fifth plot series.
\colorlet{stAF}{npgGreen!9}\colorlet{stAE}{npgGreen}
\colorlet{stBF}{npgOrange!11}\colorlet{stBE}{npgOrange}
\colorlet{stCF}{npgBlue!9}\colorlet{stCE}{npgBlue}
\colorlet{stDF}{npgPurple!10}\colorlet{stDE}{npgPurple!85!black}
\colorlet{stEF}{npgVermillion!9}\colorlet{stEE}{npgVermillion}

% Semantic aliases.  Figures should speak in these, not in stAE/stBE/...
\colorlet{scAccent}{stAE}\colorlet{scAccentF}{stAF}   % structure / normal
\colorlet{scMark}{stBE}\colorlet{scMarkF}{stBF}       % highlight / certificate
\colorlet{scRef}{stCE}\colorlet{scRefF}{stCF}         % secondary series
\colorlet{scAlert}{stEE}\colorlet{scAlertF}{stEF}     % failure / violation
\colorlet{scRule}{figHair}                            % card and axis hairlines
\colorlet{scRuleL}{figHair!55}                        % grid, dividers

% --- text styles ----------------------------------------------------------
% Three sizes, fixed in points so they do not drift with the document class.
% Arial is shared by labels, headings, panel badges, and chemistry labels.
\providecommand{\scTitle}{\scArial\fontsize{7}{8.4}\selectfont\bfseries\color{scInk}}
\providecommand{\scBody}{\scArial\fontsize{6.5}{7.8}\selectfont\color{scInkM}}
\providecommand{\scSmall}{\scArial\fontsize{6}{7.2}\selectfont\color{scInkM}}
\providecommand{\scSub}[1]{{\scSmall\itshape #1}}
\providecommand{\scPanel}[1]{{\scTitle(#1)}}
% Numerals that carry a result: ink, not colour, so the eye reads the value.
\providecommand{\scNum}{\scArial\fontsize{10}{11}\selectfont\bfseries\color{scInk}}

% --- shared TikZ styles ---------------------------------------------------
\tikzset{
  % text never hyphenates inside a figure
  every node/.append style={execute at begin node={%
      \hyphenpenalty=10000 \exhyphenpenalty=10000 \hbadness=10000 }},
  syn/.style={font=\scArial,color=scInk,
              execute at begin picture={\MTversion*{synarial}},
              >={Stealth[round,length=2.4mm,width=1.9mm]}},
  % containers ------------------------------------------------------------
  % One card chrome for the whole paper: white, 3 pt radius, one hairline.
  panel/.style={rounded corners=2.5pt,fill=none,draw=scRule,line width=0.45pt},
  card/.style={rounded corners=2.2pt,fill=scPaper,draw=scRule,line width=0.45pt,
               inner sep=5pt,align=left,anchor=north west},
  % a quiet strip for a panel's closing fact
  note/.style={rounded corners=2.5pt,fill=scPaper,draw=scRuleL,line width=0.4pt,
               inner xsep=6pt,inner ysep=3.5pt,align=left,font=\scSmall},
  % small labelled value pill
  chip/.style={rounded corners=2pt,fill=scPaper,draw=scRule,line width=0.4pt,
               inner xsep=4pt,inner ysep=2pt,font=\scSmall},
  % Circular letter badge: the same accessible panel marker used in Style/*.
  hbullet/.style={circle,fill=scInk,text=white,draw=none,
                  font=\scArial\fontsize{6}{6}\selectfont\bfseries,
                  minimum size=4.2mm,inner sep=0pt},
  ctag/.style={font=\scTitle,anchor=north west,text=scInk},
  % CRN primitives --------------------------------------------------------
  species/.style={rounded rectangle,fill=scSpF,draw=scSpE,line width=0.45pt,
                  inner xsep=2.4pt,inner ysep=1.6pt,font=\scSmall,
                  text=scInk,minimum height=3.6mm},
  cofactor/.style={species,fill=scRefF,draw=scRef},
  catalyst/.style={species,fill=scHiF,draw=scHiE},
  rxn/.style={rectangle,fill=scRxF,draw=scRxE,line width=0.45pt,
              inner sep=1.2pt,font=\scSmall,text=scInk,
              minimum width=3.0mm,minimum height=3.0mm},
  place/.style={circle,fill=scSpF,draw=scSpE,line width=0.45pt,
                minimum size=4.0mm,inner sep=0pt},
  trans/.style={rectangle,fill=scRxE,draw=scRxE,minimum width=1.1mm,
                minimum height=3.6mm,inner sep=0pt},
  token/.style={circle,fill=scInk,minimum size=1.1mm,inner sep=0pt},
  % edges -----------------------------------------------------------------
  reac/.style={->,line width=0.5pt,draw=scReac,shorten >=0.8pt,shorten <=0.8pt},
  prod/.style={->,line width=0.5pt,draw=scProd,shorten >=0.8pt,shorten <=0.8pt},
  cat/.style={->,line width=0.5pt,draw=scHiE,shorten >=0.8pt,shorten <=0.8pt},
  flow/.style={-{Stealth[round,length=2.2mm,width=1.8mm]},line width=1.1pt,
               draw=scInkM!45},
  flab/.style={font=\scSmall,text=scInkM,inner sep=2pt},
  stoi/.style={font=\scSmall,text=scInkM,inner sep=0.8pt},
  % verdict marks: shape carries the meaning, colour only reinforces it
  vpass/.style={circle,fill=scAccentF,draw=scAccent,text=scAccent,
                font=\fontsize{6}{6}\selectfont\bfseries,
                minimum size=3.2mm,inner sep=0pt},
  vfail/.style={circle,fill=scAlertF,draw=scAlert,text=scAlert,
                font=\fontsize{6}{6}\selectfont\bfseries,
                minimum size=3.2mm,inner sep=0pt},
  callout/.style={font=\scSmall,text=scInkM,align=left},
}

% \scHead[<badge colour>]{<anchor coord>}{<letter>}{<title>}{<subtitle or empty>}
% One header treatment for every panel in the paper.
\providecommand{\scHead}[5][scInk]{%
  \node[hbullet,fill=#1,anchor=north west] (scbadge#3) at ($#2+(0,-0.015)$)
    {\MakeUppercase{#3}};
  \node[ctag,anchor=west] at ($(scbadge#3.east)+(0.14,0)$) {#4};
  \ifx\relax#5\relax\else
    \node[anchor=north west,font=\scSmall,text=scInkM]
      at ($#2+(0.56,-0.245)$) {#5};
  \fi}

% --- shared pgfplots look (only if the figure loaded pgfplots) ------------
\ifdefined\pgfplotsset
  \pgfplotsset{compat=1.18}
  \pgfplotsset{
    synaxis/.style={
      scale only axis,
      tick align=outside, tick pos=left,
      axis line style={draw=scInkM!55,line width=0.4pt},
      % One grid, on the value axis only.  No minor grid: at these sizes it
      % turns the plot frame into a grey wash.
      ymajorgrids=true, xmajorgrids=false,
      major grid style={draw=scRuleL,line width=0.3pt},
      major tick length=1.6pt,
      label style={font=\scBody,color=scInk},
      tick label style={font=\scSmall,color=scInkM},
      legend style={font=\scSmall,draw=none,fill=none,inner sep=1pt,
                    row sep=0.4pt,legend columns=1},
      legend cell align=left,
      every axis plot/.append style={line width=0.7pt,mark size=1.2pt},
    },
  }
\fi
}{%
  \IfFileExists{Figure/synstyle.tex}{% ---------------------------------------------------------------------------
% synstyle.tex --- shared look for every TikZ figure of the SynKit CRN paper.
%
% Input it from a figure file:
%     \IfFileExists{synstyle.tex}{% ---------------------------------------------------------------------------
% synstyle.tex --- shared look for every TikZ figure of the SynKit CRN paper.
%
% Input it from a figure file:
%     \IfFileExists{synstyle.tex}{% ---------------------------------------------------------------------------
% synstyle.tex --- shared look for every TikZ figure of the SynKit CRN paper.
%
% Input it from a figure file:
%     \IfFileExists{synstyle.tex}{\input{synstyle}}{\input{Figure/synstyle}}
% so that the figure compiles both from Figure/ and from the paper root.
%
% The visual vocabulary follows ../../../../Style/style.tex: the
% colour-blind-safe Okabe--Ito accents, quiet neutral rules and wash, circular
% panel badges, and a compact technical type scale are shared by the companion
% papers.  Semantic aliases below keep the CRN meaning explicit.
%
% DESIGN CONTRACT -----------------------------------------------------------
% 1. Geometry.  Every figure is laid out on a canvas exactly \scW wide (the
%    sn-jnl text width, 372 pt) and is included unscaled, so one point of type
%    is the same physical size in all five figures.  Open each picture with
%    \scCanvas{<height in cm>} to pin the bounding box.
% 2. Colour is semantic, never decorative.  Four roles only:
%       scAccent (green)       structure, the normal/expected path
%       scMark   (orange)      catalyst, certificate, or highlight
%       scAlert  (vermillion)  failure, violation, or uncovered set
%       scRef    (blue)        secondary or reference series
%    A panel gets at most one accent.  Titles are always ink; panel identity
%    is also encoded by a letter inside a circular badge, not by colour alone.
% 3. Type scale.  Three sizes, no others, and nothing below 6 pt:
%       \scTitle (7 pt bold)  panel titles
%       \scBody  (6.5 pt)     labels, values, in-figure prose
%       \scSmall (6 pt)       units, footnotes, dense graph labels
% ---------------------------------------------------------------------------
% Load guard.  The paper inputs every figure source inline, so this file is
% reached once per figure; without the guard the second \newlength below
% would abort the run.
\ifdefined\SynKitStyleLoaded\endinput\fi
\newcommand{\SynKitStyleLoaded}{}

% Load in the document preamble. XeLaTeX/LuaLaTeX and an installed Arial
% family are required; fontspec reports an error instead of substituting.
\RequirePackage[no-math]{fontspec}
\newfontfamily\scArial{Arial}[NFSSFamily=scArial,Ligatures=TeX]
% Keep manuscript mathematics unchanged; select Arial letters and numerals
% only inside figures. TeX supplies specialist mathematical symbols.
\RequirePackage[subdued,italic]{mathastext}
\MTDeclareVersion{synarial}{TU}{scArial}{m}{n}
\SetMathAlphabet{\Mathsf}{synarial}{TU}{scArial}{m}{n}
\SetMathAlphabet{\Mathtt}{synarial}{TU}{scArial}{m}{n}

\usetikzlibrary{positioning,calc,fit,backgrounds,arrows.meta,shapes.geometric,
                shapes.misc,decorations.pathreplacing,patterns,math}

% --- canvas ---------------------------------------------------------------
% 372 pt text width minus the 1 pt standalone border on each side.
\newlength{\scW}\setlength{\scW}{13.0326cm}
\providecommand{\scWnum}{13.0326}
% \scCanvas{<height>} pins the picture to the full text width, y running down
% from 0, so every figure shares one coordinate convention.  It *sets* the
% bounding box rather than contributing to it, so the emitted page is exactly
% the text width and every figure is included at 1:1 -- otherwise a stray
% invisible extent (pgfplots reserves a little around an axis) rescales one
% figure relative to the others and the same \scSmall prints at two sizes.
% Anything drawn outside the canvas is therefore cropped: keep content inside,
% and check with
%     \providecommand{\scCanvas}[1]{\path[draw=red] (0,0) rectangle (\scWnum,-#1);}
\providecommand{\scCanvas}[1]{%
  \path[use as bounding box] (0,0) rectangle (\scWnum,-#1);}

% --- palette --------------------------------------------------------------
% Canonical companion-paper values (../../../../Style/style.tex).
\definecolor{npgOrange}{HTML}{E69F00}
\definecolor{npgSky}{HTML}{56B4E9}
\definecolor{npgGreen}{HTML}{009E73}
\definecolor{npgYellow}{HTML}{F0E442}
\definecolor{npgBlue}{HTML}{0072B2}
\definecolor{npgVermillion}{HTML}{E64B35}
\definecolor{npgPurple}{HTML}{CC79A7}
\definecolor{figInk}{HTML}{1A1A1A}
\definecolor{figMuted}{HTML}{666666}
\definecolor{figHair}{HTML}{D9D9D9}
\definecolor{figWash}{HTML}{F7F7F7}

\colorlet{scInk}{figInk}              % primary text
\colorlet{scInkM}{figMuted}           % secondary text
\colorlet{scPaper}{figWash}           % quiet cards and note strips
\colorlet{scSpF}{npgSky!12}           % species fill
\colorlet{scSpE}{npgBlue}             % species edge
\colorlet{scRxF}{npgOrange!12}        % reaction fill
\colorlet{scRxE}{npgOrange!88!black}  % reaction edge
\colorlet{scReac}{npgBlue!78!black}   % reactant edge
\colorlet{scProd}{npgGreen!82!black}  % product edge
\colorlet{scHiF}{npgOrange!13}        % highlight fill
\colorlet{scHiE}{npgOrange!82!black}  % highlight edge
% Stage / accent tints.  stDE remains available for the fifth plot series.
\colorlet{stAF}{npgGreen!9}\colorlet{stAE}{npgGreen}
\colorlet{stBF}{npgOrange!11}\colorlet{stBE}{npgOrange}
\colorlet{stCF}{npgBlue!9}\colorlet{stCE}{npgBlue}
\colorlet{stDF}{npgPurple!10}\colorlet{stDE}{npgPurple!85!black}
\colorlet{stEF}{npgVermillion!9}\colorlet{stEE}{npgVermillion}

% Semantic aliases.  Figures should speak in these, not in stAE/stBE/...
\colorlet{scAccent}{stAE}\colorlet{scAccentF}{stAF}   % structure / normal
\colorlet{scMark}{stBE}\colorlet{scMarkF}{stBF}       % highlight / certificate
\colorlet{scRef}{stCE}\colorlet{scRefF}{stCF}         % secondary series
\colorlet{scAlert}{stEE}\colorlet{scAlertF}{stEF}     % failure / violation
\colorlet{scRule}{figHair}                            % card and axis hairlines
\colorlet{scRuleL}{figHair!55}                        % grid, dividers

% --- text styles ----------------------------------------------------------
% Three sizes, fixed in points so they do not drift with the document class.
% Arial is shared by labels, headings, panel badges, and chemistry labels.
\providecommand{\scTitle}{\scArial\fontsize{7}{8.4}\selectfont\bfseries\color{scInk}}
\providecommand{\scBody}{\scArial\fontsize{6.5}{7.8}\selectfont\color{scInkM}}
\providecommand{\scSmall}{\scArial\fontsize{6}{7.2}\selectfont\color{scInkM}}
\providecommand{\scSub}[1]{{\scSmall\itshape #1}}
\providecommand{\scPanel}[1]{{\scTitle(#1)}}
% Numerals that carry a result: ink, not colour, so the eye reads the value.
\providecommand{\scNum}{\scArial\fontsize{10}{11}\selectfont\bfseries\color{scInk}}

% --- shared TikZ styles ---------------------------------------------------
\tikzset{
  % text never hyphenates inside a figure
  every node/.append style={execute at begin node={%
      \hyphenpenalty=10000 \exhyphenpenalty=10000 \hbadness=10000 }},
  syn/.style={font=\scArial,color=scInk,
              execute at begin picture={\MTversion*{synarial}},
              >={Stealth[round,length=2.4mm,width=1.9mm]}},
  % containers ------------------------------------------------------------
  % One card chrome for the whole paper: white, 3 pt radius, one hairline.
  panel/.style={rounded corners=2.5pt,fill=none,draw=scRule,line width=0.45pt},
  card/.style={rounded corners=2.2pt,fill=scPaper,draw=scRule,line width=0.45pt,
               inner sep=5pt,align=left,anchor=north west},
  % a quiet strip for a panel's closing fact
  note/.style={rounded corners=2.5pt,fill=scPaper,draw=scRuleL,line width=0.4pt,
               inner xsep=6pt,inner ysep=3.5pt,align=left,font=\scSmall},
  % small labelled value pill
  chip/.style={rounded corners=2pt,fill=scPaper,draw=scRule,line width=0.4pt,
               inner xsep=4pt,inner ysep=2pt,font=\scSmall},
  % Circular letter badge: the same accessible panel marker used in Style/*.
  hbullet/.style={circle,fill=scInk,text=white,draw=none,
                  font=\scArial\fontsize{6}{6}\selectfont\bfseries,
                  minimum size=4.2mm,inner sep=0pt},
  ctag/.style={font=\scTitle,anchor=north west,text=scInk},
  % CRN primitives --------------------------------------------------------
  species/.style={rounded rectangle,fill=scSpF,draw=scSpE,line width=0.45pt,
                  inner xsep=2.4pt,inner ysep=1.6pt,font=\scSmall,
                  text=scInk,minimum height=3.6mm},
  cofactor/.style={species,fill=scRefF,draw=scRef},
  catalyst/.style={species,fill=scHiF,draw=scHiE},
  rxn/.style={rectangle,fill=scRxF,draw=scRxE,line width=0.45pt,
              inner sep=1.2pt,font=\scSmall,text=scInk,
              minimum width=3.0mm,minimum height=3.0mm},
  place/.style={circle,fill=scSpF,draw=scSpE,line width=0.45pt,
                minimum size=4.0mm,inner sep=0pt},
  trans/.style={rectangle,fill=scRxE,draw=scRxE,minimum width=1.1mm,
                minimum height=3.6mm,inner sep=0pt},
  token/.style={circle,fill=scInk,minimum size=1.1mm,inner sep=0pt},
  % edges -----------------------------------------------------------------
  reac/.style={->,line width=0.5pt,draw=scReac,shorten >=0.8pt,shorten <=0.8pt},
  prod/.style={->,line width=0.5pt,draw=scProd,shorten >=0.8pt,shorten <=0.8pt},
  cat/.style={->,line width=0.5pt,draw=scHiE,shorten >=0.8pt,shorten <=0.8pt},
  flow/.style={-{Stealth[round,length=2.2mm,width=1.8mm]},line width=1.1pt,
               draw=scInkM!45},
  flab/.style={font=\scSmall,text=scInkM,inner sep=2pt},
  stoi/.style={font=\scSmall,text=scInkM,inner sep=0.8pt},
  % verdict marks: shape carries the meaning, colour only reinforces it
  vpass/.style={circle,fill=scAccentF,draw=scAccent,text=scAccent,
                font=\fontsize{6}{6}\selectfont\bfseries,
                minimum size=3.2mm,inner sep=0pt},
  vfail/.style={circle,fill=scAlertF,draw=scAlert,text=scAlert,
                font=\fontsize{6}{6}\selectfont\bfseries,
                minimum size=3.2mm,inner sep=0pt},
  callout/.style={font=\scSmall,text=scInkM,align=left},
}

% \scHead[<badge colour>]{<anchor coord>}{<letter>}{<title>}{<subtitle or empty>}
% One header treatment for every panel in the paper.
\providecommand{\scHead}[5][scInk]{%
  \node[hbullet,fill=#1,anchor=north west] (scbadge#3) at ($#2+(0,-0.015)$)
    {\MakeUppercase{#3}};
  \node[ctag,anchor=west] at ($(scbadge#3.east)+(0.14,0)$) {#4};
  \ifx\relax#5\relax\else
    \node[anchor=north west,font=\scSmall,text=scInkM]
      at ($#2+(0.56,-0.245)$) {#5};
  \fi}

% --- shared pgfplots look (only if the figure loaded pgfplots) ------------
\ifdefined\pgfplotsset
  \pgfplotsset{compat=1.18}
  \pgfplotsset{
    synaxis/.style={
      scale only axis,
      tick align=outside, tick pos=left,
      axis line style={draw=scInkM!55,line width=0.4pt},
      % One grid, on the value axis only.  No minor grid: at these sizes it
      % turns the plot frame into a grey wash.
      ymajorgrids=true, xmajorgrids=false,
      major grid style={draw=scRuleL,line width=0.3pt},
      major tick length=1.6pt,
      label style={font=\scBody,color=scInk},
      tick label style={font=\scSmall,color=scInkM},
      legend style={font=\scSmall,draw=none,fill=none,inner sep=1pt,
                    row sep=0.4pt,legend columns=1},
      legend cell align=left,
      every axis plot/.append style={line width=0.7pt,mark size=1.2pt},
    },
  }
\fi
}{% ---------------------------------------------------------------------------
% synstyle.tex --- shared look for every TikZ figure of the SynKit CRN paper.
%
% Input it from a figure file:
%     \IfFileExists{synstyle.tex}{\input{synstyle}}{\input{Figure/synstyle}}
% so that the figure compiles both from Figure/ and from the paper root.
%
% The visual vocabulary follows ../../../../Style/style.tex: the
% colour-blind-safe Okabe--Ito accents, quiet neutral rules and wash, circular
% panel badges, and a compact technical type scale are shared by the companion
% papers.  Semantic aliases below keep the CRN meaning explicit.
%
% DESIGN CONTRACT -----------------------------------------------------------
% 1. Geometry.  Every figure is laid out on a canvas exactly \scW wide (the
%    sn-jnl text width, 372 pt) and is included unscaled, so one point of type
%    is the same physical size in all five figures.  Open each picture with
%    \scCanvas{<height in cm>} to pin the bounding box.
% 2. Colour is semantic, never decorative.  Four roles only:
%       scAccent (green)       structure, the normal/expected path
%       scMark   (orange)      catalyst, certificate, or highlight
%       scAlert  (vermillion)  failure, violation, or uncovered set
%       scRef    (blue)        secondary or reference series
%    A panel gets at most one accent.  Titles are always ink; panel identity
%    is also encoded by a letter inside a circular badge, not by colour alone.
% 3. Type scale.  Three sizes, no others, and nothing below 6 pt:
%       \scTitle (7 pt bold)  panel titles
%       \scBody  (6.5 pt)     labels, values, in-figure prose
%       \scSmall (6 pt)       units, footnotes, dense graph labels
% ---------------------------------------------------------------------------
% Load guard.  The paper inputs every figure source inline, so this file is
% reached once per figure; without the guard the second \newlength below
% would abort the run.
\ifdefined\SynKitStyleLoaded\endinput\fi
\newcommand{\SynKitStyleLoaded}{}

% Load in the document preamble. XeLaTeX/LuaLaTeX and an installed Arial
% family are required; fontspec reports an error instead of substituting.
\RequirePackage[no-math]{fontspec}
\newfontfamily\scArial{Arial}[NFSSFamily=scArial,Ligatures=TeX]
% Keep manuscript mathematics unchanged; select Arial letters and numerals
% only inside figures. TeX supplies specialist mathematical symbols.
\RequirePackage[subdued,italic]{mathastext}
\MTDeclareVersion{synarial}{TU}{scArial}{m}{n}
\SetMathAlphabet{\Mathsf}{synarial}{TU}{scArial}{m}{n}
\SetMathAlphabet{\Mathtt}{synarial}{TU}{scArial}{m}{n}

\usetikzlibrary{positioning,calc,fit,backgrounds,arrows.meta,shapes.geometric,
                shapes.misc,decorations.pathreplacing,patterns,math}

% --- canvas ---------------------------------------------------------------
% 372 pt text width minus the 1 pt standalone border on each side.
\newlength{\scW}\setlength{\scW}{13.0326cm}
\providecommand{\scWnum}{13.0326}
% \scCanvas{<height>} pins the picture to the full text width, y running down
% from 0, so every figure shares one coordinate convention.  It *sets* the
% bounding box rather than contributing to it, so the emitted page is exactly
% the text width and every figure is included at 1:1 -- otherwise a stray
% invisible extent (pgfplots reserves a little around an axis) rescales one
% figure relative to the others and the same \scSmall prints at two sizes.
% Anything drawn outside the canvas is therefore cropped: keep content inside,
% and check with
%     \providecommand{\scCanvas}[1]{\path[draw=red] (0,0) rectangle (\scWnum,-#1);}
\providecommand{\scCanvas}[1]{%
  \path[use as bounding box] (0,0) rectangle (\scWnum,-#1);}

% --- palette --------------------------------------------------------------
% Canonical companion-paper values (../../../../Style/style.tex).
\definecolor{npgOrange}{HTML}{E69F00}
\definecolor{npgSky}{HTML}{56B4E9}
\definecolor{npgGreen}{HTML}{009E73}
\definecolor{npgYellow}{HTML}{F0E442}
\definecolor{npgBlue}{HTML}{0072B2}
\definecolor{npgVermillion}{HTML}{E64B35}
\definecolor{npgPurple}{HTML}{CC79A7}
\definecolor{figInk}{HTML}{1A1A1A}
\definecolor{figMuted}{HTML}{666666}
\definecolor{figHair}{HTML}{D9D9D9}
\definecolor{figWash}{HTML}{F7F7F7}

\colorlet{scInk}{figInk}              % primary text
\colorlet{scInkM}{figMuted}           % secondary text
\colorlet{scPaper}{figWash}           % quiet cards and note strips
\colorlet{scSpF}{npgSky!12}           % species fill
\colorlet{scSpE}{npgBlue}             % species edge
\colorlet{scRxF}{npgOrange!12}        % reaction fill
\colorlet{scRxE}{npgOrange!88!black}  % reaction edge
\colorlet{scReac}{npgBlue!78!black}   % reactant edge
\colorlet{scProd}{npgGreen!82!black}  % product edge
\colorlet{scHiF}{npgOrange!13}        % highlight fill
\colorlet{scHiE}{npgOrange!82!black}  % highlight edge
% Stage / accent tints.  stDE remains available for the fifth plot series.
\colorlet{stAF}{npgGreen!9}\colorlet{stAE}{npgGreen}
\colorlet{stBF}{npgOrange!11}\colorlet{stBE}{npgOrange}
\colorlet{stCF}{npgBlue!9}\colorlet{stCE}{npgBlue}
\colorlet{stDF}{npgPurple!10}\colorlet{stDE}{npgPurple!85!black}
\colorlet{stEF}{npgVermillion!9}\colorlet{stEE}{npgVermillion}

% Semantic aliases.  Figures should speak in these, not in stAE/stBE/...
\colorlet{scAccent}{stAE}\colorlet{scAccentF}{stAF}   % structure / normal
\colorlet{scMark}{stBE}\colorlet{scMarkF}{stBF}       % highlight / certificate
\colorlet{scRef}{stCE}\colorlet{scRefF}{stCF}         % secondary series
\colorlet{scAlert}{stEE}\colorlet{scAlertF}{stEF}     % failure / violation
\colorlet{scRule}{figHair}                            % card and axis hairlines
\colorlet{scRuleL}{figHair!55}                        % grid, dividers

% --- text styles ----------------------------------------------------------
% Three sizes, fixed in points so they do not drift with the document class.
% Arial is shared by labels, headings, panel badges, and chemistry labels.
\providecommand{\scTitle}{\scArial\fontsize{7}{8.4}\selectfont\bfseries\color{scInk}}
\providecommand{\scBody}{\scArial\fontsize{6.5}{7.8}\selectfont\color{scInkM}}
\providecommand{\scSmall}{\scArial\fontsize{6}{7.2}\selectfont\color{scInkM}}
\providecommand{\scSub}[1]{{\scSmall\itshape #1}}
\providecommand{\scPanel}[1]{{\scTitle(#1)}}
% Numerals that carry a result: ink, not colour, so the eye reads the value.
\providecommand{\scNum}{\scArial\fontsize{10}{11}\selectfont\bfseries\color{scInk}}

% --- shared TikZ styles ---------------------------------------------------
\tikzset{
  % text never hyphenates inside a figure
  every node/.append style={execute at begin node={%
      \hyphenpenalty=10000 \exhyphenpenalty=10000 \hbadness=10000 }},
  syn/.style={font=\scArial,color=scInk,
              execute at begin picture={\MTversion*{synarial}},
              >={Stealth[round,length=2.4mm,width=1.9mm]}},
  % containers ------------------------------------------------------------
  % One card chrome for the whole paper: white, 3 pt radius, one hairline.
  panel/.style={rounded corners=2.5pt,fill=none,draw=scRule,line width=0.45pt},
  card/.style={rounded corners=2.2pt,fill=scPaper,draw=scRule,line width=0.45pt,
               inner sep=5pt,align=left,anchor=north west},
  % a quiet strip for a panel's closing fact
  note/.style={rounded corners=2.5pt,fill=scPaper,draw=scRuleL,line width=0.4pt,
               inner xsep=6pt,inner ysep=3.5pt,align=left,font=\scSmall},
  % small labelled value pill
  chip/.style={rounded corners=2pt,fill=scPaper,draw=scRule,line width=0.4pt,
               inner xsep=4pt,inner ysep=2pt,font=\scSmall},
  % Circular letter badge: the same accessible panel marker used in Style/*.
  hbullet/.style={circle,fill=scInk,text=white,draw=none,
                  font=\scArial\fontsize{6}{6}\selectfont\bfseries,
                  minimum size=4.2mm,inner sep=0pt},
  ctag/.style={font=\scTitle,anchor=north west,text=scInk},
  % CRN primitives --------------------------------------------------------
  species/.style={rounded rectangle,fill=scSpF,draw=scSpE,line width=0.45pt,
                  inner xsep=2.4pt,inner ysep=1.6pt,font=\scSmall,
                  text=scInk,minimum height=3.6mm},
  cofactor/.style={species,fill=scRefF,draw=scRef},
  catalyst/.style={species,fill=scHiF,draw=scHiE},
  rxn/.style={rectangle,fill=scRxF,draw=scRxE,line width=0.45pt,
              inner sep=1.2pt,font=\scSmall,text=scInk,
              minimum width=3.0mm,minimum height=3.0mm},
  place/.style={circle,fill=scSpF,draw=scSpE,line width=0.45pt,
                minimum size=4.0mm,inner sep=0pt},
  trans/.style={rectangle,fill=scRxE,draw=scRxE,minimum width=1.1mm,
                minimum height=3.6mm,inner sep=0pt},
  token/.style={circle,fill=scInk,minimum size=1.1mm,inner sep=0pt},
  % edges -----------------------------------------------------------------
  reac/.style={->,line width=0.5pt,draw=scReac,shorten >=0.8pt,shorten <=0.8pt},
  prod/.style={->,line width=0.5pt,draw=scProd,shorten >=0.8pt,shorten <=0.8pt},
  cat/.style={->,line width=0.5pt,draw=scHiE,shorten >=0.8pt,shorten <=0.8pt},
  flow/.style={-{Stealth[round,length=2.2mm,width=1.8mm]},line width=1.1pt,
               draw=scInkM!45},
  flab/.style={font=\scSmall,text=scInkM,inner sep=2pt},
  stoi/.style={font=\scSmall,text=scInkM,inner sep=0.8pt},
  % verdict marks: shape carries the meaning, colour only reinforces it
  vpass/.style={circle,fill=scAccentF,draw=scAccent,text=scAccent,
                font=\fontsize{6}{6}\selectfont\bfseries,
                minimum size=3.2mm,inner sep=0pt},
  vfail/.style={circle,fill=scAlertF,draw=scAlert,text=scAlert,
                font=\fontsize{6}{6}\selectfont\bfseries,
                minimum size=3.2mm,inner sep=0pt},
  callout/.style={font=\scSmall,text=scInkM,align=left},
}

% \scHead[<badge colour>]{<anchor coord>}{<letter>}{<title>}{<subtitle or empty>}
% One header treatment for every panel in the paper.
\providecommand{\scHead}[5][scInk]{%
  \node[hbullet,fill=#1,anchor=north west] (scbadge#3) at ($#2+(0,-0.015)$)
    {\MakeUppercase{#3}};
  \node[ctag,anchor=west] at ($(scbadge#3.east)+(0.14,0)$) {#4};
  \ifx\relax#5\relax\else
    \node[anchor=north west,font=\scSmall,text=scInkM]
      at ($#2+(0.56,-0.245)$) {#5};
  \fi}

% --- shared pgfplots look (only if the figure loaded pgfplots) ------------
\ifdefined\pgfplotsset
  \pgfplotsset{compat=1.18}
  \pgfplotsset{
    synaxis/.style={
      scale only axis,
      tick align=outside, tick pos=left,
      axis line style={draw=scInkM!55,line width=0.4pt},
      % One grid, on the value axis only.  No minor grid: at these sizes it
      % turns the plot frame into a grey wash.
      ymajorgrids=true, xmajorgrids=false,
      major grid style={draw=scRuleL,line width=0.3pt},
      major tick length=1.6pt,
      label style={font=\scBody,color=scInk},
      tick label style={font=\scSmall,color=scInkM},
      legend style={font=\scSmall,draw=none,fill=none,inner sep=1pt,
                    row sep=0.4pt,legend columns=1},
      legend cell align=left,
      every axis plot/.append style={line width=0.7pt,mark size=1.2pt},
    },
  }
\fi
}
% so that the figure compiles both from Figure/ and from the paper root.
%
% The visual vocabulary follows ../../../../Style/style.tex: the
% colour-blind-safe Okabe--Ito accents, quiet neutral rules and wash, circular
% panel badges, and a compact technical type scale are shared by the companion
% papers.  Semantic aliases below keep the CRN meaning explicit.
%
% DESIGN CONTRACT -----------------------------------------------------------
% 1. Geometry.  Every figure is laid out on a canvas exactly \scW wide (the
%    sn-jnl text width, 372 pt) and is included unscaled, so one point of type
%    is the same physical size in all five figures.  Open each picture with
%    \scCanvas{<height in cm>} to pin the bounding box.
% 2. Colour is semantic, never decorative.  Four roles only:
%       scAccent (green)       structure, the normal/expected path
%       scMark   (orange)      catalyst, certificate, or highlight
%       scAlert  (vermillion)  failure, violation, or uncovered set
%       scRef    (blue)        secondary or reference series
%    A panel gets at most one accent.  Titles are always ink; panel identity
%    is also encoded by a letter inside a circular badge, not by colour alone.
% 3. Type scale.  Three sizes, no others, and nothing below 6 pt:
%       \scTitle (7 pt bold)  panel titles
%       \scBody  (6.5 pt)     labels, values, in-figure prose
%       \scSmall (6 pt)       units, footnotes, dense graph labels
% ---------------------------------------------------------------------------
% Load guard.  The paper inputs every figure source inline, so this file is
% reached once per figure; without the guard the second \newlength below
% would abort the run.
\ifdefined\SynKitStyleLoaded\endinput\fi
\newcommand{\SynKitStyleLoaded}{}

% Load in the document preamble. XeLaTeX/LuaLaTeX and an installed Arial
% family are required; fontspec reports an error instead of substituting.
\RequirePackage[no-math]{fontspec}
\newfontfamily\scArial{Arial}[NFSSFamily=scArial,Ligatures=TeX]
% Keep manuscript mathematics unchanged; select Arial letters and numerals
% only inside figures. TeX supplies specialist mathematical symbols.
\RequirePackage[subdued,italic]{mathastext}
\MTDeclareVersion{synarial}{TU}{scArial}{m}{n}
\SetMathAlphabet{\Mathsf}{synarial}{TU}{scArial}{m}{n}
\SetMathAlphabet{\Mathtt}{synarial}{TU}{scArial}{m}{n}

\usetikzlibrary{positioning,calc,fit,backgrounds,arrows.meta,shapes.geometric,
                shapes.misc,decorations.pathreplacing,patterns,math}

% --- canvas ---------------------------------------------------------------
% 372 pt text width minus the 1 pt standalone border on each side.
\newlength{\scW}\setlength{\scW}{13.0326cm}
\providecommand{\scWnum}{13.0326}
% \scCanvas{<height>} pins the picture to the full text width, y running down
% from 0, so every figure shares one coordinate convention.  It *sets* the
% bounding box rather than contributing to it, so the emitted page is exactly
% the text width and every figure is included at 1:1 -- otherwise a stray
% invisible extent (pgfplots reserves a little around an axis) rescales one
% figure relative to the others and the same \scSmall prints at two sizes.
% Anything drawn outside the canvas is therefore cropped: keep content inside,
% and check with
%     \providecommand{\scCanvas}[1]{\path[draw=red] (0,0) rectangle (\scWnum,-#1);}
\providecommand{\scCanvas}[1]{%
  \path[use as bounding box] (0,0) rectangle (\scWnum,-#1);}

% --- palette --------------------------------------------------------------
% Canonical companion-paper values (../../../../Style/style.tex).
\definecolor{npgOrange}{HTML}{E69F00}
\definecolor{npgSky}{HTML}{56B4E9}
\definecolor{npgGreen}{HTML}{009E73}
\definecolor{npgYellow}{HTML}{F0E442}
\definecolor{npgBlue}{HTML}{0072B2}
\definecolor{npgVermillion}{HTML}{E64B35}
\definecolor{npgPurple}{HTML}{CC79A7}
\definecolor{figInk}{HTML}{1A1A1A}
\definecolor{figMuted}{HTML}{666666}
\definecolor{figHair}{HTML}{D9D9D9}
\definecolor{figWash}{HTML}{F7F7F7}

\colorlet{scInk}{figInk}              % primary text
\colorlet{scInkM}{figMuted}           % secondary text
\colorlet{scPaper}{figWash}           % quiet cards and note strips
\colorlet{scSpF}{npgSky!12}           % species fill
\colorlet{scSpE}{npgBlue}             % species edge
\colorlet{scRxF}{npgOrange!12}        % reaction fill
\colorlet{scRxE}{npgOrange!88!black}  % reaction edge
\colorlet{scReac}{npgBlue!78!black}   % reactant edge
\colorlet{scProd}{npgGreen!82!black}  % product edge
\colorlet{scHiF}{npgOrange!13}        % highlight fill
\colorlet{scHiE}{npgOrange!82!black}  % highlight edge
% Stage / accent tints.  stDE remains available for the fifth plot series.
\colorlet{stAF}{npgGreen!9}\colorlet{stAE}{npgGreen}
\colorlet{stBF}{npgOrange!11}\colorlet{stBE}{npgOrange}
\colorlet{stCF}{npgBlue!9}\colorlet{stCE}{npgBlue}
\colorlet{stDF}{npgPurple!10}\colorlet{stDE}{npgPurple!85!black}
\colorlet{stEF}{npgVermillion!9}\colorlet{stEE}{npgVermillion}

% Semantic aliases.  Figures should speak in these, not in stAE/stBE/...
\colorlet{scAccent}{stAE}\colorlet{scAccentF}{stAF}   % structure / normal
\colorlet{scMark}{stBE}\colorlet{scMarkF}{stBF}       % highlight / certificate
\colorlet{scRef}{stCE}\colorlet{scRefF}{stCF}         % secondary series
\colorlet{scAlert}{stEE}\colorlet{scAlertF}{stEF}     % failure / violation
\colorlet{scRule}{figHair}                            % card and axis hairlines
\colorlet{scRuleL}{figHair!55}                        % grid, dividers

% --- text styles ----------------------------------------------------------
% Three sizes, fixed in points so they do not drift with the document class.
% Arial is shared by labels, headings, panel badges, and chemistry labels.
\providecommand{\scTitle}{\scArial\fontsize{7}{8.4}\selectfont\bfseries\color{scInk}}
\providecommand{\scBody}{\scArial\fontsize{6.5}{7.8}\selectfont\color{scInkM}}
\providecommand{\scSmall}{\scArial\fontsize{6}{7.2}\selectfont\color{scInkM}}
\providecommand{\scSub}[1]{{\scSmall\itshape #1}}
\providecommand{\scPanel}[1]{{\scTitle(#1)}}
% Numerals that carry a result: ink, not colour, so the eye reads the value.
\providecommand{\scNum}{\scArial\fontsize{10}{11}\selectfont\bfseries\color{scInk}}

% --- shared TikZ styles ---------------------------------------------------
\tikzset{
  % text never hyphenates inside a figure
  every node/.append style={execute at begin node={%
      \hyphenpenalty=10000 \exhyphenpenalty=10000 \hbadness=10000 }},
  syn/.style={font=\scArial,color=scInk,
              execute at begin picture={\MTversion*{synarial}},
              >={Stealth[round,length=2.4mm,width=1.9mm]}},
  % containers ------------------------------------------------------------
  % One card chrome for the whole paper: white, 3 pt radius, one hairline.
  panel/.style={rounded corners=2.5pt,fill=none,draw=scRule,line width=0.45pt},
  card/.style={rounded corners=2.2pt,fill=scPaper,draw=scRule,line width=0.45pt,
               inner sep=5pt,align=left,anchor=north west},
  % a quiet strip for a panel's closing fact
  note/.style={rounded corners=2.5pt,fill=scPaper,draw=scRuleL,line width=0.4pt,
               inner xsep=6pt,inner ysep=3.5pt,align=left,font=\scSmall},
  % small labelled value pill
  chip/.style={rounded corners=2pt,fill=scPaper,draw=scRule,line width=0.4pt,
               inner xsep=4pt,inner ysep=2pt,font=\scSmall},
  % Circular letter badge: the same accessible panel marker used in Style/*.
  hbullet/.style={circle,fill=scInk,text=white,draw=none,
                  font=\scArial\fontsize{6}{6}\selectfont\bfseries,
                  minimum size=4.2mm,inner sep=0pt},
  ctag/.style={font=\scTitle,anchor=north west,text=scInk},
  % CRN primitives --------------------------------------------------------
  species/.style={rounded rectangle,fill=scSpF,draw=scSpE,line width=0.45pt,
                  inner xsep=2.4pt,inner ysep=1.6pt,font=\scSmall,
                  text=scInk,minimum height=3.6mm},
  cofactor/.style={species,fill=scRefF,draw=scRef},
  catalyst/.style={species,fill=scHiF,draw=scHiE},
  rxn/.style={rectangle,fill=scRxF,draw=scRxE,line width=0.45pt,
              inner sep=1.2pt,font=\scSmall,text=scInk,
              minimum width=3.0mm,minimum height=3.0mm},
  place/.style={circle,fill=scSpF,draw=scSpE,line width=0.45pt,
                minimum size=4.0mm,inner sep=0pt},
  trans/.style={rectangle,fill=scRxE,draw=scRxE,minimum width=1.1mm,
                minimum height=3.6mm,inner sep=0pt},
  token/.style={circle,fill=scInk,minimum size=1.1mm,inner sep=0pt},
  % edges -----------------------------------------------------------------
  reac/.style={->,line width=0.5pt,draw=scReac,shorten >=0.8pt,shorten <=0.8pt},
  prod/.style={->,line width=0.5pt,draw=scProd,shorten >=0.8pt,shorten <=0.8pt},
  cat/.style={->,line width=0.5pt,draw=scHiE,shorten >=0.8pt,shorten <=0.8pt},
  flow/.style={-{Stealth[round,length=2.2mm,width=1.8mm]},line width=1.1pt,
               draw=scInkM!45},
  flab/.style={font=\scSmall,text=scInkM,inner sep=2pt},
  stoi/.style={font=\scSmall,text=scInkM,inner sep=0.8pt},
  % verdict marks: shape carries the meaning, colour only reinforces it
  vpass/.style={circle,fill=scAccentF,draw=scAccent,text=scAccent,
                font=\fontsize{6}{6}\selectfont\bfseries,
                minimum size=3.2mm,inner sep=0pt},
  vfail/.style={circle,fill=scAlertF,draw=scAlert,text=scAlert,
                font=\fontsize{6}{6}\selectfont\bfseries,
                minimum size=3.2mm,inner sep=0pt},
  callout/.style={font=\scSmall,text=scInkM,align=left},
}

% \scHead[<badge colour>]{<anchor coord>}{<letter>}{<title>}{<subtitle or empty>}
% One header treatment for every panel in the paper.
\providecommand{\scHead}[5][scInk]{%
  \node[hbullet,fill=#1,anchor=north west] (scbadge#3) at ($#2+(0,-0.015)$)
    {\MakeUppercase{#3}};
  \node[ctag,anchor=west] at ($(scbadge#3.east)+(0.14,0)$) {#4};
  \ifx\relax#5\relax\else
    \node[anchor=north west,font=\scSmall,text=scInkM]
      at ($#2+(0.56,-0.245)$) {#5};
  \fi}

% --- shared pgfplots look (only if the figure loaded pgfplots) ------------
\ifdefined\pgfplotsset
  \pgfplotsset{compat=1.18}
  \pgfplotsset{
    synaxis/.style={
      scale only axis,
      tick align=outside, tick pos=left,
      axis line style={draw=scInkM!55,line width=0.4pt},
      % One grid, on the value axis only.  No minor grid: at these sizes it
      % turns the plot frame into a grey wash.
      ymajorgrids=true, xmajorgrids=false,
      major grid style={draw=scRuleL,line width=0.3pt},
      major tick length=1.6pt,
      label style={font=\scBody,color=scInk},
      tick label style={font=\scSmall,color=scInkM},
      legend style={font=\scSmall,draw=none,fill=none,inner sep=1pt,
                    row sep=0.4pt,legend columns=1},
      legend cell align=left,
      every axis plot/.append style={line width=0.7pt,mark size=1.2pt},
    },
  }
\fi
}{% ---------------------------------------------------------------------------
% synstyle.tex --- shared look for every TikZ figure of the SynKit CRN paper.
%
% Input it from a figure file:
%     \IfFileExists{synstyle.tex}{% ---------------------------------------------------------------------------
% synstyle.tex --- shared look for every TikZ figure of the SynKit CRN paper.
%
% Input it from a figure file:
%     \IfFileExists{synstyle.tex}{\input{synstyle}}{\input{Figure/synstyle}}
% so that the figure compiles both from Figure/ and from the paper root.
%
% The visual vocabulary follows ../../../../Style/style.tex: the
% colour-blind-safe Okabe--Ito accents, quiet neutral rules and wash, circular
% panel badges, and a compact technical type scale are shared by the companion
% papers.  Semantic aliases below keep the CRN meaning explicit.
%
% DESIGN CONTRACT -----------------------------------------------------------
% 1. Geometry.  Every figure is laid out on a canvas exactly \scW wide (the
%    sn-jnl text width, 372 pt) and is included unscaled, so one point of type
%    is the same physical size in all five figures.  Open each picture with
%    \scCanvas{<height in cm>} to pin the bounding box.
% 2. Colour is semantic, never decorative.  Four roles only:
%       scAccent (green)       structure, the normal/expected path
%       scMark   (orange)      catalyst, certificate, or highlight
%       scAlert  (vermillion)  failure, violation, or uncovered set
%       scRef    (blue)        secondary or reference series
%    A panel gets at most one accent.  Titles are always ink; panel identity
%    is also encoded by a letter inside a circular badge, not by colour alone.
% 3. Type scale.  Three sizes, no others, and nothing below 6 pt:
%       \scTitle (7 pt bold)  panel titles
%       \scBody  (6.5 pt)     labels, values, in-figure prose
%       \scSmall (6 pt)       units, footnotes, dense graph labels
% ---------------------------------------------------------------------------
% Load guard.  The paper inputs every figure source inline, so this file is
% reached once per figure; without the guard the second \newlength below
% would abort the run.
\ifdefined\SynKitStyleLoaded\endinput\fi
\newcommand{\SynKitStyleLoaded}{}

% Load in the document preamble. XeLaTeX/LuaLaTeX and an installed Arial
% family are required; fontspec reports an error instead of substituting.
\RequirePackage[no-math]{fontspec}
\newfontfamily\scArial{Arial}[NFSSFamily=scArial,Ligatures=TeX]
% Keep manuscript mathematics unchanged; select Arial letters and numerals
% only inside figures. TeX supplies specialist mathematical symbols.
\RequirePackage[subdued,italic]{mathastext}
\MTDeclareVersion{synarial}{TU}{scArial}{m}{n}
\SetMathAlphabet{\Mathsf}{synarial}{TU}{scArial}{m}{n}
\SetMathAlphabet{\Mathtt}{synarial}{TU}{scArial}{m}{n}

\usetikzlibrary{positioning,calc,fit,backgrounds,arrows.meta,shapes.geometric,
                shapes.misc,decorations.pathreplacing,patterns,math}

% --- canvas ---------------------------------------------------------------
% 372 pt text width minus the 1 pt standalone border on each side.
\newlength{\scW}\setlength{\scW}{13.0326cm}
\providecommand{\scWnum}{13.0326}
% \scCanvas{<height>} pins the picture to the full text width, y running down
% from 0, so every figure shares one coordinate convention.  It *sets* the
% bounding box rather than contributing to it, so the emitted page is exactly
% the text width and every figure is included at 1:1 -- otherwise a stray
% invisible extent (pgfplots reserves a little around an axis) rescales one
% figure relative to the others and the same \scSmall prints at two sizes.
% Anything drawn outside the canvas is therefore cropped: keep content inside,
% and check with
%     \providecommand{\scCanvas}[1]{\path[draw=red] (0,0) rectangle (\scWnum,-#1);}
\providecommand{\scCanvas}[1]{%
  \path[use as bounding box] (0,0) rectangle (\scWnum,-#1);}

% --- palette --------------------------------------------------------------
% Canonical companion-paper values (../../../../Style/style.tex).
\definecolor{npgOrange}{HTML}{E69F00}
\definecolor{npgSky}{HTML}{56B4E9}
\definecolor{npgGreen}{HTML}{009E73}
\definecolor{npgYellow}{HTML}{F0E442}
\definecolor{npgBlue}{HTML}{0072B2}
\definecolor{npgVermillion}{HTML}{E64B35}
\definecolor{npgPurple}{HTML}{CC79A7}
\definecolor{figInk}{HTML}{1A1A1A}
\definecolor{figMuted}{HTML}{666666}
\definecolor{figHair}{HTML}{D9D9D9}
\definecolor{figWash}{HTML}{F7F7F7}

\colorlet{scInk}{figInk}              % primary text
\colorlet{scInkM}{figMuted}           % secondary text
\colorlet{scPaper}{figWash}           % quiet cards and note strips
\colorlet{scSpF}{npgSky!12}           % species fill
\colorlet{scSpE}{npgBlue}             % species edge
\colorlet{scRxF}{npgOrange!12}        % reaction fill
\colorlet{scRxE}{npgOrange!88!black}  % reaction edge
\colorlet{scReac}{npgBlue!78!black}   % reactant edge
\colorlet{scProd}{npgGreen!82!black}  % product edge
\colorlet{scHiF}{npgOrange!13}        % highlight fill
\colorlet{scHiE}{npgOrange!82!black}  % highlight edge
% Stage / accent tints.  stDE remains available for the fifth plot series.
\colorlet{stAF}{npgGreen!9}\colorlet{stAE}{npgGreen}
\colorlet{stBF}{npgOrange!11}\colorlet{stBE}{npgOrange}
\colorlet{stCF}{npgBlue!9}\colorlet{stCE}{npgBlue}
\colorlet{stDF}{npgPurple!10}\colorlet{stDE}{npgPurple!85!black}
\colorlet{stEF}{npgVermillion!9}\colorlet{stEE}{npgVermillion}

% Semantic aliases.  Figures should speak in these, not in stAE/stBE/...
\colorlet{scAccent}{stAE}\colorlet{scAccentF}{stAF}   % structure / normal
\colorlet{scMark}{stBE}\colorlet{scMarkF}{stBF}       % highlight / certificate
\colorlet{scRef}{stCE}\colorlet{scRefF}{stCF}         % secondary series
\colorlet{scAlert}{stEE}\colorlet{scAlertF}{stEF}     % failure / violation
\colorlet{scRule}{figHair}                            % card and axis hairlines
\colorlet{scRuleL}{figHair!55}                        % grid, dividers

% --- text styles ----------------------------------------------------------
% Three sizes, fixed in points so they do not drift with the document class.
% Arial is shared by labels, headings, panel badges, and chemistry labels.
\providecommand{\scTitle}{\scArial\fontsize{7}{8.4}\selectfont\bfseries\color{scInk}}
\providecommand{\scBody}{\scArial\fontsize{6.5}{7.8}\selectfont\color{scInkM}}
\providecommand{\scSmall}{\scArial\fontsize{6}{7.2}\selectfont\color{scInkM}}
\providecommand{\scSub}[1]{{\scSmall\itshape #1}}
\providecommand{\scPanel}[1]{{\scTitle(#1)}}
% Numerals that carry a result: ink, not colour, so the eye reads the value.
\providecommand{\scNum}{\scArial\fontsize{10}{11}\selectfont\bfseries\color{scInk}}

% --- shared TikZ styles ---------------------------------------------------
\tikzset{
  % text never hyphenates inside a figure
  every node/.append style={execute at begin node={%
      \hyphenpenalty=10000 \exhyphenpenalty=10000 \hbadness=10000 }},
  syn/.style={font=\scArial,color=scInk,
              execute at begin picture={\MTversion*{synarial}},
              >={Stealth[round,length=2.4mm,width=1.9mm]}},
  % containers ------------------------------------------------------------
  % One card chrome for the whole paper: white, 3 pt radius, one hairline.
  panel/.style={rounded corners=2.5pt,fill=none,draw=scRule,line width=0.45pt},
  card/.style={rounded corners=2.2pt,fill=scPaper,draw=scRule,line width=0.45pt,
               inner sep=5pt,align=left,anchor=north west},
  % a quiet strip for a panel's closing fact
  note/.style={rounded corners=2.5pt,fill=scPaper,draw=scRuleL,line width=0.4pt,
               inner xsep=6pt,inner ysep=3.5pt,align=left,font=\scSmall},
  % small labelled value pill
  chip/.style={rounded corners=2pt,fill=scPaper,draw=scRule,line width=0.4pt,
               inner xsep=4pt,inner ysep=2pt,font=\scSmall},
  % Circular letter badge: the same accessible panel marker used in Style/*.
  hbullet/.style={circle,fill=scInk,text=white,draw=none,
                  font=\scArial\fontsize{6}{6}\selectfont\bfseries,
                  minimum size=4.2mm,inner sep=0pt},
  ctag/.style={font=\scTitle,anchor=north west,text=scInk},
  % CRN primitives --------------------------------------------------------
  species/.style={rounded rectangle,fill=scSpF,draw=scSpE,line width=0.45pt,
                  inner xsep=2.4pt,inner ysep=1.6pt,font=\scSmall,
                  text=scInk,minimum height=3.6mm},
  cofactor/.style={species,fill=scRefF,draw=scRef},
  catalyst/.style={species,fill=scHiF,draw=scHiE},
  rxn/.style={rectangle,fill=scRxF,draw=scRxE,line width=0.45pt,
              inner sep=1.2pt,font=\scSmall,text=scInk,
              minimum width=3.0mm,minimum height=3.0mm},
  place/.style={circle,fill=scSpF,draw=scSpE,line width=0.45pt,
                minimum size=4.0mm,inner sep=0pt},
  trans/.style={rectangle,fill=scRxE,draw=scRxE,minimum width=1.1mm,
                minimum height=3.6mm,inner sep=0pt},
  token/.style={circle,fill=scInk,minimum size=1.1mm,inner sep=0pt},
  % edges -----------------------------------------------------------------
  reac/.style={->,line width=0.5pt,draw=scReac,shorten >=0.8pt,shorten <=0.8pt},
  prod/.style={->,line width=0.5pt,draw=scProd,shorten >=0.8pt,shorten <=0.8pt},
  cat/.style={->,line width=0.5pt,draw=scHiE,shorten >=0.8pt,shorten <=0.8pt},
  flow/.style={-{Stealth[round,length=2.2mm,width=1.8mm]},line width=1.1pt,
               draw=scInkM!45},
  flab/.style={font=\scSmall,text=scInkM,inner sep=2pt},
  stoi/.style={font=\scSmall,text=scInkM,inner sep=0.8pt},
  % verdict marks: shape carries the meaning, colour only reinforces it
  vpass/.style={circle,fill=scAccentF,draw=scAccent,text=scAccent,
                font=\fontsize{6}{6}\selectfont\bfseries,
                minimum size=3.2mm,inner sep=0pt},
  vfail/.style={circle,fill=scAlertF,draw=scAlert,text=scAlert,
                font=\fontsize{6}{6}\selectfont\bfseries,
                minimum size=3.2mm,inner sep=0pt},
  callout/.style={font=\scSmall,text=scInkM,align=left},
}

% \scHead[<badge colour>]{<anchor coord>}{<letter>}{<title>}{<subtitle or empty>}
% One header treatment for every panel in the paper.
\providecommand{\scHead}[5][scInk]{%
  \node[hbullet,fill=#1,anchor=north west] (scbadge#3) at ($#2+(0,-0.015)$)
    {\MakeUppercase{#3}};
  \node[ctag,anchor=west] at ($(scbadge#3.east)+(0.14,0)$) {#4};
  \ifx\relax#5\relax\else
    \node[anchor=north west,font=\scSmall,text=scInkM]
      at ($#2+(0.56,-0.245)$) {#5};
  \fi}

% --- shared pgfplots look (only if the figure loaded pgfplots) ------------
\ifdefined\pgfplotsset
  \pgfplotsset{compat=1.18}
  \pgfplotsset{
    synaxis/.style={
      scale only axis,
      tick align=outside, tick pos=left,
      axis line style={draw=scInkM!55,line width=0.4pt},
      % One grid, on the value axis only.  No minor grid: at these sizes it
      % turns the plot frame into a grey wash.
      ymajorgrids=true, xmajorgrids=false,
      major grid style={draw=scRuleL,line width=0.3pt},
      major tick length=1.6pt,
      label style={font=\scBody,color=scInk},
      tick label style={font=\scSmall,color=scInkM},
      legend style={font=\scSmall,draw=none,fill=none,inner sep=1pt,
                    row sep=0.4pt,legend columns=1},
      legend cell align=left,
      every axis plot/.append style={line width=0.7pt,mark size=1.2pt},
    },
  }
\fi
}{% ---------------------------------------------------------------------------
% synstyle.tex --- shared look for every TikZ figure of the SynKit CRN paper.
%
% Input it from a figure file:
%     \IfFileExists{synstyle.tex}{\input{synstyle}}{\input{Figure/synstyle}}
% so that the figure compiles both from Figure/ and from the paper root.
%
% The visual vocabulary follows ../../../../Style/style.tex: the
% colour-blind-safe Okabe--Ito accents, quiet neutral rules and wash, circular
% panel badges, and a compact technical type scale are shared by the companion
% papers.  Semantic aliases below keep the CRN meaning explicit.
%
% DESIGN CONTRACT -----------------------------------------------------------
% 1. Geometry.  Every figure is laid out on a canvas exactly \scW wide (the
%    sn-jnl text width, 372 pt) and is included unscaled, so one point of type
%    is the same physical size in all five figures.  Open each picture with
%    \scCanvas{<height in cm>} to pin the bounding box.
% 2. Colour is semantic, never decorative.  Four roles only:
%       scAccent (green)       structure, the normal/expected path
%       scMark   (orange)      catalyst, certificate, or highlight
%       scAlert  (vermillion)  failure, violation, or uncovered set
%       scRef    (blue)        secondary or reference series
%    A panel gets at most one accent.  Titles are always ink; panel identity
%    is also encoded by a letter inside a circular badge, not by colour alone.
% 3. Type scale.  Three sizes, no others, and nothing below 6 pt:
%       \scTitle (7 pt bold)  panel titles
%       \scBody  (6.5 pt)     labels, values, in-figure prose
%       \scSmall (6 pt)       units, footnotes, dense graph labels
% ---------------------------------------------------------------------------
% Load guard.  The paper inputs every figure source inline, so this file is
% reached once per figure; without the guard the second \newlength below
% would abort the run.
\ifdefined\SynKitStyleLoaded\endinput\fi
\newcommand{\SynKitStyleLoaded}{}

% Load in the document preamble. XeLaTeX/LuaLaTeX and an installed Arial
% family are required; fontspec reports an error instead of substituting.
\RequirePackage[no-math]{fontspec}
\newfontfamily\scArial{Arial}[NFSSFamily=scArial,Ligatures=TeX]
% Keep manuscript mathematics unchanged; select Arial letters and numerals
% only inside figures. TeX supplies specialist mathematical symbols.
\RequirePackage[subdued,italic]{mathastext}
\MTDeclareVersion{synarial}{TU}{scArial}{m}{n}
\SetMathAlphabet{\Mathsf}{synarial}{TU}{scArial}{m}{n}
\SetMathAlphabet{\Mathtt}{synarial}{TU}{scArial}{m}{n}

\usetikzlibrary{positioning,calc,fit,backgrounds,arrows.meta,shapes.geometric,
                shapes.misc,decorations.pathreplacing,patterns,math}

% --- canvas ---------------------------------------------------------------
% 372 pt text width minus the 1 pt standalone border on each side.
\newlength{\scW}\setlength{\scW}{13.0326cm}
\providecommand{\scWnum}{13.0326}
% \scCanvas{<height>} pins the picture to the full text width, y running down
% from 0, so every figure shares one coordinate convention.  It *sets* the
% bounding box rather than contributing to it, so the emitted page is exactly
% the text width and every figure is included at 1:1 -- otherwise a stray
% invisible extent (pgfplots reserves a little around an axis) rescales one
% figure relative to the others and the same \scSmall prints at two sizes.
% Anything drawn outside the canvas is therefore cropped: keep content inside,
% and check with
%     \providecommand{\scCanvas}[1]{\path[draw=red] (0,0) rectangle (\scWnum,-#1);}
\providecommand{\scCanvas}[1]{%
  \path[use as bounding box] (0,0) rectangle (\scWnum,-#1);}

% --- palette --------------------------------------------------------------
% Canonical companion-paper values (../../../../Style/style.tex).
\definecolor{npgOrange}{HTML}{E69F00}
\definecolor{npgSky}{HTML}{56B4E9}
\definecolor{npgGreen}{HTML}{009E73}
\definecolor{npgYellow}{HTML}{F0E442}
\definecolor{npgBlue}{HTML}{0072B2}
\definecolor{npgVermillion}{HTML}{E64B35}
\definecolor{npgPurple}{HTML}{CC79A7}
\definecolor{figInk}{HTML}{1A1A1A}
\definecolor{figMuted}{HTML}{666666}
\definecolor{figHair}{HTML}{D9D9D9}
\definecolor{figWash}{HTML}{F7F7F7}

\colorlet{scInk}{figInk}              % primary text
\colorlet{scInkM}{figMuted}           % secondary text
\colorlet{scPaper}{figWash}           % quiet cards and note strips
\colorlet{scSpF}{npgSky!12}           % species fill
\colorlet{scSpE}{npgBlue}             % species edge
\colorlet{scRxF}{npgOrange!12}        % reaction fill
\colorlet{scRxE}{npgOrange!88!black}  % reaction edge
\colorlet{scReac}{npgBlue!78!black}   % reactant edge
\colorlet{scProd}{npgGreen!82!black}  % product edge
\colorlet{scHiF}{npgOrange!13}        % highlight fill
\colorlet{scHiE}{npgOrange!82!black}  % highlight edge
% Stage / accent tints.  stDE remains available for the fifth plot series.
\colorlet{stAF}{npgGreen!9}\colorlet{stAE}{npgGreen}
\colorlet{stBF}{npgOrange!11}\colorlet{stBE}{npgOrange}
\colorlet{stCF}{npgBlue!9}\colorlet{stCE}{npgBlue}
\colorlet{stDF}{npgPurple!10}\colorlet{stDE}{npgPurple!85!black}
\colorlet{stEF}{npgVermillion!9}\colorlet{stEE}{npgVermillion}

% Semantic aliases.  Figures should speak in these, not in stAE/stBE/...
\colorlet{scAccent}{stAE}\colorlet{scAccentF}{stAF}   % structure / normal
\colorlet{scMark}{stBE}\colorlet{scMarkF}{stBF}       % highlight / certificate
\colorlet{scRef}{stCE}\colorlet{scRefF}{stCF}         % secondary series
\colorlet{scAlert}{stEE}\colorlet{scAlertF}{stEF}     % failure / violation
\colorlet{scRule}{figHair}                            % card and axis hairlines
\colorlet{scRuleL}{figHair!55}                        % grid, dividers

% --- text styles ----------------------------------------------------------
% Three sizes, fixed in points so they do not drift with the document class.
% Arial is shared by labels, headings, panel badges, and chemistry labels.
\providecommand{\scTitle}{\scArial\fontsize{7}{8.4}\selectfont\bfseries\color{scInk}}
\providecommand{\scBody}{\scArial\fontsize{6.5}{7.8}\selectfont\color{scInkM}}
\providecommand{\scSmall}{\scArial\fontsize{6}{7.2}\selectfont\color{scInkM}}
\providecommand{\scSub}[1]{{\scSmall\itshape #1}}
\providecommand{\scPanel}[1]{{\scTitle(#1)}}
% Numerals that carry a result: ink, not colour, so the eye reads the value.
\providecommand{\scNum}{\scArial\fontsize{10}{11}\selectfont\bfseries\color{scInk}}

% --- shared TikZ styles ---------------------------------------------------
\tikzset{
  % text never hyphenates inside a figure
  every node/.append style={execute at begin node={%
      \hyphenpenalty=10000 \exhyphenpenalty=10000 \hbadness=10000 }},
  syn/.style={font=\scArial,color=scInk,
              execute at begin picture={\MTversion*{synarial}},
              >={Stealth[round,length=2.4mm,width=1.9mm]}},
  % containers ------------------------------------------------------------
  % One card chrome for the whole paper: white, 3 pt radius, one hairline.
  panel/.style={rounded corners=2.5pt,fill=none,draw=scRule,line width=0.45pt},
  card/.style={rounded corners=2.2pt,fill=scPaper,draw=scRule,line width=0.45pt,
               inner sep=5pt,align=left,anchor=north west},
  % a quiet strip for a panel's closing fact
  note/.style={rounded corners=2.5pt,fill=scPaper,draw=scRuleL,line width=0.4pt,
               inner xsep=6pt,inner ysep=3.5pt,align=left,font=\scSmall},
  % small labelled value pill
  chip/.style={rounded corners=2pt,fill=scPaper,draw=scRule,line width=0.4pt,
               inner xsep=4pt,inner ysep=2pt,font=\scSmall},
  % Circular letter badge: the same accessible panel marker used in Style/*.
  hbullet/.style={circle,fill=scInk,text=white,draw=none,
                  font=\scArial\fontsize{6}{6}\selectfont\bfseries,
                  minimum size=4.2mm,inner sep=0pt},
  ctag/.style={font=\scTitle,anchor=north west,text=scInk},
  % CRN primitives --------------------------------------------------------
  species/.style={rounded rectangle,fill=scSpF,draw=scSpE,line width=0.45pt,
                  inner xsep=2.4pt,inner ysep=1.6pt,font=\scSmall,
                  text=scInk,minimum height=3.6mm},
  cofactor/.style={species,fill=scRefF,draw=scRef},
  catalyst/.style={species,fill=scHiF,draw=scHiE},
  rxn/.style={rectangle,fill=scRxF,draw=scRxE,line width=0.45pt,
              inner sep=1.2pt,font=\scSmall,text=scInk,
              minimum width=3.0mm,minimum height=3.0mm},
  place/.style={circle,fill=scSpF,draw=scSpE,line width=0.45pt,
                minimum size=4.0mm,inner sep=0pt},
  trans/.style={rectangle,fill=scRxE,draw=scRxE,minimum width=1.1mm,
                minimum height=3.6mm,inner sep=0pt},
  token/.style={circle,fill=scInk,minimum size=1.1mm,inner sep=0pt},
  % edges -----------------------------------------------------------------
  reac/.style={->,line width=0.5pt,draw=scReac,shorten >=0.8pt,shorten <=0.8pt},
  prod/.style={->,line width=0.5pt,draw=scProd,shorten >=0.8pt,shorten <=0.8pt},
  cat/.style={->,line width=0.5pt,draw=scHiE,shorten >=0.8pt,shorten <=0.8pt},
  flow/.style={-{Stealth[round,length=2.2mm,width=1.8mm]},line width=1.1pt,
               draw=scInkM!45},
  flab/.style={font=\scSmall,text=scInkM,inner sep=2pt},
  stoi/.style={font=\scSmall,text=scInkM,inner sep=0.8pt},
  % verdict marks: shape carries the meaning, colour only reinforces it
  vpass/.style={circle,fill=scAccentF,draw=scAccent,text=scAccent,
                font=\fontsize{6}{6}\selectfont\bfseries,
                minimum size=3.2mm,inner sep=0pt},
  vfail/.style={circle,fill=scAlertF,draw=scAlert,text=scAlert,
                font=\fontsize{6}{6}\selectfont\bfseries,
                minimum size=3.2mm,inner sep=0pt},
  callout/.style={font=\scSmall,text=scInkM,align=left},
}

% \scHead[<badge colour>]{<anchor coord>}{<letter>}{<title>}{<subtitle or empty>}
% One header treatment for every panel in the paper.
\providecommand{\scHead}[5][scInk]{%
  \node[hbullet,fill=#1,anchor=north west] (scbadge#3) at ($#2+(0,-0.015)$)
    {\MakeUppercase{#3}};
  \node[ctag,anchor=west] at ($(scbadge#3.east)+(0.14,0)$) {#4};
  \ifx\relax#5\relax\else
    \node[anchor=north west,font=\scSmall,text=scInkM]
      at ($#2+(0.56,-0.245)$) {#5};
  \fi}

% --- shared pgfplots look (only if the figure loaded pgfplots) ------------
\ifdefined\pgfplotsset
  \pgfplotsset{compat=1.18}
  \pgfplotsset{
    synaxis/.style={
      scale only axis,
      tick align=outside, tick pos=left,
      axis line style={draw=scInkM!55,line width=0.4pt},
      % One grid, on the value axis only.  No minor grid: at these sizes it
      % turns the plot frame into a grey wash.
      ymajorgrids=true, xmajorgrids=false,
      major grid style={draw=scRuleL,line width=0.3pt},
      major tick length=1.6pt,
      label style={font=\scBody,color=scInk},
      tick label style={font=\scSmall,color=scInkM},
      legend style={font=\scSmall,draw=none,fill=none,inner sep=1pt,
                    row sep=0.4pt,legend columns=1},
      legend cell align=left,
      every axis plot/.append style={line width=0.7pt,mark size=1.2pt},
    },
  }
\fi
}
% so that the figure compiles both from Figure/ and from the paper root.
%
% The visual vocabulary follows ../../../../Style/style.tex: the
% colour-blind-safe Okabe--Ito accents, quiet neutral rules and wash, circular
% panel badges, and a compact technical type scale are shared by the companion
% papers.  Semantic aliases below keep the CRN meaning explicit.
%
% DESIGN CONTRACT -----------------------------------------------------------
% 1. Geometry.  Every figure is laid out on a canvas exactly \scW wide (the
%    sn-jnl text width, 372 pt) and is included unscaled, so one point of type
%    is the same physical size in all five figures.  Open each picture with
%    \scCanvas{<height in cm>} to pin the bounding box.
% 2. Colour is semantic, never decorative.  Four roles only:
%       scAccent (green)       structure, the normal/expected path
%       scMark   (orange)      catalyst, certificate, or highlight
%       scAlert  (vermillion)  failure, violation, or uncovered set
%       scRef    (blue)        secondary or reference series
%    A panel gets at most one accent.  Titles are always ink; panel identity
%    is also encoded by a letter inside a circular badge, not by colour alone.
% 3. Type scale.  Three sizes, no others, and nothing below 6 pt:
%       \scTitle (7 pt bold)  panel titles
%       \scBody  (6.5 pt)     labels, values, in-figure prose
%       \scSmall (6 pt)       units, footnotes, dense graph labels
% ---------------------------------------------------------------------------
% Load guard.  The paper inputs every figure source inline, so this file is
% reached once per figure; without the guard the second \newlength below
% would abort the run.
\ifdefined\SynKitStyleLoaded\endinput\fi
\newcommand{\SynKitStyleLoaded}{}

% Load in the document preamble. XeLaTeX/LuaLaTeX and an installed Arial
% family are required; fontspec reports an error instead of substituting.
\RequirePackage[no-math]{fontspec}
\newfontfamily\scArial{Arial}[NFSSFamily=scArial,Ligatures=TeX]
% Keep manuscript mathematics unchanged; select Arial letters and numerals
% only inside figures. TeX supplies specialist mathematical symbols.
\RequirePackage[subdued,italic]{mathastext}
\MTDeclareVersion{synarial}{TU}{scArial}{m}{n}
\SetMathAlphabet{\Mathsf}{synarial}{TU}{scArial}{m}{n}
\SetMathAlphabet{\Mathtt}{synarial}{TU}{scArial}{m}{n}

\usetikzlibrary{positioning,calc,fit,backgrounds,arrows.meta,shapes.geometric,
                shapes.misc,decorations.pathreplacing,patterns,math}

% --- canvas ---------------------------------------------------------------
% 372 pt text width minus the 1 pt standalone border on each side.
\newlength{\scW}\setlength{\scW}{13.0326cm}
\providecommand{\scWnum}{13.0326}
% \scCanvas{<height>} pins the picture to the full text width, y running down
% from 0, so every figure shares one coordinate convention.  It *sets* the
% bounding box rather than contributing to it, so the emitted page is exactly
% the text width and every figure is included at 1:1 -- otherwise a stray
% invisible extent (pgfplots reserves a little around an axis) rescales one
% figure relative to the others and the same \scSmall prints at two sizes.
% Anything drawn outside the canvas is therefore cropped: keep content inside,
% and check with
%     \providecommand{\scCanvas}[1]{\path[draw=red] (0,0) rectangle (\scWnum,-#1);}
\providecommand{\scCanvas}[1]{%
  \path[use as bounding box] (0,0) rectangle (\scWnum,-#1);}

% --- palette --------------------------------------------------------------
% Canonical companion-paper values (../../../../Style/style.tex).
\definecolor{npgOrange}{HTML}{E69F00}
\definecolor{npgSky}{HTML}{56B4E9}
\definecolor{npgGreen}{HTML}{009E73}
\definecolor{npgYellow}{HTML}{F0E442}
\definecolor{npgBlue}{HTML}{0072B2}
\definecolor{npgVermillion}{HTML}{E64B35}
\definecolor{npgPurple}{HTML}{CC79A7}
\definecolor{figInk}{HTML}{1A1A1A}
\definecolor{figMuted}{HTML}{666666}
\definecolor{figHair}{HTML}{D9D9D9}
\definecolor{figWash}{HTML}{F7F7F7}

\colorlet{scInk}{figInk}              % primary text
\colorlet{scInkM}{figMuted}           % secondary text
\colorlet{scPaper}{figWash}           % quiet cards and note strips
\colorlet{scSpF}{npgSky!12}           % species fill
\colorlet{scSpE}{npgBlue}             % species edge
\colorlet{scRxF}{npgOrange!12}        % reaction fill
\colorlet{scRxE}{npgOrange!88!black}  % reaction edge
\colorlet{scReac}{npgBlue!78!black}   % reactant edge
\colorlet{scProd}{npgGreen!82!black}  % product edge
\colorlet{scHiF}{npgOrange!13}        % highlight fill
\colorlet{scHiE}{npgOrange!82!black}  % highlight edge
% Stage / accent tints.  stDE remains available for the fifth plot series.
\colorlet{stAF}{npgGreen!9}\colorlet{stAE}{npgGreen}
\colorlet{stBF}{npgOrange!11}\colorlet{stBE}{npgOrange}
\colorlet{stCF}{npgBlue!9}\colorlet{stCE}{npgBlue}
\colorlet{stDF}{npgPurple!10}\colorlet{stDE}{npgPurple!85!black}
\colorlet{stEF}{npgVermillion!9}\colorlet{stEE}{npgVermillion}

% Semantic aliases.  Figures should speak in these, not in stAE/stBE/...
\colorlet{scAccent}{stAE}\colorlet{scAccentF}{stAF}   % structure / normal
\colorlet{scMark}{stBE}\colorlet{scMarkF}{stBF}       % highlight / certificate
\colorlet{scRef}{stCE}\colorlet{scRefF}{stCF}         % secondary series
\colorlet{scAlert}{stEE}\colorlet{scAlertF}{stEF}     % failure / violation
\colorlet{scRule}{figHair}                            % card and axis hairlines
\colorlet{scRuleL}{figHair!55}                        % grid, dividers

% --- text styles ----------------------------------------------------------
% Three sizes, fixed in points so they do not drift with the document class.
% Arial is shared by labels, headings, panel badges, and chemistry labels.
\providecommand{\scTitle}{\scArial\fontsize{7}{8.4}\selectfont\bfseries\color{scInk}}
\providecommand{\scBody}{\scArial\fontsize{6.5}{7.8}\selectfont\color{scInkM}}
\providecommand{\scSmall}{\scArial\fontsize{6}{7.2}\selectfont\color{scInkM}}
\providecommand{\scSub}[1]{{\scSmall\itshape #1}}
\providecommand{\scPanel}[1]{{\scTitle(#1)}}
% Numerals that carry a result: ink, not colour, so the eye reads the value.
\providecommand{\scNum}{\scArial\fontsize{10}{11}\selectfont\bfseries\color{scInk}}

% --- shared TikZ styles ---------------------------------------------------
\tikzset{
  % text never hyphenates inside a figure
  every node/.append style={execute at begin node={%
      \hyphenpenalty=10000 \exhyphenpenalty=10000 \hbadness=10000 }},
  syn/.style={font=\scArial,color=scInk,
              execute at begin picture={\MTversion*{synarial}},
              >={Stealth[round,length=2.4mm,width=1.9mm]}},
  % containers ------------------------------------------------------------
  % One card chrome for the whole paper: white, 3 pt radius, one hairline.
  panel/.style={rounded corners=2.5pt,fill=none,draw=scRule,line width=0.45pt},
  card/.style={rounded corners=2.2pt,fill=scPaper,draw=scRule,line width=0.45pt,
               inner sep=5pt,align=left,anchor=north west},
  % a quiet strip for a panel's closing fact
  note/.style={rounded corners=2.5pt,fill=scPaper,draw=scRuleL,line width=0.4pt,
               inner xsep=6pt,inner ysep=3.5pt,align=left,font=\scSmall},
  % small labelled value pill
  chip/.style={rounded corners=2pt,fill=scPaper,draw=scRule,line width=0.4pt,
               inner xsep=4pt,inner ysep=2pt,font=\scSmall},
  % Circular letter badge: the same accessible panel marker used in Style/*.
  hbullet/.style={circle,fill=scInk,text=white,draw=none,
                  font=\scArial\fontsize{6}{6}\selectfont\bfseries,
                  minimum size=4.2mm,inner sep=0pt},
  ctag/.style={font=\scTitle,anchor=north west,text=scInk},
  % CRN primitives --------------------------------------------------------
  species/.style={rounded rectangle,fill=scSpF,draw=scSpE,line width=0.45pt,
                  inner xsep=2.4pt,inner ysep=1.6pt,font=\scSmall,
                  text=scInk,minimum height=3.6mm},
  cofactor/.style={species,fill=scRefF,draw=scRef},
  catalyst/.style={species,fill=scHiF,draw=scHiE},
  rxn/.style={rectangle,fill=scRxF,draw=scRxE,line width=0.45pt,
              inner sep=1.2pt,font=\scSmall,text=scInk,
              minimum width=3.0mm,minimum height=3.0mm},
  place/.style={circle,fill=scSpF,draw=scSpE,line width=0.45pt,
                minimum size=4.0mm,inner sep=0pt},
  trans/.style={rectangle,fill=scRxE,draw=scRxE,minimum width=1.1mm,
                minimum height=3.6mm,inner sep=0pt},
  token/.style={circle,fill=scInk,minimum size=1.1mm,inner sep=0pt},
  % edges -----------------------------------------------------------------
  reac/.style={->,line width=0.5pt,draw=scReac,shorten >=0.8pt,shorten <=0.8pt},
  prod/.style={->,line width=0.5pt,draw=scProd,shorten >=0.8pt,shorten <=0.8pt},
  cat/.style={->,line width=0.5pt,draw=scHiE,shorten >=0.8pt,shorten <=0.8pt},
  flow/.style={-{Stealth[round,length=2.2mm,width=1.8mm]},line width=1.1pt,
               draw=scInkM!45},
  flab/.style={font=\scSmall,text=scInkM,inner sep=2pt},
  stoi/.style={font=\scSmall,text=scInkM,inner sep=0.8pt},
  % verdict marks: shape carries the meaning, colour only reinforces it
  vpass/.style={circle,fill=scAccentF,draw=scAccent,text=scAccent,
                font=\fontsize{6}{6}\selectfont\bfseries,
                minimum size=3.2mm,inner sep=0pt},
  vfail/.style={circle,fill=scAlertF,draw=scAlert,text=scAlert,
                font=\fontsize{6}{6}\selectfont\bfseries,
                minimum size=3.2mm,inner sep=0pt},
  callout/.style={font=\scSmall,text=scInkM,align=left},
}

% \scHead[<badge colour>]{<anchor coord>}{<letter>}{<title>}{<subtitle or empty>}
% One header treatment for every panel in the paper.
\providecommand{\scHead}[5][scInk]{%
  \node[hbullet,fill=#1,anchor=north west] (scbadge#3) at ($#2+(0,-0.015)$)
    {\MakeUppercase{#3}};
  \node[ctag,anchor=west] at ($(scbadge#3.east)+(0.14,0)$) {#4};
  \ifx\relax#5\relax\else
    \node[anchor=north west,font=\scSmall,text=scInkM]
      at ($#2+(0.56,-0.245)$) {#5};
  \fi}

% --- shared pgfplots look (only if the figure loaded pgfplots) ------------
\ifdefined\pgfplotsset
  \pgfplotsset{compat=1.18}
  \pgfplotsset{
    synaxis/.style={
      scale only axis,
      tick align=outside, tick pos=left,
      axis line style={draw=scInkM!55,line width=0.4pt},
      % One grid, on the value axis only.  No minor grid: at these sizes it
      % turns the plot frame into a grey wash.
      ymajorgrids=true, xmajorgrids=false,
      major grid style={draw=scRuleL,line width=0.3pt},
      major tick length=1.6pt,
      label style={font=\scBody,color=scInk},
      tick label style={font=\scSmall,color=scInkM},
      legend style={font=\scSmall,draw=none,fill=none,inner sep=1pt,
                    row sep=0.4pt,legend columns=1},
      legend cell align=left,
      every axis plot/.append style={line width=0.7pt,mark size=1.2pt},
    },
  }
\fi
}
% so that the figure compiles both from Figure/ and from the paper root.
%
% The visual vocabulary follows ../../../../Style/style.tex: the
% colour-blind-safe Okabe--Ito accents, quiet neutral rules and wash, circular
% panel badges, and a compact technical type scale are shared by the companion
% papers.  Semantic aliases below keep the CRN meaning explicit.
%
% DESIGN CONTRACT -----------------------------------------------------------
% 1. Geometry.  Every figure is laid out on a canvas exactly \scW wide (the
%    sn-jnl text width, 372 pt) and is included unscaled, so one point of type
%    is the same physical size in all five figures.  Open each picture with
%    \scCanvas{<height in cm>} to pin the bounding box.
% 2. Colour is semantic, never decorative.  Four roles only:
%       scAccent (green)       structure, the normal/expected path
%       scMark   (orange)      catalyst, certificate, or highlight
%       scAlert  (vermillion)  failure, violation, or uncovered set
%       scRef    (blue)        secondary or reference series
%    A panel gets at most one accent.  Titles are always ink; panel identity
%    is also encoded by a letter inside a circular badge, not by colour alone.
% 3. Type scale.  Three sizes, no others, and nothing below 6 pt:
%       \scTitle (7 pt bold)  panel titles
%       \scBody  (6.5 pt)     labels, values, in-figure prose
%       \scSmall (6 pt)       units, footnotes, dense graph labels
% ---------------------------------------------------------------------------
% Load guard.  The paper inputs every figure source inline, so this file is
% reached once per figure; without the guard the second \newlength below
% would abort the run.
\ifdefined\SynKitStyleLoaded\endinput\fi
\newcommand{\SynKitStyleLoaded}{}

% Load in the document preamble. XeLaTeX/LuaLaTeX and an installed Arial
% family are required; fontspec reports an error instead of substituting.
\RequirePackage[no-math]{fontspec}
\newfontfamily\scArial{Arial}[NFSSFamily=scArial,Ligatures=TeX]
% Keep manuscript mathematics unchanged; select Arial letters and numerals
% only inside figures. TeX supplies specialist mathematical symbols.
\RequirePackage[subdued,italic]{mathastext}
\MTDeclareVersion{synarial}{TU}{scArial}{m}{n}
\SetMathAlphabet{\Mathsf}{synarial}{TU}{scArial}{m}{n}
\SetMathAlphabet{\Mathtt}{synarial}{TU}{scArial}{m}{n}

\usetikzlibrary{positioning,calc,fit,backgrounds,arrows.meta,shapes.geometric,
                shapes.misc,decorations.pathreplacing,patterns,math}

% --- canvas ---------------------------------------------------------------
% 372 pt text width minus the 1 pt standalone border on each side.
\newlength{\scW}\setlength{\scW}{13.0326cm}
\providecommand{\scWnum}{13.0326}
% \scCanvas{<height>} pins the picture to the full text width, y running down
% from 0, so every figure shares one coordinate convention.  It *sets* the
% bounding box rather than contributing to it, so the emitted page is exactly
% the text width and every figure is included at 1:1 -- otherwise a stray
% invisible extent (pgfplots reserves a little around an axis) rescales one
% figure relative to the others and the same \scSmall prints at two sizes.
% Anything drawn outside the canvas is therefore cropped: keep content inside,
% and check with
%     \providecommand{\scCanvas}[1]{\path[draw=red] (0,0) rectangle (\scWnum,-#1);}
\providecommand{\scCanvas}[1]{%
  \path[use as bounding box] (0,0) rectangle (\scWnum,-#1);}

% --- palette --------------------------------------------------------------
% Canonical companion-paper values (../../../../Style/style.tex).
\definecolor{npgOrange}{HTML}{E69F00}
\definecolor{npgSky}{HTML}{56B4E9}
\definecolor{npgGreen}{HTML}{009E73}
\definecolor{npgYellow}{HTML}{F0E442}
\definecolor{npgBlue}{HTML}{0072B2}
\definecolor{npgVermillion}{HTML}{E64B35}
\definecolor{npgPurple}{HTML}{CC79A7}
\definecolor{figInk}{HTML}{1A1A1A}
\definecolor{figMuted}{HTML}{666666}
\definecolor{figHair}{HTML}{D9D9D9}
\definecolor{figWash}{HTML}{F7F7F7}

\colorlet{scInk}{figInk}              % primary text
\colorlet{scInkM}{figMuted}           % secondary text
\colorlet{scPaper}{figWash}           % quiet cards and note strips
\colorlet{scSpF}{npgSky!12}           % species fill
\colorlet{scSpE}{npgBlue}             % species edge
\colorlet{scRxF}{npgOrange!12}        % reaction fill
\colorlet{scRxE}{npgOrange!88!black}  % reaction edge
\colorlet{scReac}{npgBlue!78!black}   % reactant edge
\colorlet{scProd}{npgGreen!82!black}  % product edge
\colorlet{scHiF}{npgOrange!13}        % highlight fill
\colorlet{scHiE}{npgOrange!82!black}  % highlight edge
% Stage / accent tints.  stDE remains available for the fifth plot series.
\colorlet{stAF}{npgGreen!9}\colorlet{stAE}{npgGreen}
\colorlet{stBF}{npgOrange!11}\colorlet{stBE}{npgOrange}
\colorlet{stCF}{npgBlue!9}\colorlet{stCE}{npgBlue}
\colorlet{stDF}{npgPurple!10}\colorlet{stDE}{npgPurple!85!black}
\colorlet{stEF}{npgVermillion!9}\colorlet{stEE}{npgVermillion}

% Semantic aliases.  Figures should speak in these, not in stAE/stBE/...
\colorlet{scAccent}{stAE}\colorlet{scAccentF}{stAF}   % structure / normal
\colorlet{scMark}{stBE}\colorlet{scMarkF}{stBF}       % highlight / certificate
\colorlet{scRef}{stCE}\colorlet{scRefF}{stCF}         % secondary series
\colorlet{scAlert}{stEE}\colorlet{scAlertF}{stEF}     % failure / violation
\colorlet{scRule}{figHair}                            % card and axis hairlines
\colorlet{scRuleL}{figHair!55}                        % grid, dividers

% --- text styles ----------------------------------------------------------
% Three sizes, fixed in points so they do not drift with the document class.
% Arial is shared by labels, headings, panel badges, and chemistry labels.
\providecommand{\scTitle}{\scArial\fontsize{7}{8.4}\selectfont\bfseries\color{scInk}}
\providecommand{\scBody}{\scArial\fontsize{6.5}{7.8}\selectfont\color{scInkM}}
\providecommand{\scSmall}{\scArial\fontsize{6}{7.2}\selectfont\color{scInkM}}
\providecommand{\scSub}[1]{{\scSmall\itshape #1}}
\providecommand{\scPanel}[1]{{\scTitle(#1)}}
% Numerals that carry a result: ink, not colour, so the eye reads the value.
\providecommand{\scNum}{\scArial\fontsize{10}{11}\selectfont\bfseries\color{scInk}}

% --- shared TikZ styles ---------------------------------------------------
\tikzset{
  % text never hyphenates inside a figure
  every node/.append style={execute at begin node={%
      \hyphenpenalty=10000 \exhyphenpenalty=10000 \hbadness=10000 }},
  syn/.style={font=\scArial,color=scInk,
              execute at begin picture={\MTversion*{synarial}},
              >={Stealth[round,length=2.4mm,width=1.9mm]}},
  % containers ------------------------------------------------------------
  % One card chrome for the whole paper: white, 3 pt radius, one hairline.
  panel/.style={rounded corners=2.5pt,fill=none,draw=scRule,line width=0.45pt},
  card/.style={rounded corners=2.2pt,fill=scPaper,draw=scRule,line width=0.45pt,
               inner sep=5pt,align=left,anchor=north west},
  % a quiet strip for a panel's closing fact
  note/.style={rounded corners=2.5pt,fill=scPaper,draw=scRuleL,line width=0.4pt,
               inner xsep=6pt,inner ysep=3.5pt,align=left,font=\scSmall},
  % small labelled value pill
  chip/.style={rounded corners=2pt,fill=scPaper,draw=scRule,line width=0.4pt,
               inner xsep=4pt,inner ysep=2pt,font=\scSmall},
  % Circular letter badge: the same accessible panel marker used in Style/*.
  hbullet/.style={circle,fill=scInk,text=white,draw=none,
                  font=\scArial\fontsize{6}{6}\selectfont\bfseries,
                  minimum size=4.2mm,inner sep=0pt},
  ctag/.style={font=\scTitle,anchor=north west,text=scInk},
  % CRN primitives --------------------------------------------------------
  species/.style={rounded rectangle,fill=scSpF,draw=scSpE,line width=0.45pt,
                  inner xsep=2.4pt,inner ysep=1.6pt,font=\scSmall,
                  text=scInk,minimum height=3.6mm},
  cofactor/.style={species,fill=scRefF,draw=scRef},
  catalyst/.style={species,fill=scHiF,draw=scHiE},
  rxn/.style={rectangle,fill=scRxF,draw=scRxE,line width=0.45pt,
              inner sep=1.2pt,font=\scSmall,text=scInk,
              minimum width=3.0mm,minimum height=3.0mm},
  place/.style={circle,fill=scSpF,draw=scSpE,line width=0.45pt,
                minimum size=4.0mm,inner sep=0pt},
  trans/.style={rectangle,fill=scRxE,draw=scRxE,minimum width=1.1mm,
                minimum height=3.6mm,inner sep=0pt},
  token/.style={circle,fill=scInk,minimum size=1.1mm,inner sep=0pt},
  % edges -----------------------------------------------------------------
  reac/.style={->,line width=0.5pt,draw=scReac,shorten >=0.8pt,shorten <=0.8pt},
  prod/.style={->,line width=0.5pt,draw=scProd,shorten >=0.8pt,shorten <=0.8pt},
  cat/.style={->,line width=0.5pt,draw=scHiE,shorten >=0.8pt,shorten <=0.8pt},
  flow/.style={-{Stealth[round,length=2.2mm,width=1.8mm]},line width=1.1pt,
               draw=scInkM!45},
  flab/.style={font=\scSmall,text=scInkM,inner sep=2pt},
  stoi/.style={font=\scSmall,text=scInkM,inner sep=0.8pt},
  % verdict marks: shape carries the meaning, colour only reinforces it
  vpass/.style={circle,fill=scAccentF,draw=scAccent,text=scAccent,
                font=\fontsize{6}{6}\selectfont\bfseries,
                minimum size=3.2mm,inner sep=0pt},
  vfail/.style={circle,fill=scAlertF,draw=scAlert,text=scAlert,
                font=\fontsize{6}{6}\selectfont\bfseries,
                minimum size=3.2mm,inner sep=0pt},
  callout/.style={font=\scSmall,text=scInkM,align=left},
}

% \scHead[<badge colour>]{<anchor coord>}{<letter>}{<title>}{<subtitle or empty>}
% One header treatment for every panel in the paper.
\providecommand{\scHead}[5][scInk]{%
  \node[hbullet,fill=#1,anchor=north west] (scbadge#3) at ($#2+(0,-0.015)$)
    {\MakeUppercase{#3}};
  \node[ctag,anchor=west] at ($(scbadge#3.east)+(0.14,0)$) {#4};
  \ifx\relax#5\relax\else
    \node[anchor=north west,font=\scSmall,text=scInkM]
      at ($#2+(0.56,-0.245)$) {#5};
  \fi}

% --- shared pgfplots look (only if the figure loaded pgfplots) ------------
\ifdefined\pgfplotsset
  \pgfplotsset{compat=1.18}
  \pgfplotsset{
    synaxis/.style={
      scale only axis,
      tick align=outside, tick pos=left,
      axis line style={draw=scInkM!55,line width=0.4pt},
      % One grid, on the value axis only.  No minor grid: at these sizes it
      % turns the plot frame into a grey wash.
      ymajorgrids=true, xmajorgrids=false,
      major grid style={draw=scRuleL,line width=0.3pt},
      major tick length=1.6pt,
      label style={font=\scBody,color=scInk},
      tick label style={font=\scSmall,color=scInkM},
      legend style={font=\scSmall,draw=none,fill=none,inner sep=1pt,
                    row sep=0.4pt,legend columns=1},
      legend cell align=left,
      every axis plot/.append style={line width=0.7pt,mark size=1.2pt},
    },
  }
\fi
}{%
    % ---------------------------------------------------------------------------
% synstyle.tex --- shared look for every TikZ figure of the SynKit CRN paper.
%
% Input it from a figure file:
%     \IfFileExists{synstyle.tex}{% ---------------------------------------------------------------------------
% synstyle.tex --- shared look for every TikZ figure of the SynKit CRN paper.
%
% Input it from a figure file:
%     \IfFileExists{synstyle.tex}{% ---------------------------------------------------------------------------
% synstyle.tex --- shared look for every TikZ figure of the SynKit CRN paper.
%
% Input it from a figure file:
%     \IfFileExists{synstyle.tex}{\input{synstyle}}{\input{Figure/synstyle}}
% so that the figure compiles both from Figure/ and from the paper root.
%
% The visual vocabulary follows ../../../../Style/style.tex: the
% colour-blind-safe Okabe--Ito accents, quiet neutral rules and wash, circular
% panel badges, and a compact technical type scale are shared by the companion
% papers.  Semantic aliases below keep the CRN meaning explicit.
%
% DESIGN CONTRACT -----------------------------------------------------------
% 1. Geometry.  Every figure is laid out on a canvas exactly \scW wide (the
%    sn-jnl text width, 372 pt) and is included unscaled, so one point of type
%    is the same physical size in all five figures.  Open each picture with
%    \scCanvas{<height in cm>} to pin the bounding box.
% 2. Colour is semantic, never decorative.  Four roles only:
%       scAccent (green)       structure, the normal/expected path
%       scMark   (orange)      catalyst, certificate, or highlight
%       scAlert  (vermillion)  failure, violation, or uncovered set
%       scRef    (blue)        secondary or reference series
%    A panel gets at most one accent.  Titles are always ink; panel identity
%    is also encoded by a letter inside a circular badge, not by colour alone.
% 3. Type scale.  Three sizes, no others, and nothing below 6 pt:
%       \scTitle (7 pt bold)  panel titles
%       \scBody  (6.5 pt)     labels, values, in-figure prose
%       \scSmall (6 pt)       units, footnotes, dense graph labels
% ---------------------------------------------------------------------------
% Load guard.  The paper inputs every figure source inline, so this file is
% reached once per figure; without the guard the second \newlength below
% would abort the run.
\ifdefined\SynKitStyleLoaded\endinput\fi
\newcommand{\SynKitStyleLoaded}{}

% Load in the document preamble. XeLaTeX/LuaLaTeX and an installed Arial
% family are required; fontspec reports an error instead of substituting.
\RequirePackage[no-math]{fontspec}
\newfontfamily\scArial{Arial}[NFSSFamily=scArial,Ligatures=TeX]
% Keep manuscript mathematics unchanged; select Arial letters and numerals
% only inside figures. TeX supplies specialist mathematical symbols.
\RequirePackage[subdued,italic]{mathastext}
\MTDeclareVersion{synarial}{TU}{scArial}{m}{n}
\SetMathAlphabet{\Mathsf}{synarial}{TU}{scArial}{m}{n}
\SetMathAlphabet{\Mathtt}{synarial}{TU}{scArial}{m}{n}

\usetikzlibrary{positioning,calc,fit,backgrounds,arrows.meta,shapes.geometric,
                shapes.misc,decorations.pathreplacing,patterns,math}

% --- canvas ---------------------------------------------------------------
% 372 pt text width minus the 1 pt standalone border on each side.
\newlength{\scW}\setlength{\scW}{13.0326cm}
\providecommand{\scWnum}{13.0326}
% \scCanvas{<height>} pins the picture to the full text width, y running down
% from 0, so every figure shares one coordinate convention.  It *sets* the
% bounding box rather than contributing to it, so the emitted page is exactly
% the text width and every figure is included at 1:1 -- otherwise a stray
% invisible extent (pgfplots reserves a little around an axis) rescales one
% figure relative to the others and the same \scSmall prints at two sizes.
% Anything drawn outside the canvas is therefore cropped: keep content inside,
% and check with
%     \providecommand{\scCanvas}[1]{\path[draw=red] (0,0) rectangle (\scWnum,-#1);}
\providecommand{\scCanvas}[1]{%
  \path[use as bounding box] (0,0) rectangle (\scWnum,-#1);}

% --- palette --------------------------------------------------------------
% Canonical companion-paper values (../../../../Style/style.tex).
\definecolor{npgOrange}{HTML}{E69F00}
\definecolor{npgSky}{HTML}{56B4E9}
\definecolor{npgGreen}{HTML}{009E73}
\definecolor{npgYellow}{HTML}{F0E442}
\definecolor{npgBlue}{HTML}{0072B2}
\definecolor{npgVermillion}{HTML}{E64B35}
\definecolor{npgPurple}{HTML}{CC79A7}
\definecolor{figInk}{HTML}{1A1A1A}
\definecolor{figMuted}{HTML}{666666}
\definecolor{figHair}{HTML}{D9D9D9}
\definecolor{figWash}{HTML}{F7F7F7}

\colorlet{scInk}{figInk}              % primary text
\colorlet{scInkM}{figMuted}           % secondary text
\colorlet{scPaper}{figWash}           % quiet cards and note strips
\colorlet{scSpF}{npgSky!12}           % species fill
\colorlet{scSpE}{npgBlue}             % species edge
\colorlet{scRxF}{npgOrange!12}        % reaction fill
\colorlet{scRxE}{npgOrange!88!black}  % reaction edge
\colorlet{scReac}{npgBlue!78!black}   % reactant edge
\colorlet{scProd}{npgGreen!82!black}  % product edge
\colorlet{scHiF}{npgOrange!13}        % highlight fill
\colorlet{scHiE}{npgOrange!82!black}  % highlight edge
% Stage / accent tints.  stDE remains available for the fifth plot series.
\colorlet{stAF}{npgGreen!9}\colorlet{stAE}{npgGreen}
\colorlet{stBF}{npgOrange!11}\colorlet{stBE}{npgOrange}
\colorlet{stCF}{npgBlue!9}\colorlet{stCE}{npgBlue}
\colorlet{stDF}{npgPurple!10}\colorlet{stDE}{npgPurple!85!black}
\colorlet{stEF}{npgVermillion!9}\colorlet{stEE}{npgVermillion}

% Semantic aliases.  Figures should speak in these, not in stAE/stBE/...
\colorlet{scAccent}{stAE}\colorlet{scAccentF}{stAF}   % structure / normal
\colorlet{scMark}{stBE}\colorlet{scMarkF}{stBF}       % highlight / certificate
\colorlet{scRef}{stCE}\colorlet{scRefF}{stCF}         % secondary series
\colorlet{scAlert}{stEE}\colorlet{scAlertF}{stEF}     % failure / violation
\colorlet{scRule}{figHair}                            % card and axis hairlines
\colorlet{scRuleL}{figHair!55}                        % grid, dividers

% --- text styles ----------------------------------------------------------
% Three sizes, fixed in points so they do not drift with the document class.
% Arial is shared by labels, headings, panel badges, and chemistry labels.
\providecommand{\scTitle}{\scArial\fontsize{7}{8.4}\selectfont\bfseries\color{scInk}}
\providecommand{\scBody}{\scArial\fontsize{6.5}{7.8}\selectfont\color{scInkM}}
\providecommand{\scSmall}{\scArial\fontsize{6}{7.2}\selectfont\color{scInkM}}
\providecommand{\scSub}[1]{{\scSmall\itshape #1}}
\providecommand{\scPanel}[1]{{\scTitle(#1)}}
% Numerals that carry a result: ink, not colour, so the eye reads the value.
\providecommand{\scNum}{\scArial\fontsize{10}{11}\selectfont\bfseries\color{scInk}}

% --- shared TikZ styles ---------------------------------------------------
\tikzset{
  % text never hyphenates inside a figure
  every node/.append style={execute at begin node={%
      \hyphenpenalty=10000 \exhyphenpenalty=10000 \hbadness=10000 }},
  syn/.style={font=\scArial,color=scInk,
              execute at begin picture={\MTversion*{synarial}},
              >={Stealth[round,length=2.4mm,width=1.9mm]}},
  % containers ------------------------------------------------------------
  % One card chrome for the whole paper: white, 3 pt radius, one hairline.
  panel/.style={rounded corners=2.5pt,fill=none,draw=scRule,line width=0.45pt},
  card/.style={rounded corners=2.2pt,fill=scPaper,draw=scRule,line width=0.45pt,
               inner sep=5pt,align=left,anchor=north west},
  % a quiet strip for a panel's closing fact
  note/.style={rounded corners=2.5pt,fill=scPaper,draw=scRuleL,line width=0.4pt,
               inner xsep=6pt,inner ysep=3.5pt,align=left,font=\scSmall},
  % small labelled value pill
  chip/.style={rounded corners=2pt,fill=scPaper,draw=scRule,line width=0.4pt,
               inner xsep=4pt,inner ysep=2pt,font=\scSmall},
  % Circular letter badge: the same accessible panel marker used in Style/*.
  hbullet/.style={circle,fill=scInk,text=white,draw=none,
                  font=\scArial\fontsize{6}{6}\selectfont\bfseries,
                  minimum size=4.2mm,inner sep=0pt},
  ctag/.style={font=\scTitle,anchor=north west,text=scInk},
  % CRN primitives --------------------------------------------------------
  species/.style={rounded rectangle,fill=scSpF,draw=scSpE,line width=0.45pt,
                  inner xsep=2.4pt,inner ysep=1.6pt,font=\scSmall,
                  text=scInk,minimum height=3.6mm},
  cofactor/.style={species,fill=scRefF,draw=scRef},
  catalyst/.style={species,fill=scHiF,draw=scHiE},
  rxn/.style={rectangle,fill=scRxF,draw=scRxE,line width=0.45pt,
              inner sep=1.2pt,font=\scSmall,text=scInk,
              minimum width=3.0mm,minimum height=3.0mm},
  place/.style={circle,fill=scSpF,draw=scSpE,line width=0.45pt,
                minimum size=4.0mm,inner sep=0pt},
  trans/.style={rectangle,fill=scRxE,draw=scRxE,minimum width=1.1mm,
                minimum height=3.6mm,inner sep=0pt},
  token/.style={circle,fill=scInk,minimum size=1.1mm,inner sep=0pt},
  % edges -----------------------------------------------------------------
  reac/.style={->,line width=0.5pt,draw=scReac,shorten >=0.8pt,shorten <=0.8pt},
  prod/.style={->,line width=0.5pt,draw=scProd,shorten >=0.8pt,shorten <=0.8pt},
  cat/.style={->,line width=0.5pt,draw=scHiE,shorten >=0.8pt,shorten <=0.8pt},
  flow/.style={-{Stealth[round,length=2.2mm,width=1.8mm]},line width=1.1pt,
               draw=scInkM!45},
  flab/.style={font=\scSmall,text=scInkM,inner sep=2pt},
  stoi/.style={font=\scSmall,text=scInkM,inner sep=0.8pt},
  % verdict marks: shape carries the meaning, colour only reinforces it
  vpass/.style={circle,fill=scAccentF,draw=scAccent,text=scAccent,
                font=\fontsize{6}{6}\selectfont\bfseries,
                minimum size=3.2mm,inner sep=0pt},
  vfail/.style={circle,fill=scAlertF,draw=scAlert,text=scAlert,
                font=\fontsize{6}{6}\selectfont\bfseries,
                minimum size=3.2mm,inner sep=0pt},
  callout/.style={font=\scSmall,text=scInkM,align=left},
}

% \scHead[<badge colour>]{<anchor coord>}{<letter>}{<title>}{<subtitle or empty>}
% One header treatment for every panel in the paper.
\providecommand{\scHead}[5][scInk]{%
  \node[hbullet,fill=#1,anchor=north west] (scbadge#3) at ($#2+(0,-0.015)$)
    {\MakeUppercase{#3}};
  \node[ctag,anchor=west] at ($(scbadge#3.east)+(0.14,0)$) {#4};
  \ifx\relax#5\relax\else
    \node[anchor=north west,font=\scSmall,text=scInkM]
      at ($#2+(0.56,-0.245)$) {#5};
  \fi}

% --- shared pgfplots look (only if the figure loaded pgfplots) ------------
\ifdefined\pgfplotsset
  \pgfplotsset{compat=1.18}
  \pgfplotsset{
    synaxis/.style={
      scale only axis,
      tick align=outside, tick pos=left,
      axis line style={draw=scInkM!55,line width=0.4pt},
      % One grid, on the value axis only.  No minor grid: at these sizes it
      % turns the plot frame into a grey wash.
      ymajorgrids=true, xmajorgrids=false,
      major grid style={draw=scRuleL,line width=0.3pt},
      major tick length=1.6pt,
      label style={font=\scBody,color=scInk},
      tick label style={font=\scSmall,color=scInkM},
      legend style={font=\scSmall,draw=none,fill=none,inner sep=1pt,
                    row sep=0.4pt,legend columns=1},
      legend cell align=left,
      every axis plot/.append style={line width=0.7pt,mark size=1.2pt},
    },
  }
\fi
}{% ---------------------------------------------------------------------------
% synstyle.tex --- shared look for every TikZ figure of the SynKit CRN paper.
%
% Input it from a figure file:
%     \IfFileExists{synstyle.tex}{\input{synstyle}}{\input{Figure/synstyle}}
% so that the figure compiles both from Figure/ and from the paper root.
%
% The visual vocabulary follows ../../../../Style/style.tex: the
% colour-blind-safe Okabe--Ito accents, quiet neutral rules and wash, circular
% panel badges, and a compact technical type scale are shared by the companion
% papers.  Semantic aliases below keep the CRN meaning explicit.
%
% DESIGN CONTRACT -----------------------------------------------------------
% 1. Geometry.  Every figure is laid out on a canvas exactly \scW wide (the
%    sn-jnl text width, 372 pt) and is included unscaled, so one point of type
%    is the same physical size in all five figures.  Open each picture with
%    \scCanvas{<height in cm>} to pin the bounding box.
% 2. Colour is semantic, never decorative.  Four roles only:
%       scAccent (green)       structure, the normal/expected path
%       scMark   (orange)      catalyst, certificate, or highlight
%       scAlert  (vermillion)  failure, violation, or uncovered set
%       scRef    (blue)        secondary or reference series
%    A panel gets at most one accent.  Titles are always ink; panel identity
%    is also encoded by a letter inside a circular badge, not by colour alone.
% 3. Type scale.  Three sizes, no others, and nothing below 6 pt:
%       \scTitle (7 pt bold)  panel titles
%       \scBody  (6.5 pt)     labels, values, in-figure prose
%       \scSmall (6 pt)       units, footnotes, dense graph labels
% ---------------------------------------------------------------------------
% Load guard.  The paper inputs every figure source inline, so this file is
% reached once per figure; without the guard the second \newlength below
% would abort the run.
\ifdefined\SynKitStyleLoaded\endinput\fi
\newcommand{\SynKitStyleLoaded}{}

% Load in the document preamble. XeLaTeX/LuaLaTeX and an installed Arial
% family are required; fontspec reports an error instead of substituting.
\RequirePackage[no-math]{fontspec}
\newfontfamily\scArial{Arial}[NFSSFamily=scArial,Ligatures=TeX]
% Keep manuscript mathematics unchanged; select Arial letters and numerals
% only inside figures. TeX supplies specialist mathematical symbols.
\RequirePackage[subdued,italic]{mathastext}
\MTDeclareVersion{synarial}{TU}{scArial}{m}{n}
\SetMathAlphabet{\Mathsf}{synarial}{TU}{scArial}{m}{n}
\SetMathAlphabet{\Mathtt}{synarial}{TU}{scArial}{m}{n}

\usetikzlibrary{positioning,calc,fit,backgrounds,arrows.meta,shapes.geometric,
                shapes.misc,decorations.pathreplacing,patterns,math}

% --- canvas ---------------------------------------------------------------
% 372 pt text width minus the 1 pt standalone border on each side.
\newlength{\scW}\setlength{\scW}{13.0326cm}
\providecommand{\scWnum}{13.0326}
% \scCanvas{<height>} pins the picture to the full text width, y running down
% from 0, so every figure shares one coordinate convention.  It *sets* the
% bounding box rather than contributing to it, so the emitted page is exactly
% the text width and every figure is included at 1:1 -- otherwise a stray
% invisible extent (pgfplots reserves a little around an axis) rescales one
% figure relative to the others and the same \scSmall prints at two sizes.
% Anything drawn outside the canvas is therefore cropped: keep content inside,
% and check with
%     \providecommand{\scCanvas}[1]{\path[draw=red] (0,0) rectangle (\scWnum,-#1);}
\providecommand{\scCanvas}[1]{%
  \path[use as bounding box] (0,0) rectangle (\scWnum,-#1);}

% --- palette --------------------------------------------------------------
% Canonical companion-paper values (../../../../Style/style.tex).
\definecolor{npgOrange}{HTML}{E69F00}
\definecolor{npgSky}{HTML}{56B4E9}
\definecolor{npgGreen}{HTML}{009E73}
\definecolor{npgYellow}{HTML}{F0E442}
\definecolor{npgBlue}{HTML}{0072B2}
\definecolor{npgVermillion}{HTML}{E64B35}
\definecolor{npgPurple}{HTML}{CC79A7}
\definecolor{figInk}{HTML}{1A1A1A}
\definecolor{figMuted}{HTML}{666666}
\definecolor{figHair}{HTML}{D9D9D9}
\definecolor{figWash}{HTML}{F7F7F7}

\colorlet{scInk}{figInk}              % primary text
\colorlet{scInkM}{figMuted}           % secondary text
\colorlet{scPaper}{figWash}           % quiet cards and note strips
\colorlet{scSpF}{npgSky!12}           % species fill
\colorlet{scSpE}{npgBlue}             % species edge
\colorlet{scRxF}{npgOrange!12}        % reaction fill
\colorlet{scRxE}{npgOrange!88!black}  % reaction edge
\colorlet{scReac}{npgBlue!78!black}   % reactant edge
\colorlet{scProd}{npgGreen!82!black}  % product edge
\colorlet{scHiF}{npgOrange!13}        % highlight fill
\colorlet{scHiE}{npgOrange!82!black}  % highlight edge
% Stage / accent tints.  stDE remains available for the fifth plot series.
\colorlet{stAF}{npgGreen!9}\colorlet{stAE}{npgGreen}
\colorlet{stBF}{npgOrange!11}\colorlet{stBE}{npgOrange}
\colorlet{stCF}{npgBlue!9}\colorlet{stCE}{npgBlue}
\colorlet{stDF}{npgPurple!10}\colorlet{stDE}{npgPurple!85!black}
\colorlet{stEF}{npgVermillion!9}\colorlet{stEE}{npgVermillion}

% Semantic aliases.  Figures should speak in these, not in stAE/stBE/...
\colorlet{scAccent}{stAE}\colorlet{scAccentF}{stAF}   % structure / normal
\colorlet{scMark}{stBE}\colorlet{scMarkF}{stBF}       % highlight / certificate
\colorlet{scRef}{stCE}\colorlet{scRefF}{stCF}         % secondary series
\colorlet{scAlert}{stEE}\colorlet{scAlertF}{stEF}     % failure / violation
\colorlet{scRule}{figHair}                            % card and axis hairlines
\colorlet{scRuleL}{figHair!55}                        % grid, dividers

% --- text styles ----------------------------------------------------------
% Three sizes, fixed in points so they do not drift with the document class.
% Arial is shared by labels, headings, panel badges, and chemistry labels.
\providecommand{\scTitle}{\scArial\fontsize{7}{8.4}\selectfont\bfseries\color{scInk}}
\providecommand{\scBody}{\scArial\fontsize{6.5}{7.8}\selectfont\color{scInkM}}
\providecommand{\scSmall}{\scArial\fontsize{6}{7.2}\selectfont\color{scInkM}}
\providecommand{\scSub}[1]{{\scSmall\itshape #1}}
\providecommand{\scPanel}[1]{{\scTitle(#1)}}
% Numerals that carry a result: ink, not colour, so the eye reads the value.
\providecommand{\scNum}{\scArial\fontsize{10}{11}\selectfont\bfseries\color{scInk}}

% --- shared TikZ styles ---------------------------------------------------
\tikzset{
  % text never hyphenates inside a figure
  every node/.append style={execute at begin node={%
      \hyphenpenalty=10000 \exhyphenpenalty=10000 \hbadness=10000 }},
  syn/.style={font=\scArial,color=scInk,
              execute at begin picture={\MTversion*{synarial}},
              >={Stealth[round,length=2.4mm,width=1.9mm]}},
  % containers ------------------------------------------------------------
  % One card chrome for the whole paper: white, 3 pt radius, one hairline.
  panel/.style={rounded corners=2.5pt,fill=none,draw=scRule,line width=0.45pt},
  card/.style={rounded corners=2.2pt,fill=scPaper,draw=scRule,line width=0.45pt,
               inner sep=5pt,align=left,anchor=north west},
  % a quiet strip for a panel's closing fact
  note/.style={rounded corners=2.5pt,fill=scPaper,draw=scRuleL,line width=0.4pt,
               inner xsep=6pt,inner ysep=3.5pt,align=left,font=\scSmall},
  % small labelled value pill
  chip/.style={rounded corners=2pt,fill=scPaper,draw=scRule,line width=0.4pt,
               inner xsep=4pt,inner ysep=2pt,font=\scSmall},
  % Circular letter badge: the same accessible panel marker used in Style/*.
  hbullet/.style={circle,fill=scInk,text=white,draw=none,
                  font=\scArial\fontsize{6}{6}\selectfont\bfseries,
                  minimum size=4.2mm,inner sep=0pt},
  ctag/.style={font=\scTitle,anchor=north west,text=scInk},
  % CRN primitives --------------------------------------------------------
  species/.style={rounded rectangle,fill=scSpF,draw=scSpE,line width=0.45pt,
                  inner xsep=2.4pt,inner ysep=1.6pt,font=\scSmall,
                  text=scInk,minimum height=3.6mm},
  cofactor/.style={species,fill=scRefF,draw=scRef},
  catalyst/.style={species,fill=scHiF,draw=scHiE},
  rxn/.style={rectangle,fill=scRxF,draw=scRxE,line width=0.45pt,
              inner sep=1.2pt,font=\scSmall,text=scInk,
              minimum width=3.0mm,minimum height=3.0mm},
  place/.style={circle,fill=scSpF,draw=scSpE,line width=0.45pt,
                minimum size=4.0mm,inner sep=0pt},
  trans/.style={rectangle,fill=scRxE,draw=scRxE,minimum width=1.1mm,
                minimum height=3.6mm,inner sep=0pt},
  token/.style={circle,fill=scInk,minimum size=1.1mm,inner sep=0pt},
  % edges -----------------------------------------------------------------
  reac/.style={->,line width=0.5pt,draw=scReac,shorten >=0.8pt,shorten <=0.8pt},
  prod/.style={->,line width=0.5pt,draw=scProd,shorten >=0.8pt,shorten <=0.8pt},
  cat/.style={->,line width=0.5pt,draw=scHiE,shorten >=0.8pt,shorten <=0.8pt},
  flow/.style={-{Stealth[round,length=2.2mm,width=1.8mm]},line width=1.1pt,
               draw=scInkM!45},
  flab/.style={font=\scSmall,text=scInkM,inner sep=2pt},
  stoi/.style={font=\scSmall,text=scInkM,inner sep=0.8pt},
  % verdict marks: shape carries the meaning, colour only reinforces it
  vpass/.style={circle,fill=scAccentF,draw=scAccent,text=scAccent,
                font=\fontsize{6}{6}\selectfont\bfseries,
                minimum size=3.2mm,inner sep=0pt},
  vfail/.style={circle,fill=scAlertF,draw=scAlert,text=scAlert,
                font=\fontsize{6}{6}\selectfont\bfseries,
                minimum size=3.2mm,inner sep=0pt},
  callout/.style={font=\scSmall,text=scInkM,align=left},
}

% \scHead[<badge colour>]{<anchor coord>}{<letter>}{<title>}{<subtitle or empty>}
% One header treatment for every panel in the paper.
\providecommand{\scHead}[5][scInk]{%
  \node[hbullet,fill=#1,anchor=north west] (scbadge#3) at ($#2+(0,-0.015)$)
    {\MakeUppercase{#3}};
  \node[ctag,anchor=west] at ($(scbadge#3.east)+(0.14,0)$) {#4};
  \ifx\relax#5\relax\else
    \node[anchor=north west,font=\scSmall,text=scInkM]
      at ($#2+(0.56,-0.245)$) {#5};
  \fi}

% --- shared pgfplots look (only if the figure loaded pgfplots) ------------
\ifdefined\pgfplotsset
  \pgfplotsset{compat=1.18}
  \pgfplotsset{
    synaxis/.style={
      scale only axis,
      tick align=outside, tick pos=left,
      axis line style={draw=scInkM!55,line width=0.4pt},
      % One grid, on the value axis only.  No minor grid: at these sizes it
      % turns the plot frame into a grey wash.
      ymajorgrids=true, xmajorgrids=false,
      major grid style={draw=scRuleL,line width=0.3pt},
      major tick length=1.6pt,
      label style={font=\scBody,color=scInk},
      tick label style={font=\scSmall,color=scInkM},
      legend style={font=\scSmall,draw=none,fill=none,inner sep=1pt,
                    row sep=0.4pt,legend columns=1},
      legend cell align=left,
      every axis plot/.append style={line width=0.7pt,mark size=1.2pt},
    },
  }
\fi
}
% so that the figure compiles both from Figure/ and from the paper root.
%
% The visual vocabulary follows ../../../../Style/style.tex: the
% colour-blind-safe Okabe--Ito accents, quiet neutral rules and wash, circular
% panel badges, and a compact technical type scale are shared by the companion
% papers.  Semantic aliases below keep the CRN meaning explicit.
%
% DESIGN CONTRACT -----------------------------------------------------------
% 1. Geometry.  Every figure is laid out on a canvas exactly \scW wide (the
%    sn-jnl text width, 372 pt) and is included unscaled, so one point of type
%    is the same physical size in all five figures.  Open each picture with
%    \scCanvas{<height in cm>} to pin the bounding box.
% 2. Colour is semantic, never decorative.  Four roles only:
%       scAccent (green)       structure, the normal/expected path
%       scMark   (orange)      catalyst, certificate, or highlight
%       scAlert  (vermillion)  failure, violation, or uncovered set
%       scRef    (blue)        secondary or reference series
%    A panel gets at most one accent.  Titles are always ink; panel identity
%    is also encoded by a letter inside a circular badge, not by colour alone.
% 3. Type scale.  Three sizes, no others, and nothing below 6 pt:
%       \scTitle (7 pt bold)  panel titles
%       \scBody  (6.5 pt)     labels, values, in-figure prose
%       \scSmall (6 pt)       units, footnotes, dense graph labels
% ---------------------------------------------------------------------------
% Load guard.  The paper inputs every figure source inline, so this file is
% reached once per figure; without the guard the second \newlength below
% would abort the run.
\ifdefined\SynKitStyleLoaded\endinput\fi
\newcommand{\SynKitStyleLoaded}{}

% Load in the document preamble. XeLaTeX/LuaLaTeX and an installed Arial
% family are required; fontspec reports an error instead of substituting.
\RequirePackage[no-math]{fontspec}
\newfontfamily\scArial{Arial}[NFSSFamily=scArial,Ligatures=TeX]
% Keep manuscript mathematics unchanged; select Arial letters and numerals
% only inside figures. TeX supplies specialist mathematical symbols.
\RequirePackage[subdued,italic]{mathastext}
\MTDeclareVersion{synarial}{TU}{scArial}{m}{n}
\SetMathAlphabet{\Mathsf}{synarial}{TU}{scArial}{m}{n}
\SetMathAlphabet{\Mathtt}{synarial}{TU}{scArial}{m}{n}

\usetikzlibrary{positioning,calc,fit,backgrounds,arrows.meta,shapes.geometric,
                shapes.misc,decorations.pathreplacing,patterns,math}

% --- canvas ---------------------------------------------------------------
% 372 pt text width minus the 1 pt standalone border on each side.
\newlength{\scW}\setlength{\scW}{13.0326cm}
\providecommand{\scWnum}{13.0326}
% \scCanvas{<height>} pins the picture to the full text width, y running down
% from 0, so every figure shares one coordinate convention.  It *sets* the
% bounding box rather than contributing to it, so the emitted page is exactly
% the text width and every figure is included at 1:1 -- otherwise a stray
% invisible extent (pgfplots reserves a little around an axis) rescales one
% figure relative to the others and the same \scSmall prints at two sizes.
% Anything drawn outside the canvas is therefore cropped: keep content inside,
% and check with
%     \providecommand{\scCanvas}[1]{\path[draw=red] (0,0) rectangle (\scWnum,-#1);}
\providecommand{\scCanvas}[1]{%
  \path[use as bounding box] (0,0) rectangle (\scWnum,-#1);}

% --- palette --------------------------------------------------------------
% Canonical companion-paper values (../../../../Style/style.tex).
\definecolor{npgOrange}{HTML}{E69F00}
\definecolor{npgSky}{HTML}{56B4E9}
\definecolor{npgGreen}{HTML}{009E73}
\definecolor{npgYellow}{HTML}{F0E442}
\definecolor{npgBlue}{HTML}{0072B2}
\definecolor{npgVermillion}{HTML}{E64B35}
\definecolor{npgPurple}{HTML}{CC79A7}
\definecolor{figInk}{HTML}{1A1A1A}
\definecolor{figMuted}{HTML}{666666}
\definecolor{figHair}{HTML}{D9D9D9}
\definecolor{figWash}{HTML}{F7F7F7}

\colorlet{scInk}{figInk}              % primary text
\colorlet{scInkM}{figMuted}           % secondary text
\colorlet{scPaper}{figWash}           % quiet cards and note strips
\colorlet{scSpF}{npgSky!12}           % species fill
\colorlet{scSpE}{npgBlue}             % species edge
\colorlet{scRxF}{npgOrange!12}        % reaction fill
\colorlet{scRxE}{npgOrange!88!black}  % reaction edge
\colorlet{scReac}{npgBlue!78!black}   % reactant edge
\colorlet{scProd}{npgGreen!82!black}  % product edge
\colorlet{scHiF}{npgOrange!13}        % highlight fill
\colorlet{scHiE}{npgOrange!82!black}  % highlight edge
% Stage / accent tints.  stDE remains available for the fifth plot series.
\colorlet{stAF}{npgGreen!9}\colorlet{stAE}{npgGreen}
\colorlet{stBF}{npgOrange!11}\colorlet{stBE}{npgOrange}
\colorlet{stCF}{npgBlue!9}\colorlet{stCE}{npgBlue}
\colorlet{stDF}{npgPurple!10}\colorlet{stDE}{npgPurple!85!black}
\colorlet{stEF}{npgVermillion!9}\colorlet{stEE}{npgVermillion}

% Semantic aliases.  Figures should speak in these, not in stAE/stBE/...
\colorlet{scAccent}{stAE}\colorlet{scAccentF}{stAF}   % structure / normal
\colorlet{scMark}{stBE}\colorlet{scMarkF}{stBF}       % highlight / certificate
\colorlet{scRef}{stCE}\colorlet{scRefF}{stCF}         % secondary series
\colorlet{scAlert}{stEE}\colorlet{scAlertF}{stEF}     % failure / violation
\colorlet{scRule}{figHair}                            % card and axis hairlines
\colorlet{scRuleL}{figHair!55}                        % grid, dividers

% --- text styles ----------------------------------------------------------
% Three sizes, fixed in points so they do not drift with the document class.
% Arial is shared by labels, headings, panel badges, and chemistry labels.
\providecommand{\scTitle}{\scArial\fontsize{7}{8.4}\selectfont\bfseries\color{scInk}}
\providecommand{\scBody}{\scArial\fontsize{6.5}{7.8}\selectfont\color{scInkM}}
\providecommand{\scSmall}{\scArial\fontsize{6}{7.2}\selectfont\color{scInkM}}
\providecommand{\scSub}[1]{{\scSmall\itshape #1}}
\providecommand{\scPanel}[1]{{\scTitle(#1)}}
% Numerals that carry a result: ink, not colour, so the eye reads the value.
\providecommand{\scNum}{\scArial\fontsize{10}{11}\selectfont\bfseries\color{scInk}}

% --- shared TikZ styles ---------------------------------------------------
\tikzset{
  % text never hyphenates inside a figure
  every node/.append style={execute at begin node={%
      \hyphenpenalty=10000 \exhyphenpenalty=10000 \hbadness=10000 }},
  syn/.style={font=\scArial,color=scInk,
              execute at begin picture={\MTversion*{synarial}},
              >={Stealth[round,length=2.4mm,width=1.9mm]}},
  % containers ------------------------------------------------------------
  % One card chrome for the whole paper: white, 3 pt radius, one hairline.
  panel/.style={rounded corners=2.5pt,fill=none,draw=scRule,line width=0.45pt},
  card/.style={rounded corners=2.2pt,fill=scPaper,draw=scRule,line width=0.45pt,
               inner sep=5pt,align=left,anchor=north west},
  % a quiet strip for a panel's closing fact
  note/.style={rounded corners=2.5pt,fill=scPaper,draw=scRuleL,line width=0.4pt,
               inner xsep=6pt,inner ysep=3.5pt,align=left,font=\scSmall},
  % small labelled value pill
  chip/.style={rounded corners=2pt,fill=scPaper,draw=scRule,line width=0.4pt,
               inner xsep=4pt,inner ysep=2pt,font=\scSmall},
  % Circular letter badge: the same accessible panel marker used in Style/*.
  hbullet/.style={circle,fill=scInk,text=white,draw=none,
                  font=\scArial\fontsize{6}{6}\selectfont\bfseries,
                  minimum size=4.2mm,inner sep=0pt},
  ctag/.style={font=\scTitle,anchor=north west,text=scInk},
  % CRN primitives --------------------------------------------------------
  species/.style={rounded rectangle,fill=scSpF,draw=scSpE,line width=0.45pt,
                  inner xsep=2.4pt,inner ysep=1.6pt,font=\scSmall,
                  text=scInk,minimum height=3.6mm},
  cofactor/.style={species,fill=scRefF,draw=scRef},
  catalyst/.style={species,fill=scHiF,draw=scHiE},
  rxn/.style={rectangle,fill=scRxF,draw=scRxE,line width=0.45pt,
              inner sep=1.2pt,font=\scSmall,text=scInk,
              minimum width=3.0mm,minimum height=3.0mm},
  place/.style={circle,fill=scSpF,draw=scSpE,line width=0.45pt,
                minimum size=4.0mm,inner sep=0pt},
  trans/.style={rectangle,fill=scRxE,draw=scRxE,minimum width=1.1mm,
                minimum height=3.6mm,inner sep=0pt},
  token/.style={circle,fill=scInk,minimum size=1.1mm,inner sep=0pt},
  % edges -----------------------------------------------------------------
  reac/.style={->,line width=0.5pt,draw=scReac,shorten >=0.8pt,shorten <=0.8pt},
  prod/.style={->,line width=0.5pt,draw=scProd,shorten >=0.8pt,shorten <=0.8pt},
  cat/.style={->,line width=0.5pt,draw=scHiE,shorten >=0.8pt,shorten <=0.8pt},
  flow/.style={-{Stealth[round,length=2.2mm,width=1.8mm]},line width=1.1pt,
               draw=scInkM!45},
  flab/.style={font=\scSmall,text=scInkM,inner sep=2pt},
  stoi/.style={font=\scSmall,text=scInkM,inner sep=0.8pt},
  % verdict marks: shape carries the meaning, colour only reinforces it
  vpass/.style={circle,fill=scAccentF,draw=scAccent,text=scAccent,
                font=\fontsize{6}{6}\selectfont\bfseries,
                minimum size=3.2mm,inner sep=0pt},
  vfail/.style={circle,fill=scAlertF,draw=scAlert,text=scAlert,
                font=\fontsize{6}{6}\selectfont\bfseries,
                minimum size=3.2mm,inner sep=0pt},
  callout/.style={font=\scSmall,text=scInkM,align=left},
}

% \scHead[<badge colour>]{<anchor coord>}{<letter>}{<title>}{<subtitle or empty>}
% One header treatment for every panel in the paper.
\providecommand{\scHead}[5][scInk]{%
  \node[hbullet,fill=#1,anchor=north west] (scbadge#3) at ($#2+(0,-0.015)$)
    {\MakeUppercase{#3}};
  \node[ctag,anchor=west] at ($(scbadge#3.east)+(0.14,0)$) {#4};
  \ifx\relax#5\relax\else
    \node[anchor=north west,font=\scSmall,text=scInkM]
      at ($#2+(0.56,-0.245)$) {#5};
  \fi}

% --- shared pgfplots look (only if the figure loaded pgfplots) ------------
\ifdefined\pgfplotsset
  \pgfplotsset{compat=1.18}
  \pgfplotsset{
    synaxis/.style={
      scale only axis,
      tick align=outside, tick pos=left,
      axis line style={draw=scInkM!55,line width=0.4pt},
      % One grid, on the value axis only.  No minor grid: at these sizes it
      % turns the plot frame into a grey wash.
      ymajorgrids=true, xmajorgrids=false,
      major grid style={draw=scRuleL,line width=0.3pt},
      major tick length=1.6pt,
      label style={font=\scBody,color=scInk},
      tick label style={font=\scSmall,color=scInkM},
      legend style={font=\scSmall,draw=none,fill=none,inner sep=1pt,
                    row sep=0.4pt,legend columns=1},
      legend cell align=left,
      every axis plot/.append style={line width=0.7pt,mark size=1.2pt},
    },
  }
\fi
}{% ---------------------------------------------------------------------------
% synstyle.tex --- shared look for every TikZ figure of the SynKit CRN paper.
%
% Input it from a figure file:
%     \IfFileExists{synstyle.tex}{% ---------------------------------------------------------------------------
% synstyle.tex --- shared look for every TikZ figure of the SynKit CRN paper.
%
% Input it from a figure file:
%     \IfFileExists{synstyle.tex}{\input{synstyle}}{\input{Figure/synstyle}}
% so that the figure compiles both from Figure/ and from the paper root.
%
% The visual vocabulary follows ../../../../Style/style.tex: the
% colour-blind-safe Okabe--Ito accents, quiet neutral rules and wash, circular
% panel badges, and a compact technical type scale are shared by the companion
% papers.  Semantic aliases below keep the CRN meaning explicit.
%
% DESIGN CONTRACT -----------------------------------------------------------
% 1. Geometry.  Every figure is laid out on a canvas exactly \scW wide (the
%    sn-jnl text width, 372 pt) and is included unscaled, so one point of type
%    is the same physical size in all five figures.  Open each picture with
%    \scCanvas{<height in cm>} to pin the bounding box.
% 2. Colour is semantic, never decorative.  Four roles only:
%       scAccent (green)       structure, the normal/expected path
%       scMark   (orange)      catalyst, certificate, or highlight
%       scAlert  (vermillion)  failure, violation, or uncovered set
%       scRef    (blue)        secondary or reference series
%    A panel gets at most one accent.  Titles are always ink; panel identity
%    is also encoded by a letter inside a circular badge, not by colour alone.
% 3. Type scale.  Three sizes, no others, and nothing below 6 pt:
%       \scTitle (7 pt bold)  panel titles
%       \scBody  (6.5 pt)     labels, values, in-figure prose
%       \scSmall (6 pt)       units, footnotes, dense graph labels
% ---------------------------------------------------------------------------
% Load guard.  The paper inputs every figure source inline, so this file is
% reached once per figure; without the guard the second \newlength below
% would abort the run.
\ifdefined\SynKitStyleLoaded\endinput\fi
\newcommand{\SynKitStyleLoaded}{}

% Load in the document preamble. XeLaTeX/LuaLaTeX and an installed Arial
% family are required; fontspec reports an error instead of substituting.
\RequirePackage[no-math]{fontspec}
\newfontfamily\scArial{Arial}[NFSSFamily=scArial,Ligatures=TeX]
% Keep manuscript mathematics unchanged; select Arial letters and numerals
% only inside figures. TeX supplies specialist mathematical symbols.
\RequirePackage[subdued,italic]{mathastext}
\MTDeclareVersion{synarial}{TU}{scArial}{m}{n}
\SetMathAlphabet{\Mathsf}{synarial}{TU}{scArial}{m}{n}
\SetMathAlphabet{\Mathtt}{synarial}{TU}{scArial}{m}{n}

\usetikzlibrary{positioning,calc,fit,backgrounds,arrows.meta,shapes.geometric,
                shapes.misc,decorations.pathreplacing,patterns,math}

% --- canvas ---------------------------------------------------------------
% 372 pt text width minus the 1 pt standalone border on each side.
\newlength{\scW}\setlength{\scW}{13.0326cm}
\providecommand{\scWnum}{13.0326}
% \scCanvas{<height>} pins the picture to the full text width, y running down
% from 0, so every figure shares one coordinate convention.  It *sets* the
% bounding box rather than contributing to it, so the emitted page is exactly
% the text width and every figure is included at 1:1 -- otherwise a stray
% invisible extent (pgfplots reserves a little around an axis) rescales one
% figure relative to the others and the same \scSmall prints at two sizes.
% Anything drawn outside the canvas is therefore cropped: keep content inside,
% and check with
%     \providecommand{\scCanvas}[1]{\path[draw=red] (0,0) rectangle (\scWnum,-#1);}
\providecommand{\scCanvas}[1]{%
  \path[use as bounding box] (0,0) rectangle (\scWnum,-#1);}

% --- palette --------------------------------------------------------------
% Canonical companion-paper values (../../../../Style/style.tex).
\definecolor{npgOrange}{HTML}{E69F00}
\definecolor{npgSky}{HTML}{56B4E9}
\definecolor{npgGreen}{HTML}{009E73}
\definecolor{npgYellow}{HTML}{F0E442}
\definecolor{npgBlue}{HTML}{0072B2}
\definecolor{npgVermillion}{HTML}{E64B35}
\definecolor{npgPurple}{HTML}{CC79A7}
\definecolor{figInk}{HTML}{1A1A1A}
\definecolor{figMuted}{HTML}{666666}
\definecolor{figHair}{HTML}{D9D9D9}
\definecolor{figWash}{HTML}{F7F7F7}

\colorlet{scInk}{figInk}              % primary text
\colorlet{scInkM}{figMuted}           % secondary text
\colorlet{scPaper}{figWash}           % quiet cards and note strips
\colorlet{scSpF}{npgSky!12}           % species fill
\colorlet{scSpE}{npgBlue}             % species edge
\colorlet{scRxF}{npgOrange!12}        % reaction fill
\colorlet{scRxE}{npgOrange!88!black}  % reaction edge
\colorlet{scReac}{npgBlue!78!black}   % reactant edge
\colorlet{scProd}{npgGreen!82!black}  % product edge
\colorlet{scHiF}{npgOrange!13}        % highlight fill
\colorlet{scHiE}{npgOrange!82!black}  % highlight edge
% Stage / accent tints.  stDE remains available for the fifth plot series.
\colorlet{stAF}{npgGreen!9}\colorlet{stAE}{npgGreen}
\colorlet{stBF}{npgOrange!11}\colorlet{stBE}{npgOrange}
\colorlet{stCF}{npgBlue!9}\colorlet{stCE}{npgBlue}
\colorlet{stDF}{npgPurple!10}\colorlet{stDE}{npgPurple!85!black}
\colorlet{stEF}{npgVermillion!9}\colorlet{stEE}{npgVermillion}

% Semantic aliases.  Figures should speak in these, not in stAE/stBE/...
\colorlet{scAccent}{stAE}\colorlet{scAccentF}{stAF}   % structure / normal
\colorlet{scMark}{stBE}\colorlet{scMarkF}{stBF}       % highlight / certificate
\colorlet{scRef}{stCE}\colorlet{scRefF}{stCF}         % secondary series
\colorlet{scAlert}{stEE}\colorlet{scAlertF}{stEF}     % failure / violation
\colorlet{scRule}{figHair}                            % card and axis hairlines
\colorlet{scRuleL}{figHair!55}                        % grid, dividers

% --- text styles ----------------------------------------------------------
% Three sizes, fixed in points so they do not drift with the document class.
% Arial is shared by labels, headings, panel badges, and chemistry labels.
\providecommand{\scTitle}{\scArial\fontsize{7}{8.4}\selectfont\bfseries\color{scInk}}
\providecommand{\scBody}{\scArial\fontsize{6.5}{7.8}\selectfont\color{scInkM}}
\providecommand{\scSmall}{\scArial\fontsize{6}{7.2}\selectfont\color{scInkM}}
\providecommand{\scSub}[1]{{\scSmall\itshape #1}}
\providecommand{\scPanel}[1]{{\scTitle(#1)}}
% Numerals that carry a result: ink, not colour, so the eye reads the value.
\providecommand{\scNum}{\scArial\fontsize{10}{11}\selectfont\bfseries\color{scInk}}

% --- shared TikZ styles ---------------------------------------------------
\tikzset{
  % text never hyphenates inside a figure
  every node/.append style={execute at begin node={%
      \hyphenpenalty=10000 \exhyphenpenalty=10000 \hbadness=10000 }},
  syn/.style={font=\scArial,color=scInk,
              execute at begin picture={\MTversion*{synarial}},
              >={Stealth[round,length=2.4mm,width=1.9mm]}},
  % containers ------------------------------------------------------------
  % One card chrome for the whole paper: white, 3 pt radius, one hairline.
  panel/.style={rounded corners=2.5pt,fill=none,draw=scRule,line width=0.45pt},
  card/.style={rounded corners=2.2pt,fill=scPaper,draw=scRule,line width=0.45pt,
               inner sep=5pt,align=left,anchor=north west},
  % a quiet strip for a panel's closing fact
  note/.style={rounded corners=2.5pt,fill=scPaper,draw=scRuleL,line width=0.4pt,
               inner xsep=6pt,inner ysep=3.5pt,align=left,font=\scSmall},
  % small labelled value pill
  chip/.style={rounded corners=2pt,fill=scPaper,draw=scRule,line width=0.4pt,
               inner xsep=4pt,inner ysep=2pt,font=\scSmall},
  % Circular letter badge: the same accessible panel marker used in Style/*.
  hbullet/.style={circle,fill=scInk,text=white,draw=none,
                  font=\scArial\fontsize{6}{6}\selectfont\bfseries,
                  minimum size=4.2mm,inner sep=0pt},
  ctag/.style={font=\scTitle,anchor=north west,text=scInk},
  % CRN primitives --------------------------------------------------------
  species/.style={rounded rectangle,fill=scSpF,draw=scSpE,line width=0.45pt,
                  inner xsep=2.4pt,inner ysep=1.6pt,font=\scSmall,
                  text=scInk,minimum height=3.6mm},
  cofactor/.style={species,fill=scRefF,draw=scRef},
  catalyst/.style={species,fill=scHiF,draw=scHiE},
  rxn/.style={rectangle,fill=scRxF,draw=scRxE,line width=0.45pt,
              inner sep=1.2pt,font=\scSmall,text=scInk,
              minimum width=3.0mm,minimum height=3.0mm},
  place/.style={circle,fill=scSpF,draw=scSpE,line width=0.45pt,
                minimum size=4.0mm,inner sep=0pt},
  trans/.style={rectangle,fill=scRxE,draw=scRxE,minimum width=1.1mm,
                minimum height=3.6mm,inner sep=0pt},
  token/.style={circle,fill=scInk,minimum size=1.1mm,inner sep=0pt},
  % edges -----------------------------------------------------------------
  reac/.style={->,line width=0.5pt,draw=scReac,shorten >=0.8pt,shorten <=0.8pt},
  prod/.style={->,line width=0.5pt,draw=scProd,shorten >=0.8pt,shorten <=0.8pt},
  cat/.style={->,line width=0.5pt,draw=scHiE,shorten >=0.8pt,shorten <=0.8pt},
  flow/.style={-{Stealth[round,length=2.2mm,width=1.8mm]},line width=1.1pt,
               draw=scInkM!45},
  flab/.style={font=\scSmall,text=scInkM,inner sep=2pt},
  stoi/.style={font=\scSmall,text=scInkM,inner sep=0.8pt},
  % verdict marks: shape carries the meaning, colour only reinforces it
  vpass/.style={circle,fill=scAccentF,draw=scAccent,text=scAccent,
                font=\fontsize{6}{6}\selectfont\bfseries,
                minimum size=3.2mm,inner sep=0pt},
  vfail/.style={circle,fill=scAlertF,draw=scAlert,text=scAlert,
                font=\fontsize{6}{6}\selectfont\bfseries,
                minimum size=3.2mm,inner sep=0pt},
  callout/.style={font=\scSmall,text=scInkM,align=left},
}

% \scHead[<badge colour>]{<anchor coord>}{<letter>}{<title>}{<subtitle or empty>}
% One header treatment for every panel in the paper.
\providecommand{\scHead}[5][scInk]{%
  \node[hbullet,fill=#1,anchor=north west] (scbadge#3) at ($#2+(0,-0.015)$)
    {\MakeUppercase{#3}};
  \node[ctag,anchor=west] at ($(scbadge#3.east)+(0.14,0)$) {#4};
  \ifx\relax#5\relax\else
    \node[anchor=north west,font=\scSmall,text=scInkM]
      at ($#2+(0.56,-0.245)$) {#5};
  \fi}

% --- shared pgfplots look (only if the figure loaded pgfplots) ------------
\ifdefined\pgfplotsset
  \pgfplotsset{compat=1.18}
  \pgfplotsset{
    synaxis/.style={
      scale only axis,
      tick align=outside, tick pos=left,
      axis line style={draw=scInkM!55,line width=0.4pt},
      % One grid, on the value axis only.  No minor grid: at these sizes it
      % turns the plot frame into a grey wash.
      ymajorgrids=true, xmajorgrids=false,
      major grid style={draw=scRuleL,line width=0.3pt},
      major tick length=1.6pt,
      label style={font=\scBody,color=scInk},
      tick label style={font=\scSmall,color=scInkM},
      legend style={font=\scSmall,draw=none,fill=none,inner sep=1pt,
                    row sep=0.4pt,legend columns=1},
      legend cell align=left,
      every axis plot/.append style={line width=0.7pt,mark size=1.2pt},
    },
  }
\fi
}{% ---------------------------------------------------------------------------
% synstyle.tex --- shared look for every TikZ figure of the SynKit CRN paper.
%
% Input it from a figure file:
%     \IfFileExists{synstyle.tex}{\input{synstyle}}{\input{Figure/synstyle}}
% so that the figure compiles both from Figure/ and from the paper root.
%
% The visual vocabulary follows ../../../../Style/style.tex: the
% colour-blind-safe Okabe--Ito accents, quiet neutral rules and wash, circular
% panel badges, and a compact technical type scale are shared by the companion
% papers.  Semantic aliases below keep the CRN meaning explicit.
%
% DESIGN CONTRACT -----------------------------------------------------------
% 1. Geometry.  Every figure is laid out on a canvas exactly \scW wide (the
%    sn-jnl text width, 372 pt) and is included unscaled, so one point of type
%    is the same physical size in all five figures.  Open each picture with
%    \scCanvas{<height in cm>} to pin the bounding box.
% 2. Colour is semantic, never decorative.  Four roles only:
%       scAccent (green)       structure, the normal/expected path
%       scMark   (orange)      catalyst, certificate, or highlight
%       scAlert  (vermillion)  failure, violation, or uncovered set
%       scRef    (blue)        secondary or reference series
%    A panel gets at most one accent.  Titles are always ink; panel identity
%    is also encoded by a letter inside a circular badge, not by colour alone.
% 3. Type scale.  Three sizes, no others, and nothing below 6 pt:
%       \scTitle (7 pt bold)  panel titles
%       \scBody  (6.5 pt)     labels, values, in-figure prose
%       \scSmall (6 pt)       units, footnotes, dense graph labels
% ---------------------------------------------------------------------------
% Load guard.  The paper inputs every figure source inline, so this file is
% reached once per figure; without the guard the second \newlength below
% would abort the run.
\ifdefined\SynKitStyleLoaded\endinput\fi
\newcommand{\SynKitStyleLoaded}{}

% Load in the document preamble. XeLaTeX/LuaLaTeX and an installed Arial
% family are required; fontspec reports an error instead of substituting.
\RequirePackage[no-math]{fontspec}
\newfontfamily\scArial{Arial}[NFSSFamily=scArial,Ligatures=TeX]
% Keep manuscript mathematics unchanged; select Arial letters and numerals
% only inside figures. TeX supplies specialist mathematical symbols.
\RequirePackage[subdued,italic]{mathastext}
\MTDeclareVersion{synarial}{TU}{scArial}{m}{n}
\SetMathAlphabet{\Mathsf}{synarial}{TU}{scArial}{m}{n}
\SetMathAlphabet{\Mathtt}{synarial}{TU}{scArial}{m}{n}

\usetikzlibrary{positioning,calc,fit,backgrounds,arrows.meta,shapes.geometric,
                shapes.misc,decorations.pathreplacing,patterns,math}

% --- canvas ---------------------------------------------------------------
% 372 pt text width minus the 1 pt standalone border on each side.
\newlength{\scW}\setlength{\scW}{13.0326cm}
\providecommand{\scWnum}{13.0326}
% \scCanvas{<height>} pins the picture to the full text width, y running down
% from 0, so every figure shares one coordinate convention.  It *sets* the
% bounding box rather than contributing to it, so the emitted page is exactly
% the text width and every figure is included at 1:1 -- otherwise a stray
% invisible extent (pgfplots reserves a little around an axis) rescales one
% figure relative to the others and the same \scSmall prints at two sizes.
% Anything drawn outside the canvas is therefore cropped: keep content inside,
% and check with
%     \providecommand{\scCanvas}[1]{\path[draw=red] (0,0) rectangle (\scWnum,-#1);}
\providecommand{\scCanvas}[1]{%
  \path[use as bounding box] (0,0) rectangle (\scWnum,-#1);}

% --- palette --------------------------------------------------------------
% Canonical companion-paper values (../../../../Style/style.tex).
\definecolor{npgOrange}{HTML}{E69F00}
\definecolor{npgSky}{HTML}{56B4E9}
\definecolor{npgGreen}{HTML}{009E73}
\definecolor{npgYellow}{HTML}{F0E442}
\definecolor{npgBlue}{HTML}{0072B2}
\definecolor{npgVermillion}{HTML}{E64B35}
\definecolor{npgPurple}{HTML}{CC79A7}
\definecolor{figInk}{HTML}{1A1A1A}
\definecolor{figMuted}{HTML}{666666}
\definecolor{figHair}{HTML}{D9D9D9}
\definecolor{figWash}{HTML}{F7F7F7}

\colorlet{scInk}{figInk}              % primary text
\colorlet{scInkM}{figMuted}           % secondary text
\colorlet{scPaper}{figWash}           % quiet cards and note strips
\colorlet{scSpF}{npgSky!12}           % species fill
\colorlet{scSpE}{npgBlue}             % species edge
\colorlet{scRxF}{npgOrange!12}        % reaction fill
\colorlet{scRxE}{npgOrange!88!black}  % reaction edge
\colorlet{scReac}{npgBlue!78!black}   % reactant edge
\colorlet{scProd}{npgGreen!82!black}  % product edge
\colorlet{scHiF}{npgOrange!13}        % highlight fill
\colorlet{scHiE}{npgOrange!82!black}  % highlight edge
% Stage / accent tints.  stDE remains available for the fifth plot series.
\colorlet{stAF}{npgGreen!9}\colorlet{stAE}{npgGreen}
\colorlet{stBF}{npgOrange!11}\colorlet{stBE}{npgOrange}
\colorlet{stCF}{npgBlue!9}\colorlet{stCE}{npgBlue}
\colorlet{stDF}{npgPurple!10}\colorlet{stDE}{npgPurple!85!black}
\colorlet{stEF}{npgVermillion!9}\colorlet{stEE}{npgVermillion}

% Semantic aliases.  Figures should speak in these, not in stAE/stBE/...
\colorlet{scAccent}{stAE}\colorlet{scAccentF}{stAF}   % structure / normal
\colorlet{scMark}{stBE}\colorlet{scMarkF}{stBF}       % highlight / certificate
\colorlet{scRef}{stCE}\colorlet{scRefF}{stCF}         % secondary series
\colorlet{scAlert}{stEE}\colorlet{scAlertF}{stEF}     % failure / violation
\colorlet{scRule}{figHair}                            % card and axis hairlines
\colorlet{scRuleL}{figHair!55}                        % grid, dividers

% --- text styles ----------------------------------------------------------
% Three sizes, fixed in points so they do not drift with the document class.
% Arial is shared by labels, headings, panel badges, and chemistry labels.
\providecommand{\scTitle}{\scArial\fontsize{7}{8.4}\selectfont\bfseries\color{scInk}}
\providecommand{\scBody}{\scArial\fontsize{6.5}{7.8}\selectfont\color{scInkM}}
\providecommand{\scSmall}{\scArial\fontsize{6}{7.2}\selectfont\color{scInkM}}
\providecommand{\scSub}[1]{{\scSmall\itshape #1}}
\providecommand{\scPanel}[1]{{\scTitle(#1)}}
% Numerals that carry a result: ink, not colour, so the eye reads the value.
\providecommand{\scNum}{\scArial\fontsize{10}{11}\selectfont\bfseries\color{scInk}}

% --- shared TikZ styles ---------------------------------------------------
\tikzset{
  % text never hyphenates inside a figure
  every node/.append style={execute at begin node={%
      \hyphenpenalty=10000 \exhyphenpenalty=10000 \hbadness=10000 }},
  syn/.style={font=\scArial,color=scInk,
              execute at begin picture={\MTversion*{synarial}},
              >={Stealth[round,length=2.4mm,width=1.9mm]}},
  % containers ------------------------------------------------------------
  % One card chrome for the whole paper: white, 3 pt radius, one hairline.
  panel/.style={rounded corners=2.5pt,fill=none,draw=scRule,line width=0.45pt},
  card/.style={rounded corners=2.2pt,fill=scPaper,draw=scRule,line width=0.45pt,
               inner sep=5pt,align=left,anchor=north west},
  % a quiet strip for a panel's closing fact
  note/.style={rounded corners=2.5pt,fill=scPaper,draw=scRuleL,line width=0.4pt,
               inner xsep=6pt,inner ysep=3.5pt,align=left,font=\scSmall},
  % small labelled value pill
  chip/.style={rounded corners=2pt,fill=scPaper,draw=scRule,line width=0.4pt,
               inner xsep=4pt,inner ysep=2pt,font=\scSmall},
  % Circular letter badge: the same accessible panel marker used in Style/*.
  hbullet/.style={circle,fill=scInk,text=white,draw=none,
                  font=\scArial\fontsize{6}{6}\selectfont\bfseries,
                  minimum size=4.2mm,inner sep=0pt},
  ctag/.style={font=\scTitle,anchor=north west,text=scInk},
  % CRN primitives --------------------------------------------------------
  species/.style={rounded rectangle,fill=scSpF,draw=scSpE,line width=0.45pt,
                  inner xsep=2.4pt,inner ysep=1.6pt,font=\scSmall,
                  text=scInk,minimum height=3.6mm},
  cofactor/.style={species,fill=scRefF,draw=scRef},
  catalyst/.style={species,fill=scHiF,draw=scHiE},
  rxn/.style={rectangle,fill=scRxF,draw=scRxE,line width=0.45pt,
              inner sep=1.2pt,font=\scSmall,text=scInk,
              minimum width=3.0mm,minimum height=3.0mm},
  place/.style={circle,fill=scSpF,draw=scSpE,line width=0.45pt,
                minimum size=4.0mm,inner sep=0pt},
  trans/.style={rectangle,fill=scRxE,draw=scRxE,minimum width=1.1mm,
                minimum height=3.6mm,inner sep=0pt},
  token/.style={circle,fill=scInk,minimum size=1.1mm,inner sep=0pt},
  % edges -----------------------------------------------------------------
  reac/.style={->,line width=0.5pt,draw=scReac,shorten >=0.8pt,shorten <=0.8pt},
  prod/.style={->,line width=0.5pt,draw=scProd,shorten >=0.8pt,shorten <=0.8pt},
  cat/.style={->,line width=0.5pt,draw=scHiE,shorten >=0.8pt,shorten <=0.8pt},
  flow/.style={-{Stealth[round,length=2.2mm,width=1.8mm]},line width=1.1pt,
               draw=scInkM!45},
  flab/.style={font=\scSmall,text=scInkM,inner sep=2pt},
  stoi/.style={font=\scSmall,text=scInkM,inner sep=0.8pt},
  % verdict marks: shape carries the meaning, colour only reinforces it
  vpass/.style={circle,fill=scAccentF,draw=scAccent,text=scAccent,
                font=\fontsize{6}{6}\selectfont\bfseries,
                minimum size=3.2mm,inner sep=0pt},
  vfail/.style={circle,fill=scAlertF,draw=scAlert,text=scAlert,
                font=\fontsize{6}{6}\selectfont\bfseries,
                minimum size=3.2mm,inner sep=0pt},
  callout/.style={font=\scSmall,text=scInkM,align=left},
}

% \scHead[<badge colour>]{<anchor coord>}{<letter>}{<title>}{<subtitle or empty>}
% One header treatment for every panel in the paper.
\providecommand{\scHead}[5][scInk]{%
  \node[hbullet,fill=#1,anchor=north west] (scbadge#3) at ($#2+(0,-0.015)$)
    {\MakeUppercase{#3}};
  \node[ctag,anchor=west] at ($(scbadge#3.east)+(0.14,0)$) {#4};
  \ifx\relax#5\relax\else
    \node[anchor=north west,font=\scSmall,text=scInkM]
      at ($#2+(0.56,-0.245)$) {#5};
  \fi}

% --- shared pgfplots look (only if the figure loaded pgfplots) ------------
\ifdefined\pgfplotsset
  \pgfplotsset{compat=1.18}
  \pgfplotsset{
    synaxis/.style={
      scale only axis,
      tick align=outside, tick pos=left,
      axis line style={draw=scInkM!55,line width=0.4pt},
      % One grid, on the value axis only.  No minor grid: at these sizes it
      % turns the plot frame into a grey wash.
      ymajorgrids=true, xmajorgrids=false,
      major grid style={draw=scRuleL,line width=0.3pt},
      major tick length=1.6pt,
      label style={font=\scBody,color=scInk},
      tick label style={font=\scSmall,color=scInkM},
      legend style={font=\scSmall,draw=none,fill=none,inner sep=1pt,
                    row sep=0.4pt,legend columns=1},
      legend cell align=left,
      every axis plot/.append style={line width=0.7pt,mark size=1.2pt},
    },
  }
\fi
}
% so that the figure compiles both from Figure/ and from the paper root.
%
% The visual vocabulary follows ../../../../Style/style.tex: the
% colour-blind-safe Okabe--Ito accents, quiet neutral rules and wash, circular
% panel badges, and a compact technical type scale are shared by the companion
% papers.  Semantic aliases below keep the CRN meaning explicit.
%
% DESIGN CONTRACT -----------------------------------------------------------
% 1. Geometry.  Every figure is laid out on a canvas exactly \scW wide (the
%    sn-jnl text width, 372 pt) and is included unscaled, so one point of type
%    is the same physical size in all five figures.  Open each picture with
%    \scCanvas{<height in cm>} to pin the bounding box.
% 2. Colour is semantic, never decorative.  Four roles only:
%       scAccent (green)       structure, the normal/expected path
%       scMark   (orange)      catalyst, certificate, or highlight
%       scAlert  (vermillion)  failure, violation, or uncovered set
%       scRef    (blue)        secondary or reference series
%    A panel gets at most one accent.  Titles are always ink; panel identity
%    is also encoded by a letter inside a circular badge, not by colour alone.
% 3. Type scale.  Three sizes, no others, and nothing below 6 pt:
%       \scTitle (7 pt bold)  panel titles
%       \scBody  (6.5 pt)     labels, values, in-figure prose
%       \scSmall (6 pt)       units, footnotes, dense graph labels
% ---------------------------------------------------------------------------
% Load guard.  The paper inputs every figure source inline, so this file is
% reached once per figure; without the guard the second \newlength below
% would abort the run.
\ifdefined\SynKitStyleLoaded\endinput\fi
\newcommand{\SynKitStyleLoaded}{}

% Load in the document preamble. XeLaTeX/LuaLaTeX and an installed Arial
% family are required; fontspec reports an error instead of substituting.
\RequirePackage[no-math]{fontspec}
\newfontfamily\scArial{Arial}[NFSSFamily=scArial,Ligatures=TeX]
% Keep manuscript mathematics unchanged; select Arial letters and numerals
% only inside figures. TeX supplies specialist mathematical symbols.
\RequirePackage[subdued,italic]{mathastext}
\MTDeclareVersion{synarial}{TU}{scArial}{m}{n}
\SetMathAlphabet{\Mathsf}{synarial}{TU}{scArial}{m}{n}
\SetMathAlphabet{\Mathtt}{synarial}{TU}{scArial}{m}{n}

\usetikzlibrary{positioning,calc,fit,backgrounds,arrows.meta,shapes.geometric,
                shapes.misc,decorations.pathreplacing,patterns,math}

% --- canvas ---------------------------------------------------------------
% 372 pt text width minus the 1 pt standalone border on each side.
\newlength{\scW}\setlength{\scW}{13.0326cm}
\providecommand{\scWnum}{13.0326}
% \scCanvas{<height>} pins the picture to the full text width, y running down
% from 0, so every figure shares one coordinate convention.  It *sets* the
% bounding box rather than contributing to it, so the emitted page is exactly
% the text width and every figure is included at 1:1 -- otherwise a stray
% invisible extent (pgfplots reserves a little around an axis) rescales one
% figure relative to the others and the same \scSmall prints at two sizes.
% Anything drawn outside the canvas is therefore cropped: keep content inside,
% and check with
%     \providecommand{\scCanvas}[1]{\path[draw=red] (0,0) rectangle (\scWnum,-#1);}
\providecommand{\scCanvas}[1]{%
  \path[use as bounding box] (0,0) rectangle (\scWnum,-#1);}

% --- palette --------------------------------------------------------------
% Canonical companion-paper values (../../../../Style/style.tex).
\definecolor{npgOrange}{HTML}{E69F00}
\definecolor{npgSky}{HTML}{56B4E9}
\definecolor{npgGreen}{HTML}{009E73}
\definecolor{npgYellow}{HTML}{F0E442}
\definecolor{npgBlue}{HTML}{0072B2}
\definecolor{npgVermillion}{HTML}{E64B35}
\definecolor{npgPurple}{HTML}{CC79A7}
\definecolor{figInk}{HTML}{1A1A1A}
\definecolor{figMuted}{HTML}{666666}
\definecolor{figHair}{HTML}{D9D9D9}
\definecolor{figWash}{HTML}{F7F7F7}

\colorlet{scInk}{figInk}              % primary text
\colorlet{scInkM}{figMuted}           % secondary text
\colorlet{scPaper}{figWash}           % quiet cards and note strips
\colorlet{scSpF}{npgSky!12}           % species fill
\colorlet{scSpE}{npgBlue}             % species edge
\colorlet{scRxF}{npgOrange!12}        % reaction fill
\colorlet{scRxE}{npgOrange!88!black}  % reaction edge
\colorlet{scReac}{npgBlue!78!black}   % reactant edge
\colorlet{scProd}{npgGreen!82!black}  % product edge
\colorlet{scHiF}{npgOrange!13}        % highlight fill
\colorlet{scHiE}{npgOrange!82!black}  % highlight edge
% Stage / accent tints.  stDE remains available for the fifth plot series.
\colorlet{stAF}{npgGreen!9}\colorlet{stAE}{npgGreen}
\colorlet{stBF}{npgOrange!11}\colorlet{stBE}{npgOrange}
\colorlet{stCF}{npgBlue!9}\colorlet{stCE}{npgBlue}
\colorlet{stDF}{npgPurple!10}\colorlet{stDE}{npgPurple!85!black}
\colorlet{stEF}{npgVermillion!9}\colorlet{stEE}{npgVermillion}

% Semantic aliases.  Figures should speak in these, not in stAE/stBE/...
\colorlet{scAccent}{stAE}\colorlet{scAccentF}{stAF}   % structure / normal
\colorlet{scMark}{stBE}\colorlet{scMarkF}{stBF}       % highlight / certificate
\colorlet{scRef}{stCE}\colorlet{scRefF}{stCF}         % secondary series
\colorlet{scAlert}{stEE}\colorlet{scAlertF}{stEF}     % failure / violation
\colorlet{scRule}{figHair}                            % card and axis hairlines
\colorlet{scRuleL}{figHair!55}                        % grid, dividers

% --- text styles ----------------------------------------------------------
% Three sizes, fixed in points so they do not drift with the document class.
% Arial is shared by labels, headings, panel badges, and chemistry labels.
\providecommand{\scTitle}{\scArial\fontsize{7}{8.4}\selectfont\bfseries\color{scInk}}
\providecommand{\scBody}{\scArial\fontsize{6.5}{7.8}\selectfont\color{scInkM}}
\providecommand{\scSmall}{\scArial\fontsize{6}{7.2}\selectfont\color{scInkM}}
\providecommand{\scSub}[1]{{\scSmall\itshape #1}}
\providecommand{\scPanel}[1]{{\scTitle(#1)}}
% Numerals that carry a result: ink, not colour, so the eye reads the value.
\providecommand{\scNum}{\scArial\fontsize{10}{11}\selectfont\bfseries\color{scInk}}

% --- shared TikZ styles ---------------------------------------------------
\tikzset{
  % text never hyphenates inside a figure
  every node/.append style={execute at begin node={%
      \hyphenpenalty=10000 \exhyphenpenalty=10000 \hbadness=10000 }},
  syn/.style={font=\scArial,color=scInk,
              execute at begin picture={\MTversion*{synarial}},
              >={Stealth[round,length=2.4mm,width=1.9mm]}},
  % containers ------------------------------------------------------------
  % One card chrome for the whole paper: white, 3 pt radius, one hairline.
  panel/.style={rounded corners=2.5pt,fill=none,draw=scRule,line width=0.45pt},
  card/.style={rounded corners=2.2pt,fill=scPaper,draw=scRule,line width=0.45pt,
               inner sep=5pt,align=left,anchor=north west},
  % a quiet strip for a panel's closing fact
  note/.style={rounded corners=2.5pt,fill=scPaper,draw=scRuleL,line width=0.4pt,
               inner xsep=6pt,inner ysep=3.5pt,align=left,font=\scSmall},
  % small labelled value pill
  chip/.style={rounded corners=2pt,fill=scPaper,draw=scRule,line width=0.4pt,
               inner xsep=4pt,inner ysep=2pt,font=\scSmall},
  % Circular letter badge: the same accessible panel marker used in Style/*.
  hbullet/.style={circle,fill=scInk,text=white,draw=none,
                  font=\scArial\fontsize{6}{6}\selectfont\bfseries,
                  minimum size=4.2mm,inner sep=0pt},
  ctag/.style={font=\scTitle,anchor=north west,text=scInk},
  % CRN primitives --------------------------------------------------------
  species/.style={rounded rectangle,fill=scSpF,draw=scSpE,line width=0.45pt,
                  inner xsep=2.4pt,inner ysep=1.6pt,font=\scSmall,
                  text=scInk,minimum height=3.6mm},
  cofactor/.style={species,fill=scRefF,draw=scRef},
  catalyst/.style={species,fill=scHiF,draw=scHiE},
  rxn/.style={rectangle,fill=scRxF,draw=scRxE,line width=0.45pt,
              inner sep=1.2pt,font=\scSmall,text=scInk,
              minimum width=3.0mm,minimum height=3.0mm},
  place/.style={circle,fill=scSpF,draw=scSpE,line width=0.45pt,
                minimum size=4.0mm,inner sep=0pt},
  trans/.style={rectangle,fill=scRxE,draw=scRxE,minimum width=1.1mm,
                minimum height=3.6mm,inner sep=0pt},
  token/.style={circle,fill=scInk,minimum size=1.1mm,inner sep=0pt},
  % edges -----------------------------------------------------------------
  reac/.style={->,line width=0.5pt,draw=scReac,shorten >=0.8pt,shorten <=0.8pt},
  prod/.style={->,line width=0.5pt,draw=scProd,shorten >=0.8pt,shorten <=0.8pt},
  cat/.style={->,line width=0.5pt,draw=scHiE,shorten >=0.8pt,shorten <=0.8pt},
  flow/.style={-{Stealth[round,length=2.2mm,width=1.8mm]},line width=1.1pt,
               draw=scInkM!45},
  flab/.style={font=\scSmall,text=scInkM,inner sep=2pt},
  stoi/.style={font=\scSmall,text=scInkM,inner sep=0.8pt},
  % verdict marks: shape carries the meaning, colour only reinforces it
  vpass/.style={circle,fill=scAccentF,draw=scAccent,text=scAccent,
                font=\fontsize{6}{6}\selectfont\bfseries,
                minimum size=3.2mm,inner sep=0pt},
  vfail/.style={circle,fill=scAlertF,draw=scAlert,text=scAlert,
                font=\fontsize{6}{6}\selectfont\bfseries,
                minimum size=3.2mm,inner sep=0pt},
  callout/.style={font=\scSmall,text=scInkM,align=left},
}

% \scHead[<badge colour>]{<anchor coord>}{<letter>}{<title>}{<subtitle or empty>}
% One header treatment for every panel in the paper.
\providecommand{\scHead}[5][scInk]{%
  \node[hbullet,fill=#1,anchor=north west] (scbadge#3) at ($#2+(0,-0.015)$)
    {\MakeUppercase{#3}};
  \node[ctag,anchor=west] at ($(scbadge#3.east)+(0.14,0)$) {#4};
  \ifx\relax#5\relax\else
    \node[anchor=north west,font=\scSmall,text=scInkM]
      at ($#2+(0.56,-0.245)$) {#5};
  \fi}

% --- shared pgfplots look (only if the figure loaded pgfplots) ------------
\ifdefined\pgfplotsset
  \pgfplotsset{compat=1.18}
  \pgfplotsset{
    synaxis/.style={
      scale only axis,
      tick align=outside, tick pos=left,
      axis line style={draw=scInkM!55,line width=0.4pt},
      % One grid, on the value axis only.  No minor grid: at these sizes it
      % turns the plot frame into a grey wash.
      ymajorgrids=true, xmajorgrids=false,
      major grid style={draw=scRuleL,line width=0.3pt},
      major tick length=1.6pt,
      label style={font=\scBody,color=scInk},
      tick label style={font=\scSmall,color=scInkM},
      legend style={font=\scSmall,draw=none,fill=none,inner sep=1pt,
                    row sep=0.4pt,legend columns=1},
      legend cell align=left,
      every axis plot/.append style={line width=0.7pt,mark size=1.2pt},
    },
  }
\fi
}
% so that the figure compiles both from Figure/ and from the paper root.
%
% The visual vocabulary follows ../../../../Style/style.tex: the
% colour-blind-safe Okabe--Ito accents, quiet neutral rules and wash, circular
% panel badges, and a compact technical type scale are shared by the companion
% papers.  Semantic aliases below keep the CRN meaning explicit.
%
% DESIGN CONTRACT -----------------------------------------------------------
% 1. Geometry.  Every figure is laid out on a canvas exactly \scW wide (the
%    sn-jnl text width, 372 pt) and is included unscaled, so one point of type
%    is the same physical size in all five figures.  Open each picture with
%    \scCanvas{<height in cm>} to pin the bounding box.
% 2. Colour is semantic, never decorative.  Four roles only:
%       scAccent (green)       structure, the normal/expected path
%       scMark   (orange)      catalyst, certificate, or highlight
%       scAlert  (vermillion)  failure, violation, or uncovered set
%       scRef    (blue)        secondary or reference series
%    A panel gets at most one accent.  Titles are always ink; panel identity
%    is also encoded by a letter inside a circular badge, not by colour alone.
% 3. Type scale.  Three sizes, no others, and nothing below 6 pt:
%       \scTitle (7 pt bold)  panel titles
%       \scBody  (6.5 pt)     labels, values, in-figure prose
%       \scSmall (6 pt)       units, footnotes, dense graph labels
% ---------------------------------------------------------------------------
% Load guard.  The paper inputs every figure source inline, so this file is
% reached once per figure; without the guard the second \newlength below
% would abort the run.
\ifdefined\SynKitStyleLoaded\endinput\fi
\newcommand{\SynKitStyleLoaded}{}

% Load in the document preamble. XeLaTeX/LuaLaTeX and an installed Arial
% family are required; fontspec reports an error instead of substituting.
\RequirePackage[no-math]{fontspec}
\newfontfamily\scArial{Arial}[NFSSFamily=scArial,Ligatures=TeX]
% Keep manuscript mathematics unchanged; select Arial letters and numerals
% only inside figures. TeX supplies specialist mathematical symbols.
\RequirePackage[subdued,italic]{mathastext}
\MTDeclareVersion{synarial}{TU}{scArial}{m}{n}
\SetMathAlphabet{\Mathsf}{synarial}{TU}{scArial}{m}{n}
\SetMathAlphabet{\Mathtt}{synarial}{TU}{scArial}{m}{n}

\usetikzlibrary{positioning,calc,fit,backgrounds,arrows.meta,shapes.geometric,
                shapes.misc,decorations.pathreplacing,patterns,math}

% --- canvas ---------------------------------------------------------------
% 372 pt text width minus the 1 pt standalone border on each side.
\newlength{\scW}\setlength{\scW}{13.0326cm}
\providecommand{\scWnum}{13.0326}
% \scCanvas{<height>} pins the picture to the full text width, y running down
% from 0, so every figure shares one coordinate convention.  It *sets* the
% bounding box rather than contributing to it, so the emitted page is exactly
% the text width and every figure is included at 1:1 -- otherwise a stray
% invisible extent (pgfplots reserves a little around an axis) rescales one
% figure relative to the others and the same \scSmall prints at two sizes.
% Anything drawn outside the canvas is therefore cropped: keep content inside,
% and check with
%     \providecommand{\scCanvas}[1]{\path[draw=red] (0,0) rectangle (\scWnum,-#1);}
\providecommand{\scCanvas}[1]{%
  \path[use as bounding box] (0,0) rectangle (\scWnum,-#1);}

% --- palette --------------------------------------------------------------
% Canonical companion-paper values (../../../../Style/style.tex).
\definecolor{npgOrange}{HTML}{E69F00}
\definecolor{npgSky}{HTML}{56B4E9}
\definecolor{npgGreen}{HTML}{009E73}
\definecolor{npgYellow}{HTML}{F0E442}
\definecolor{npgBlue}{HTML}{0072B2}
\definecolor{npgVermillion}{HTML}{E64B35}
\definecolor{npgPurple}{HTML}{CC79A7}
\definecolor{figInk}{HTML}{1A1A1A}
\definecolor{figMuted}{HTML}{666666}
\definecolor{figHair}{HTML}{D9D9D9}
\definecolor{figWash}{HTML}{F7F7F7}

\colorlet{scInk}{figInk}              % primary text
\colorlet{scInkM}{figMuted}           % secondary text
\colorlet{scPaper}{figWash}           % quiet cards and note strips
\colorlet{scSpF}{npgSky!12}           % species fill
\colorlet{scSpE}{npgBlue}             % species edge
\colorlet{scRxF}{npgOrange!12}        % reaction fill
\colorlet{scRxE}{npgOrange!88!black}  % reaction edge
\colorlet{scReac}{npgBlue!78!black}   % reactant edge
\colorlet{scProd}{npgGreen!82!black}  % product edge
\colorlet{scHiF}{npgOrange!13}        % highlight fill
\colorlet{scHiE}{npgOrange!82!black}  % highlight edge
% Stage / accent tints.  stDE remains available for the fifth plot series.
\colorlet{stAF}{npgGreen!9}\colorlet{stAE}{npgGreen}
\colorlet{stBF}{npgOrange!11}\colorlet{stBE}{npgOrange}
\colorlet{stCF}{npgBlue!9}\colorlet{stCE}{npgBlue}
\colorlet{stDF}{npgPurple!10}\colorlet{stDE}{npgPurple!85!black}
\colorlet{stEF}{npgVermillion!9}\colorlet{stEE}{npgVermillion}

% Semantic aliases.  Figures should speak in these, not in stAE/stBE/...
\colorlet{scAccent}{stAE}\colorlet{scAccentF}{stAF}   % structure / normal
\colorlet{scMark}{stBE}\colorlet{scMarkF}{stBF}       % highlight / certificate
\colorlet{scRef}{stCE}\colorlet{scRefF}{stCF}         % secondary series
\colorlet{scAlert}{stEE}\colorlet{scAlertF}{stEF}     % failure / violation
\colorlet{scRule}{figHair}                            % card and axis hairlines
\colorlet{scRuleL}{figHair!55}                        % grid, dividers

% --- text styles ----------------------------------------------------------
% Three sizes, fixed in points so they do not drift with the document class.
% Arial is shared by labels, headings, panel badges, and chemistry labels.
\providecommand{\scTitle}{\scArial\fontsize{7}{8.4}\selectfont\bfseries\color{scInk}}
\providecommand{\scBody}{\scArial\fontsize{6.5}{7.8}\selectfont\color{scInkM}}
\providecommand{\scSmall}{\scArial\fontsize{6}{7.2}\selectfont\color{scInkM}}
\providecommand{\scSub}[1]{{\scSmall\itshape #1}}
\providecommand{\scPanel}[1]{{\scTitle(#1)}}
% Numerals that carry a result: ink, not colour, so the eye reads the value.
\providecommand{\scNum}{\scArial\fontsize{10}{11}\selectfont\bfseries\color{scInk}}

% --- shared TikZ styles ---------------------------------------------------
\tikzset{
  % text never hyphenates inside a figure
  every node/.append style={execute at begin node={%
      \hyphenpenalty=10000 \exhyphenpenalty=10000 \hbadness=10000 }},
  syn/.style={font=\scArial,color=scInk,
              execute at begin picture={\MTversion*{synarial}},
              >={Stealth[round,length=2.4mm,width=1.9mm]}},
  % containers ------------------------------------------------------------
  % One card chrome for the whole paper: white, 3 pt radius, one hairline.
  panel/.style={rounded corners=2.5pt,fill=none,draw=scRule,line width=0.45pt},
  card/.style={rounded corners=2.2pt,fill=scPaper,draw=scRule,line width=0.45pt,
               inner sep=5pt,align=left,anchor=north west},
  % a quiet strip for a panel's closing fact
  note/.style={rounded corners=2.5pt,fill=scPaper,draw=scRuleL,line width=0.4pt,
               inner xsep=6pt,inner ysep=3.5pt,align=left,font=\scSmall},
  % small labelled value pill
  chip/.style={rounded corners=2pt,fill=scPaper,draw=scRule,line width=0.4pt,
               inner xsep=4pt,inner ysep=2pt,font=\scSmall},
  % Circular letter badge: the same accessible panel marker used in Style/*.
  hbullet/.style={circle,fill=scInk,text=white,draw=none,
                  font=\scArial\fontsize{6}{6}\selectfont\bfseries,
                  minimum size=4.2mm,inner sep=0pt},
  ctag/.style={font=\scTitle,anchor=north west,text=scInk},
  % CRN primitives --------------------------------------------------------
  species/.style={rounded rectangle,fill=scSpF,draw=scSpE,line width=0.45pt,
                  inner xsep=2.4pt,inner ysep=1.6pt,font=\scSmall,
                  text=scInk,minimum height=3.6mm},
  cofactor/.style={species,fill=scRefF,draw=scRef},
  catalyst/.style={species,fill=scHiF,draw=scHiE},
  rxn/.style={rectangle,fill=scRxF,draw=scRxE,line width=0.45pt,
              inner sep=1.2pt,font=\scSmall,text=scInk,
              minimum width=3.0mm,minimum height=3.0mm},
  place/.style={circle,fill=scSpF,draw=scSpE,line width=0.45pt,
                minimum size=4.0mm,inner sep=0pt},
  trans/.style={rectangle,fill=scRxE,draw=scRxE,minimum width=1.1mm,
                minimum height=3.6mm,inner sep=0pt},
  token/.style={circle,fill=scInk,minimum size=1.1mm,inner sep=0pt},
  % edges -----------------------------------------------------------------
  reac/.style={->,line width=0.5pt,draw=scReac,shorten >=0.8pt,shorten <=0.8pt},
  prod/.style={->,line width=0.5pt,draw=scProd,shorten >=0.8pt,shorten <=0.8pt},
  cat/.style={->,line width=0.5pt,draw=scHiE,shorten >=0.8pt,shorten <=0.8pt},
  flow/.style={-{Stealth[round,length=2.2mm,width=1.8mm]},line width=1.1pt,
               draw=scInkM!45},
  flab/.style={font=\scSmall,text=scInkM,inner sep=2pt},
  stoi/.style={font=\scSmall,text=scInkM,inner sep=0.8pt},
  % verdict marks: shape carries the meaning, colour only reinforces it
  vpass/.style={circle,fill=scAccentF,draw=scAccent,text=scAccent,
                font=\fontsize{6}{6}\selectfont\bfseries,
                minimum size=3.2mm,inner sep=0pt},
  vfail/.style={circle,fill=scAlertF,draw=scAlert,text=scAlert,
                font=\fontsize{6}{6}\selectfont\bfseries,
                minimum size=3.2mm,inner sep=0pt},
  callout/.style={font=\scSmall,text=scInkM,align=left},
}

% \scHead[<badge colour>]{<anchor coord>}{<letter>}{<title>}{<subtitle or empty>}
% One header treatment for every panel in the paper.
\providecommand{\scHead}[5][scInk]{%
  \node[hbullet,fill=#1,anchor=north west] (scbadge#3) at ($#2+(0,-0.015)$)
    {\MakeUppercase{#3}};
  \node[ctag,anchor=west] at ($(scbadge#3.east)+(0.14,0)$) {#4};
  \ifx\relax#5\relax\else
    \node[anchor=north west,font=\scSmall,text=scInkM]
      at ($#2+(0.56,-0.245)$) {#5};
  \fi}

% --- shared pgfplots look (only if the figure loaded pgfplots) ------------
\ifdefined\pgfplotsset
  \pgfplotsset{compat=1.18}
  \pgfplotsset{
    synaxis/.style={
      scale only axis,
      tick align=outside, tick pos=left,
      axis line style={draw=scInkM!55,line width=0.4pt},
      % One grid, on the value axis only.  No minor grid: at these sizes it
      % turns the plot frame into a grey wash.
      ymajorgrids=true, xmajorgrids=false,
      major grid style={draw=scRuleL,line width=0.3pt},
      major tick length=1.6pt,
      label style={font=\scBody,color=scInk},
      tick label style={font=\scSmall,color=scInkM},
      legend style={font=\scSmall,draw=none,fill=none,inner sep=1pt,
                    row sep=0.4pt,legend columns=1},
      legend cell align=left,
      every axis plot/.append style={line width=0.7pt,mark size=1.2pt},
    },
  }
\fi
}}

\tikzset{
  complex/.style={rounded corners=2pt,fill=white,draw=scRxE!70,line width=0.5pt,
    inner xsep=3.4pt,inner ysep=2.4pt,font=\scBody,text=scInk},
  cxedge/.style={-{Stealth[round,length=1.7mm,width=1.4mm]},line width=0.5pt,
    draw=scRxE,shorten >=1pt,shorten <=1pt},
  lcbox/.style={rounded corners=4pt,fill=scAccentF,draw=scAccent!45,
    line width=0.45pt},
  vrow/.style={anchor=base west,font=\scBody,text=scInkM},
}

% A signed cell of S.  Sign is carried by fill *and* by the printed sign, so
% the matrix survives greyscale printing.
\newcommand{\cell}[3]{%
  \ifnum#3=0
    \node[draw=scRuleL,fill=white,line width=0.3pt,minimum size=4.2mm,
      inner sep=0pt,font=\scSmall,text=scInkM!60] at ({#1},{#2}) {$\cdot$};
  \else
    \ifnum#3>0
      \node[draw=scAccent!55,fill=scAccentF,line width=0.3pt,minimum size=4.2mm,
        inner sep=0pt,font=\scSmall,text=scInk] at ({#1},{#2}) {$+#3$};
    \else
      \node[draw=scAlert!55,fill=scAlertF,line width=0.3pt,minimum size=4.2mm,
        inner sep=0pt,font=\scSmall,text=scInk] at ({#1},{#2}) {$#3$};
    \fi
  \fi}
% A cell of the conservation vector, and a cell of the resulting zero row.
\newcommand{\ycell}[3]{%
  \node[draw=scMark!60,fill=scMarkF,line width=0.35pt,minimum size=4.2mm,
    inner sep=0pt,font=\scSmall,text=scInk] at ({#1},{#2}) {$#3$};}
\newcommand{\zcell}[2]{%
  \node[draw=scRuleL,fill=white,line width=0.3pt,minimum size=4.2mm,
    inner sep=0pt,font=\scSmall,text=scInkM] at ({#1},{#2}) {$0$};}

\def\ax{0}\def\aw{5.00}
\def\bx{5.20}\def\bw{5.90}
\def\cx{11.30}\def\cw{5.10}
\def\ph{4.60}

\begin{tikzpicture}[syn]
\def\W{16.40}
\path[use as bounding box] (0,0) rectangle (\W,-\ph);

% =========================================================================
% (a) the network and the complex graph it induces
% =========================================================================
\begin{scope}[on background layer]
\path[panel] (\ax,0) rectangle (\ax+\aw,-\ph);
\end{scope}
\scHead{(0.22,-0.16)}{a}{complex graph}{Edelstein motif, 3 species}

\node[font=\scBody,text=scInk] at (2.50,-1.00)
  {$\mathrm A\rightleftharpoons2\mathrm A
    \qquad
    \mathrm A+\mathrm B\rightleftharpoons\mathrm C\rightleftharpoons\mathrm B$};

\node[complex] (a)  at (1.60,-1.90) {A};
\node[complex] (aa) at (3.40,-1.90) {2A};
\node[complex] (ab) at (1.15,-3.10) {A$+$B};
\node[complex] (c)  at (2.55,-3.10) {C};
\node[complex] (b)  at (3.75,-3.10) {B};
\draw[cxedge] (a)  to[out=22,in=158]   (aa);
\draw[cxedge] (aa) to[out=-158,in=-22] (a);
\draw[cxedge] (ab) to[out=22,in=158]   (c);
\draw[cxedge] (c)  to[out=-158,in=-22] (ab);
\draw[cxedge] (c)  to[out=22,in=158]   (b);
\draw[cxedge] (b)  to[out=-158,in=-22] (c);
\begin{scope}[on background layer]
  \node[lcbox,fit=(a)(aa),inner sep=5pt] (lc1) {};
  \node[lcbox,fit=(ab)(c)(b),inner sep=5pt] (lc2) {};
\end{scope}
\node[font=\scSmall,text=scAccent,anchor=west,xshift=1.5pt] at (lc1.east) {$\ell_1$};
\node[font=\scSmall,text=scAccent,anchor=west,xshift=1.5pt] at (lc2.east) {$\ell_2$};

\node[chip,anchor=north,text=scAccent] at (2.50,-3.92)
  {\textbf{weakly reversible}\ \ $n=5$,\ $\ell=2$};

\draw[flow] (\ax+\aw+0.04,-2.30) -- (\bx-0.04,-2.30);

% =========================================================================
% (b) the exact stoichiometric matrix and its left kernel
% =========================================================================
\begin{scope}[on background layer]
\path[panel] (\bx,0) rectangle (\bx+\bw,-\ph);
\end{scope}
\scHead{(\bx+0.22,-0.16)}{b}{exact stoichiometry}{integer arithmetic, no rounding}

\def\mx{6.83}\def\my{-1.55}\def\cs{0.46}
\def\yx{9.85}
\foreach \name [count=\j from 0] in {r_1,r_2,r_3,r_4,r_5,r_6}
  \node[font=\scSmall,text=scInkM] at ({\mx+\j*\cs},{\my+0.36}) {$\name$};
\node[font=\scBody,text=scInk] at ({\yx},{\my+0.36}) {$y$};
\foreach \name [count=\i from 0] in {A,B,C}
  \node[font=\scBody,text=scInk,anchor=east] at ({\mx-0.30},{\my-\i*\cs}) {$\name$};
\node[font=\scBody,text=scInk,anchor=east] at ({\mx-0.30},{\my+0.36}) {$\mathbf S$};
\cell{\mx+0*\cs}{\my-0*\cs}{1} \cell{\mx+1*\cs}{\my-0*\cs}{-1}
\cell{\mx+2*\cs}{\my-0*\cs}{-1}\cell{\mx+3*\cs}{\my-0*\cs}{1}
\cell{\mx+4*\cs}{\my-0*\cs}{0} \cell{\mx+5*\cs}{\my-0*\cs}{0}
\cell{\mx+0*\cs}{\my-1*\cs}{0} \cell{\mx+1*\cs}{\my-1*\cs}{0}
\cell{\mx+2*\cs}{\my-1*\cs}{-1}\cell{\mx+3*\cs}{\my-1*\cs}{1}
\cell{\mx+4*\cs}{\my-1*\cs}{1} \cell{\mx+5*\cs}{\my-1*\cs}{-1}
\cell{\mx+0*\cs}{\my-2*\cs}{0} \cell{\mx+1*\cs}{\my-2*\cs}{0}
\cell{\mx+2*\cs}{\my-2*\cs}{1} \cell{\mx+3*\cs}{\my-2*\cs}{-1}
\cell{\mx+4*\cs}{\my-2*\cs}{-1}\cell{\mx+5*\cs}{\my-2*\cs}{1}
% the conservation vector, as a column beside the rows it weights
\ycell{\yx}{\my-0*\cs}{0}
\ycell{\yx}{\my-1*\cs}{1}
\ycell{\yx}{\my-2*\cs}{1}
% and the row of zeros it produces
\node[font=\scSmall,text=scMark,anchor=east] at ({\mx-0.30},{\my-3*\cs-0.34})
  {$y^{\mathsf T}\mathbf S$};
\foreach \j in {0,...,5} {\zcell{\mx+\j*\cs}{\my-3*\cs-0.34}}

\node[anchor=base west,font=\scBody,text=scInk] at (\bx+0.30,\my-3*\cs-0.92)
  {$s=\operatorname{rank}\mathbf S=2$};
\node[chip,anchor=north east,text=scMark] at (\bx+\bw-0.22,\my-3*\cs-0.72)
  {\textbf{conserved pool} $\mathrm B+\mathrm C$, exact};

\draw[flow] (\bx+\bw+0.04,-2.30) -- (\cx-0.04,-2.30);

% =========================================================================
% (c) what the package reports
% =========================================================================
\begin{scope}[on background layer]
\path[panel] (\cx,0) rectangle (\cx+\cw,-\ph);
\end{scope}
\scHead[scAlert]{(\cx+0.22,-0.16)}{c}{theorem applicability}{}

\path[rounded corners=2.5pt,fill=scPaper,draw=scRule,line width=0.4pt]
  (\cx+0.22,-0.80) rectangle (\cx+\cw-0.22,-1.52);
\node[font=\scNum,anchor=west] at (\cx+0.40,-1.16) {$\delta=1$};
\draw[draw=scRule,line width=0.4pt] (\cx+1.48,-0.94) -- (\cx+1.48,-1.38);
\node[anchor=west,font=\scSmall,text=scInkM] at (\cx+1.62,-1.16)
  {$=n-\ell-s=5-2-2$};

\node[anchor=base west,font=\scBody,text=scInk] at (\cx+0.24,-1.86)
  {\bfseries Deficiency Zero Theorem};
\node[vfail] at (\cx+0.38,-2.16) {$\times$};
\node[vrow] at (\cx+0.62,-2.19) {$\delta\neq0$};

\draw[draw=scRuleL,line width=0.4pt] (\cx+0.22,-2.44) -- (\cx+\cw-0.22,-2.44);

\node[anchor=base west,font=\scBody,text=scInk] at (\cx+0.24,-2.70)
  {\bfseries Deficiency One Theorem};
\node[vpass] at (\cx+0.38,-2.98) {$\checkmark$};
\node[vrow] at (\cx+0.62,-3.01) {each $\delta_i\leq1$};
\node[vpass] at (\cx+0.38,-3.26) {$\checkmark$};
\node[vrow] at (\cx+0.62,-3.29) {one terminal strong class per $\ell_i$};
\node[vfail] at (\cx+0.38,-3.54) {$\times$};
\node[vrow] at (\cx+0.62,-3.57) {$\textstyle\sum_i\delta_i=0\neq\delta$};

\node[note,anchor=north west,text width=4.26cm,minimum height=0.46cm,
      align=center,fill=scAlertF!45,draw=scAlert!30] at (\cx+0.22,-3.86)
  {\textbf{\textcolor{scAlert}{multistationarity undecided}}};

\end{tikzpicture}
